% concatenated from heat-resurgence-arxiv1.tex
\documentclass[11pt]{amsart}

\linespread{1.18}
\setlength{\textwidth}{6.9in}
\setlength{\hoffset}{-0.90in}  
\setlength{\voffset}{-0.8in}     
\setlength{\textheight}{9.1in}
\setlength{\headsep}{10pt}
\usepackage{graphicx}
\usepackage{framed}
\usepackage{ulem}
\usepackage{pgf,tikz,pgfplots}
\usepackage{float}
\usepackage{mathrsfs}
\usepackage{mathtools}
\usepackage{mathpple}
\usepackage[all,cmtip]{xy}
\usepackage{amsmath}
\usepackage{tikz-cd}
\usepackage{mathdots}
\usepackage{yhmath}
\usepackage{cancel}
\usepackage{color}
\definecolor{deepgreen}{rgb}{0.0,0.35,0.0}
\usepackage{siunitx}
\usepackage{array}
\usepackage{multirow}
\usepackage{amssymb}
\usepackage{gensymb}
\usepackage{tabularx}
\usepackage{booktabs}
\usepackage{makecell}
\usetikzlibrary{fadings}
\usetikzlibrary{patterns}
\usetikzlibrary{shapes}
\usetikzlibrary{arrows.meta,calc,decorations.markings}

\usetikzlibrary{matrix,arrows,decorations.pathmorphing}
\usepackage{enumerate}
\usepackage{enumitem}
\usepackage{extarrows}
\usepackage{appendix}

%\usepackage{circledsteps}
\usepackage[bookmarks=true, bookmarksopen=true]{hyperref}

\DeclareMathOperator{\LHS}{LHS}
\DeclareMathOperator{\RHS}{RHS}
\DeclareMathOperator{\Id}{Id}




\theoremstyle{plain}
\newtheorem{thm}{Theorem}[section]
\newtheorem{thm-defn}[thm]{Theorem/Definition}
\newtheorem{lem}[thm]{Lemma}
\newtheorem{lem-defn}[thm]{Lemma/Definition}
\newtheorem{prop}[thm]{Proposition}
\newtheorem{cor}[thm]{Corollary}
\newtheorem{conj}[thm]{Conjecture}
\newtheorem{prob}[thm]{Problem}
\newtheorem{defn}[thm]{Definition}
\newtheorem{notn}[thm]{Notation}
\newtheorem{eg}[thm]{Example}
\newtheorem{rmk}[thm]{Remark}
\newtheorem{warn}[thm]{Warning}
\newtheorem*{claim}{Claim}

\theoremstyle{definition}
\newtheorem{assumption}{Assumption}


\hypersetup{colorlinks=true,linkcolor=blue,citecolor=blue,urlcolor=blue}


\newcommand{\R}{\mathbb{R}}
\newcommand{\C}{\mathbb{C}}
\newcommand{\Z}{\mathbb Z}
\newcommand{\dd}{\mathrm d}
\newcommand{\rank}{\operatorname{rank}}
\newcommand{\J}{\mathcal J}
\newcommand{\Kk}{\mathcal K}
\newcommand{\del}{\partial}
\newcommand{\Jj}{J_j}
\newcommand{\ev}{\mathrm{ev}_0}
\newcommand{\Norm}[1]{\left\lVert #1\right\rVert}
\newcommand{\Abs}[1]{\left\lvert #1\right\rvert}
\newcommand{\ip}[2]{\left\langle #1,#2\right\rangle}
\newcommand{\norm}[1]{\left\lVert #1\right\rVert}
\renewcommand{\d}[1]{\mathrm{d}#1}
\newcommand{\g}{\mathfrak{g}}
\newcommand{\st}{\; \big| \;}
\newcommand{\PP}{\mathcal{P}}
\newcommand{\HH}{\mathbb{H}}
\newcommand{\PR}{\mathcal{P}_d^{\mathbb{R}}}
\newcommand{\var}{\text{var}}
\newcommand{\cont}{\text{cont}}
\newcommand{\rad}{\text{rad}}

\DeclareMathOperator{\re}{Re}
\DeclareMathOperator{\im}{Im}
\DeclareMathOperator{\diag}{diag}
\DeclareMathOperator{\Span}{Span}
\DeclareMathOperator{\coker}{Coker}
\DeclareMathOperator{\ind}{ind}
\DeclareMathOperator{\Spec}{Spec}
\DeclareMathOperator{\Crit}{Crit}
\DeclareMathOperator{\Lie}{Lie}
\DeclareMathOperator{\grad}{grad}
\DeclareMathOperator{\Dens}{Dens}


\title[Heat Kernel and Resurgence]{Heat Kernel and Resurgence}

\author[Si Li]{Si Li}
\author[Yong Li]{Yong Li}
\author[Xinxing Tang]{Xinxing Tang}

\thanks{S.~Li: Yau Mathematical Sciences Center, Tsinghua University, Beijing, China.
Email: \texttt{sili@mail.tsinghua.edu.cn}.}

\thanks{Y.~Li: Beijing Institute of Mathematical Sciences and Applications,
Beijing, China.
Email: \texttt{liyong@bimsa.cn}.}

\thanks{X.~Tang: Beijing Institute of Mathematical Sciences and Applications,
Beijing, China.
Email: \texttt{tangxinxing@bimsa.cn}.}

\date{}

\keywords{}


\begin{document}


\begin{abstract}
We study the resurgent structure of short-time heat kernel asymptotics from the viewpoint of Picard–Lefschetz theory. For a real analytic Riemannian manifold, we show the heat kernel admits a $1$--Gevrey small-time expansion whose Borel transform detects complex-geometric data beyond the real geodesic sector. We formulate an infinite-dimensional Picard–Lefschetz problem of Morse–Floer type for the holomorphic energy functional on the complexified path space, and propose a heat-kernel analogue of the Picard–Lefschetz/Alien correspondence. In this framework, pointed alien operators acting on the asymptotic expansion associated with the real geodesic are predicted to produce the formal heat-kernel sectors associated with other holomorphic geodesics, with coefficients given by signed counts of connecting trajectories of the Morse flow. We perform a confirming test of this proposal on the hyperbolic plane $\mathbb{H}_{\R}^2$.
\end{abstract}
\maketitle
\tableofcontents

\section{Introduction}

The purpose of this paper is to explain how the resurgent structure of a heat-kernel expansion is reflected by Picard--Lefschetz theory on an infinite-dimensional complexified path space. This is motivated by the Picard–Lefschetz/Alien correspondence in the finite-dimensional case. We start by reviewing this correspondence following \cite{LiLiTang2026}  and refer to it for further details and references therein. 

\subsection{Lefschetz thimble and Borel singularity}\label{sec:intro-Borel}
Consider a finite-dimensional exponential integral
\begin{equation*}
I_\Gamma(\hbar)=\int_\Gamma e^{-S(z)/\hbar}\,\mu
\end{equation*}
where $S(z)$ is a holomorphic function on a complex manifold, $\mu$ is a holomorphic form, and $\Gamma$ is a contour (or cycle) in a suitable relative homology class $[\Gamma]$. The asymptotic behavior of $I_\Gamma(\hbar)$ for small $\hbar$ is controlled by the critical points of the phase function
$$\dd S=0.$$

Assume $S$ has only non-degenerate isolated critical points and the corresponding critical values are pairwise distinct. For each critical point $p$, there corresponds an integration cycle $\mathcal J_p$, called the Lefschetz thimble, which can be constructed by the downward gradient flow of $\re(S/\hbar)$ (with a chosen Hermitian metric on $X$). One can decompose the cycle in homology class
$$
[\Gamma]=\sum_p n_p [\mathcal J_p]
$$
and obtain a representation of the integral by
$$
I_\Gamma(\hbar)=\sum_{p}n_p \int_{\mathcal J_p} e^{-S(z)/\hbar}\,\mu.
$$
Near each critical point $p$, one obtains a formal asymptotic expansion of the form 
$$
I_p(\hbar) :=\int_{\mathcal J_p} e^{-S(z)/\hbar}\,\mu \sim e^{-S(p)/\hbar}\tilde \phi_p(\hbar)= e^{-S(p)/\hbar}\hbar^{n/2}\left(c_{p,0}+c_{p,1}\hbar+c_{p,2}\hbar^2+\cdots\right)\quad \text{as}\ \hbar\to 0 
$$
uniformly on closed subsectors of a sufficiently small sector in the $\hbar$-plane. This is the standard local steepest-descent (or stationary-phase) expansion of a Laplace-type integral near a non-degenerate critical point (see e.g. \cite{BleisteinHandelsman,PhamSaddlepoint}); equivalently, it may be viewed as the local WKB expansion associated with the saddle $p$.
This expansion is typically divergent and only asymptotic in nature. They all together capture the asymptotic behavior of $I_\Gamma(\hbar)$ for small $\hbar$. 

The theory of resurgence developed by \'Ecalle \cite{Ecalle1981I,Ecalle1981II,Ecalle1985III} provides a general framework to study global analytic natures of divergent asymptotic expansions. Under appropriate growth condition, the Borel transform $\hat \phi_p(\xi)$ of the $\hbar$-asymptotic series $\tilde \phi_p(\hbar)$ from the critical point $p$ will be an integrable germ locally near $\xi=0$ and admits analytic continuation beyond it. For $\arg \hbar = \theta$,  one has a resummed function with the expected asymptotic expansion
\[
I_p(\hbar)
:=
\int_{\mathcal J_p} e^{-S(z)/\hbar}\,\mu
=
e^{-S(p)/\hbar}\,
\mathcal L_\theta\widehat\phi_p(\hbar) \sim e^{-S(p)/\hbar}
\widetilde\phi_p(\hbar),
\qquad \hbar\to0.
\]
Here $\mathcal L_\theta$ is the Laplace tranform along the ray  $e^{i\theta}\R_{\geq 0}$ (see Definition \ref{defn:Laplace}). %If the direction $e^{i\theta}\R_{\geq 0}$ is singular, the same statement is understood in the corresponding lateral sense. 
The Borel transform $\hat \phi_p(\xi)$ may exhibit singularities in the $\xi$-plane. Such Borel singularities lead to possible ambiguities in the resummation. One has to choose a direction $e^{i\theta}$ so that the ray in the Laplace transform avoids Borel singularities. Two different directions could give rise to different analytic objects when one ray has to pass certain singularities to reach the other. This is precisely the Stokes phenomenon. 

\begin{figure}[H]
\centering
\begin{tikzpicture}[scale=0.7,>=Latex]

\tikzset{
  stokesray/.style={
    gray!70,
    dashed,
    line width=1pt
  },
  lateralsum/.style={
    black,
    line width=1.1pt,
    -{Latex[length=2.2mm]}
  },
  stokesarcminus/.style={
    black,
    line width=1.05pt,
    -{Latex[length=2.2mm]}
  },
  stokesarcplus/.style={
    black,
    line width=1.05pt,
    -{Latex[length=2.2mm]}
  }
}
\coordinate (O) at (0,0);
\draw[stokesray] (O) -- (30:7.2);
\draw[lateralsum] (O) -- (20:7.0);
\draw[lateralsum] (O) -- (40:7.0);
\filldraw[black] (O) circle (1.3pt);
\node[below left] at (O) {$0$};
%\node[below right] at (6.8,0.25) {$\xi$-plane};
%\node[gray!75!black] at (30:7.05) {$d$};
\foreach \r in {2.0,3.6,5.2,6.5}{
  \draw[red, line width=1pt]
    ($(30:\r)+(-0.11,-0.11)$) -- ($(30:\r)+(0.11,0.11)$);
  \draw[red, line width=1pt]
    ($(30:\r)+(-0.11,0.11)$) -- ($(30:\r)+(0.11,-0.11)$);
}
\node at (20:7.15) [below right] {$\theta_2$};
\node at (40:7.15) [above left] {$\theta_1$};
%\node at (20:3.45) [below right] {$\mathcal S^{<}$};
%\node at (40:3.55) [above left] {$\mathcal S^{>}$};
%\draw[stokesarcminus]
  (25:8.35) arc[start angle=25,end angle=35,radius=8.35];
%\node at (30:8.75) {$\mathfrak S^{-}$};
%\draw[stokesarcplus]
  (35:8.2) arc[start angle=35,end angle=25,radius=7.65];
%\node at (30:7.85) {$\mathfrak S^{+}$};
\end{tikzpicture}
\caption{$\xi$-plane}
%\caption{Stokes ray $d$ in the Borel $\zeta$-plane, its neighboring lateral directions $d^{<}$ and $d^{>}$, and the corresponding Stokes automorphisms $\mathfrak S^\pm$.}
\label{fig:stokes-automorphism-zeta-plane}
\end{figure}

For the exponential integral, the ambiguity of the ray for Laplace transform is related to the ambiguity of the choice of Lefschetz thimble. Take $\hbar$ to be valued in $\C^\times$ and write
$$
\hbar=|\hbar|e^{i\theta}.
$$
Then the flow of $\re(S/\hbar)$ is the same as that of $\re(e^{-i\theta}S)$. The magic of holomorphic Morse function $S$ tells that along the  the downward gradient flow of $\re(S/\hbar)$ :
\begin{align*}
&(1)~ \re(e^{-i\theta} S)\ \text{is {decreasing}}, &\qquad (2)~ \im (e^{-i\theta} S) \ \text{is {constant}}.
\end{align*}
The corresponding Lefschetz thimble is represented by (here we specify its $\theta$-dependence)
$$
\mathcal J_p^\theta=\{z\in X\mid \text{$z$ flows to $p$ at time $t=+\infty$}\}.
$$


With our assumptions on $S$, at generic $\theta$, the values of $\im(e^{-i\theta}S)$ at critical points are all distinct, so all such constructed thimbles are disjoint. If we sit at some $\theta_*$ such that
$$
\im(e^{-i\theta_*}S(p))=\im(e^{-i\theta_*}S(q))
$$
for two critical points $p,q$, then one could have flow trajectories that connect one critical point to another.  Such $\theta_*$ is called a Stokes phase and the ray in the direction $e^{i\theta_*}$ is called a Stokes ray. 

Assume we have a trajectory that flows from $q$ to $p$ at a Stokes phase $\theta_*$ (this requires $\re(e^{-i\theta_*}S(p))<\re(e^{-i\theta_*}S(q))$ by the downward gradient flow). Then the Borel transform $\hat \phi_p(\xi)$ will see a singularity in the direction $e^{i\theta_*}$ at 
$$
\omega=S(q)-S(p).
$$
Thus Borel singularities correspond to relative critical values with respect to a reference critical point.

In  \cite{KS-2022,KS-2024}, Kontsevich and Soibelman investigated analyticity and resurgence phenomena in the context of Lefschetz thimbles and their wall-crossing behavior, and develop a Floer-theoretic framework for exponential integrals. Our study of the complexified path space in this work follows the perspective of holomorphic Morse–Floer theory and connects it to the alien calculus, as we explain next.


\subsection{Connecting trajectory and Alien operator}\label{sec:intro-Alien} Borel singularities lead to rich algebraic structures on the global analytic nature of divergence series. They are captured by the so-called Alien operators.  

Let $\widetilde\phi(\hbar)$ 
be a $1$--Gevrey formal series (see Section \ref{sec:Gevrey-Borel}). Let
\[
\widehat\phi(\xi)=\mathcal B\widetilde\phi(\xi)
\]
be its Borel transform.  Suppose that $\widehat\phi$ admits analytic
continuation from the origin to a singular point $\omega\in\C^*$, and the germ is integrable at $\omega$. Then we can obtain a formal series in $\hbar$, denoted by  $$
\Delta_{\omega}^{+}\widetilde\phi
$$
by extracting the local
variation of the Borel transform at $\omega$, recentering it at the origin,
and applying inverse Borel transform. See Section \ref{sec:heat-alien} for the precise definition and references. 



$\Delta_{\omega}^{+}$ is called the positive alien operator. The corresponding pointed alien operator is
\begin{equation*}%\label{eq:pointed-positive-alien-operator}
\dot\Delta_{\omega}^{+}
:=
e^{-\omega/\hbar}\Delta_{\omega}^{+}.
\end{equation*}
The exponential factor restores the action shift carried by the Borel
singularity.  For instance, if a formal solution with exponential weight
$e^{-\Phi_0/\hbar}$ has a Borel singularity at
\[
\omega=\Phi_1-\Phi_0,
\]
then $\dot\Delta_{\omega}^{+}$ produces a contribution with weight
\[
e^{-\Phi_1/\hbar}.
\]

The finite-dimensional analysis of \cite{LiLiTang2026} provides the guiding
principle for the comparison below.  There, the pointed alien operator acts on
the formal saddle expansion $\widetilde I_p$ attached to a critical point
$p$, and produces the expansion $\widetilde I_q$ attached to another
critical point $q$ when
\[
\omega=S(q)-S(p)
\]
is a Borel singularity.  More precisely, if
\[
\dot\Delta_{\omega}^{+}\,\widetilde I_p
=
N_{pq}\,\widetilde I_q,
\]
then $N_{pq}$ agrees with the corresponding Picard--Lefschetz
wall-crossing coefficient; in elementary situations, this is the signed
count of direct connecting trajectories at the relevant Stokes phase:
$$
N_{pq}=\text{signed}\ \sharp \{\text{flow}\ \gamma(t): \gamma(-\infty)=q, \gamma(+\infty)=p\}.
$$
This identification is established in \cite{LiLiTang2026} by comparing the
Stokes automorphism on the resurgent side with the thimble jump formula on
the Picard--Lefschetz side.

Here is an example for illustration. Consider $X=\C^*$ and the exponential integral 
\[
\int_{\Gamma}e^{-\frac{1}{2\hbar}\left(z-\frac{1}{z}\right)}\frac{\dd z}{z}.
\]
There are two critical points $p_{\pm}=\pm i$, which lead to two asymptotic series $ I_{\pm}(\hbar)$. The Stokes phases appear at $\theta_*=\frac{\pi}{2}\ \text{or}\ \frac{3\pi}{2}$. 
\begin{figure}[H]
\centering

\begin{tikzpicture}[>=Latex, scale=0.95]

\tikzset{
  cyl/.style={
    gray!65,
    line width=0.9pt
  },
  cylfill/.style={
    fill=gray!10,
    draw=none
  },
  waist/.style={
    gray!75,
    densely dashed,
    line width=0.9pt
  },
  crit/.style={
    circle,
    fill=black,
    inner sep=1.5pt
  },
  Jcurve/.style={
    black,
    line width=1.25pt,
    postaction={decorate},
    decoration={
      markings,
      mark=at position 0.75 with {\arrow{Latex[length=2.2mm]}}
    }
  },
  JcurveBack/.style={
    black,
    line width=1.25pt,
    dashed,
    postaction={decorate},
    decoration={
      markings,
      mark=at position 0.75 with {\arrow{Latex[length=2.2mm]}}
    }
  },
  traj/.style={
    blue!75!black,
    line width=2.1pt,
    postaction={decorate},
    decoration={
      markings,
      mark=at position 0.82 with {\arrow{Latex[length=2.4mm]}}
    }
  },
  trajback/.style={
    blue!75!black,
    line width=2.1pt,
    dashed,
    postaction={decorate},
    decoration={
      markings,
      mark=at position 0.82 with {\arrow{Latex[length=2.4mm]}}
    }
  }
}
\def\rx{1.15}
\def\ry{0.33}
\def\H{2.25}

\newcommand{\drawCylinderBase}{
\fill[cylfill] (-\rx,-\H) rectangle (\rx,\H);

\draw[cyl,dashed] (-\rx,\H)
    .. controls (-0.55,\H+\ry) and (0.55,\H+\ry) ..
    (\rx,\H);
\draw[cyl,dashed] (-\rx,-\H)
    .. controls (-0.55,-\H+\ry) and (0.55,-\H+\ry) ..
    (\rx,-\H);

\draw[cyl] (-\rx,-\H) -- (-\rx,\H);
\draw[cyl] (\rx,-\H) -- (\rx,\H);

\draw[cyl] (-\rx,\H)
    .. controls (-0.55,\H-\ry) and (0.55,\H-\ry) ..
    (\rx,\H);
\draw[cyl] (-\rx,-\H)
    .. controls (-0.55,-\H-\ry) and (0.55,-\H-\ry) ..
    (\rx,-\H);
\draw[waist] (0,0) ellipse[x radius=\rx, y radius=\ry];

\node[crit] (wp) at (-\rx,0) {};
\node[crit] (wm) at (\rx,0) {};
\node[left] at (-\rx-0.08,0.12) {$p_+$};
\node[right] at (\rx+0.08,0.12) {$p_-$};
}

\begin{scope}[shift={(-4.2,0)}]
  \node at (0,\H+0.85) {$\theta=\frac\pi2-\delta$};
  \drawCylinderBase

\draw[Jcurve]
(\rx,0) .. controls  (0.5,-0.3) and  (-0.5,-0.2)..
    (-0.85,0.1)
    --
    (-0.85,\H-0.1);
\draw[JcurveBack]
    (-0.85,-\H+0.1)
    --
    (-0.85,-0.1)
    .. controls (-0.5,0.2) and (0.5,0.3) ..
    (\rx,0);
\draw[Jcurve]
    (-\rx,-\H)--(-\rx,\H);
  \node[black] at (-1.55,1.35) {$\mathcal J_+^{<}$};
  \node[black] at (-0.5,1.35) {$\mathcal J_-^{<}$};
\end{scope}

\begin{scope}[shift={(0,0)}]
  \node at (0,\H+0.85) {$\theta_*=\frac{\pi}{2}$};
  \drawCylinderBase

\draw[traj]
    (-\rx,0)
    .. controls (-0.70,-0.42) and (0.70,-0.42) ..
    (\rx,0);
\draw[Jcurve]
    (-\rx,-\H)--(-\rx,\H);
\draw[trajback]
    (-\rx,0)
    .. controls (-0.70,0.42) and (0.70,0.42) ..
    (\rx,0);
\end{scope}

\begin{scope}[shift={(4.2,0)}]
  \node at (0,\H+0.85) {$\theta=\frac\pi2+\delta$};
  \drawCylinderBase
\draw[Jcurve]
(\rx,0) .. controls  (0.5,-0.2) and  (-0.5,-0.3)..
    (-0.85,-0.1)
    --
    (-0.85,-\H-0.1);
\draw[JcurveBack]
    (-0.85,\H+0.1)
    --
    (-0.85,0.1)
    .. controls (-0.5,0.2) and (0.5,0.3) ..
    (\rx,0);
\draw[Jcurve]
    (-\rx,-\H)--(-\rx,\H);
  \node[black] at (-1.55,1.35) {$\mathcal J_+^{>}$};
\node[black] at (-0.5,-1.35) {$\mathcal J_-^{>}$};
\end{scope}

\end{tikzpicture}

\end{figure}

At $\theta_*=\frac{\pi}{2}$, there are two connecting trajectories from $p_+$ to $p_-$. The two Lefschetz thimbles exhibit topological change when crossing the Stokes phase. On the other hand, explicit computation via analytic continuation shows
$$
 \dot\Delta_{\omega}^{+}\widetilde I_-=
2\,\widetilde I_+.
$$
The coefficient $2$ coincides precisely with the two connecting trajectories above. See \cite{LiLiTang2026} for details. 

\subsection{Holomorphic Morse-Floer Theory}\label{sec:intro-MF} In this paper, we formulate an infinite-dimensional Picard–Lefschetz problem of Morse–Floer type following Witten's idea of analytic continuation of quantum mechanics \cite{Witten2010}. In \cite{Witten2010}, the path integral is formulated on the phase space, and an appropriate complexification of the target is a hyperK\"{a}hler manifold. Our formulation is  instead modeled on the path integral over the configuration space, which could be understood as the adiabatic limit of the hyperK\"ahler construction.


Let $(M,g)$ be a Riemannian manifold. Fix $\tau>0$ and $a,b\in M$. Consider the Sobolev path space
\[
\PP_{a,b}(M):=
\bigl\{\gamma\in H^{1}([0,\tau],M)\,:\,\gamma(0)=a,\ \gamma(\tau)=b\bigr\}
\]
equipped with the energy functional (the $\frac{1}{4}$ convention is for convenience of resurgence analysis)
\[
\mathcal E(\gamma):=\frac{1}{4}\int_{0}^{\tau}g(\dot\gamma(t),\dot{\gamma}(t))\,\dd t.
\]
Critical points of $\mathcal E$ correspond to geodesics in $M$ from $a$ to $b$. 

We are interested in Picard–Lefschetz theory on path spaces. To set things up, we need a holomorphic structure. Assume $(M,g)$ is real analytic and extends to {an appropriate} complexification  $(M_\C,g_{\C})$ where $M_{\C}$ is a complex manifold, $g_{\C}$ is a holomorphic metric such that $M$ is a connected component of the real locus of $M_{\C}$ and $g$ is the restriction of $g_{\C}$. We discuss a few related constructions of complexification in Section \ref{sec: complexification}. Then the natural candidate of our compexified path space is 
\[
\mathcal P_{a,b}(M_{\C})
=\left\{\gamma\in H^1([0,\tau],M_{\C}) :
\gamma(0)=a,\ \gamma(\tau)=b\right\}
\]
which inherits a complex structure from that of $M_{\C}$.  The complexified energy by
\[
\mathcal E_{\C}(\gamma)
:=
\frac{1}{4}\int_0^\tau
g_{\C}(\dot\gamma(t),\dot\gamma(t))\,\dd t
\]
can be viewed as a holomorphic function on the infinite-dimensional complex manifold $\mathcal P_{a,b}(M_{\C})$. 

Note that if we normalize the path by re-parametrization 
$$
\tilde \gamma(t):=\gamma(\tau t): [0,1]\to M_{\C},
$$
then 
$$
\mathcal E_{\C}(\gamma)=\frac{1}{4\tau}\int_0^1
g_{\C}(\dot{\tilde\gamma}(t),\dot{\tilde \gamma}(t))\,\dd t.
$$
Thus $\tau$ plays the analogous role of $\hbar$. Motivated by Feynman-Kac formula, Picard–Lefschetz theory of such type is related to the resurgence problem of heat kernel at time $\tau$. The following theorem on the Borel summability of heat kernel at small time provides a first evidence. 

\begin{thm}[Theorem \ref{thm:local-borel-summability}]
\label{thm:local-borel-summability-introduction}
Let $(M,g)$ be real analytic and $\Delta$ be the  Laplace--Beltrami
operator acting on functions
\[
  \Delta f = - \mathrm{div}(\nabla f).
\]
Let $V$ be a sufficiently small compact
convex neighborhood with the heat kernel asymptotic expansion 
\[
K(\tau;x,y)\sim (4\pi \tau)^{-n/2}e^{-d(x,y)^2/(4\tau)}\sum_{j\ge0}u_j(x,y)\tau^j, \quad (x,y)\in V\times V. 
\]
Then the coefficient series in time $\tau$
\[
\widetilde u(\tau;x,y)=\sum_{j\ge0}u_j(x,y)\tau^j
\]
is $1$--Gevrey, uniformly for $(x,y)\in V'\times V'$ with
$V'\Subset V$.  Hence its Borel transform
\[
\widehat u(\xi;x,y)=\sum_{j\ge0}\frac{u_j(x,y)}{j!}\xi^j\quad \text{converges near $\xi=0$.  }
\]
\end{thm}

We prove Theorem \ref{thm:local-borel-summability-introduction} in Section \ref{subsec:proof-borel-summability}. It allows us to do resurgence analysis in the plane of Borel transform of time. In particular, the analytic continuation of the Borel transform $\widehat u(\xi;a,b)$ is expected to reflect the Picard–Lefschetz geometry of
$$
\tau\mathcal E_{\C}: \PP_{a,b}(M_{\C})\to \C. 
$$

The first surprise comes from the critical points. Standard calculus of variation shows that critical points of $\mathcal E_{\C}$ are holomorphic geodesics in $M_{\C}$ from $a$ to $b$ (see Section \ref{sec:PL-problem-path}). They contain usual real geodesics inside $M$, as well as possible new holomorphic geodesics that go out of $M$ into $M_{\C}$ before reaching $b$. Let $a,b\in M$ be close and such that there exists a unique minimal geodesic in $M$ from $a$ to $b$.  The heat kernel $K(\tau;a,b)$ is constructed purely from the data on $M$, but the Borel transform $\widehat u(\xi;a,b)$ will detect geometric information on $M_{\C}$ that is invisible from $M$. For example, consider the set of Borel singularities of $\widehat u(\xi;a,b)$. As explained in Section \ref{sec:intro-Borel}, they correspond to critical values of $\mathcal E_{\C}$ relative to the energy of the minimal geodesic, and we shall observe those new holomorphic geodesics from $\widehat u(\xi;a,b)$. This can be viewed as the Picard–Lefschetz analogue of the quantum tunneling phenomenon.


As an example, consider the two dimensional hyperboloid space
\[
\HH_{\R}^2=\{-x_0^2+x_1^2+x_2^2=-1\mid x_i\in\R, x_0>0\}, \quad \dd s^2=-\dd x_0^2+\dd x_1^2+\dd x_2^2\big|_{\HH_{\R}^2}.
\] 
It has a natural complexification 
\[
\HH_{\C}^2:=\bigl\{z=(z_0,z_1,z_2)\in \C^3:\ -z_0^2+z_1^2+z_2^2=-1\bigr\}, \quad \dd s^2_{\C}=-\dd z_0^2+\dd z_1^2+\dd z_2^2\big|_{\HH_{\C}^2}.
\]
Consider two end-points  
$$
a=(1,0,0),~ b=(\cosh d,\sinh d,0) \in \HH_{\R}^2,
$$ 
where $d$ is the real geodesic distance between $a$ and $b$. Then the set of all holomorphic geodesics in $\HH_{\C}^2$ from $a$ to $b$ is indexed by $k\in \Z$ and explicitly given by (see Example \ref{eg:H^2-PL})
\begin{equation*}
\gamma_{{k},d}(t)=\bigl(\cosh(c_{k,d}t),\,\sinh(c_{k,d}t),\,0\bigr),
\qquad 0\le t\le \tau
\end{equation*}
where $c_{k,d}:=\frac{d+2\pi i k}{\tau}$. Only $\gamma_{0,d}(t)$ lies in $\HH_{\R}^2$ and all the other $\gamma_{{k},d}(t)$'s are invisible from $\HH_{\R}^2$. The critical values of the energy functional are 
$$
\tau\mathcal E_{\C}(\gamma_{{k},d})=\frac{(d+2\pi i k)^2}{4}.
$$ 

The exact heat kernel formula on $\HH^2_{\R}$ is known by
\begin{equation*}
K_{\HH_{\R}^{2}}(\tau;a,b)
=
\frac{\sqrt{2}e^{-\tau/4}}{(4\pi \tau)^{3/2}}\int_{d^2/4}^{\infty}\frac{e^{-w/\tau}}{\sqrt{\cosh\sqrt{4w}-\cosh d}}~\dd w
\end{equation*}
which has already the form of Laplace-transform. In fact, the Borel coordinate is related to $w$ by 
$$
\xi=w-d^2/4. 
$$
The Borel singularities occur at
$$
{\cosh\sqrt{4w}}=\cosh d, \quad \text{i.e.}\quad \xi=\frac{(d+2\pi i k)^2}{4}-\frac{d^2}{4}, \quad k\in \Z.
$$
They correspond precisely to the critical values of $\tau\mathcal E_{\C}$ relative to the real geodesic $\gamma_{{0},d}$
$$
\tau\mathcal E_{\C}(\gamma_{{k},d})-\tau\mathcal E_{\C}(\gamma_{{0},d})
$$
as expected. This example was explained to us by Kontsevich \cite{Kont,Kont2020}. We make precise the following Conjecture \ref{conj:singularity} which is essentially due to Kontsevich. 

\begin{conj}[\cite{Kont, Kont2020}]\label{conj:singularity} Singularities of the Borel transform of the heat kernel $K(\tau;a,b)$ correspond to (relative) critical values of the complexified energy functional $\tau \mathcal E_{\C}$. 
\end{conj}

The Borel transform of the  short-time heat kernel asymptotics in Conjecture \ref{conj:singularity}  is established by Theorem \ref{thm:local-borel-summability-introduction}. We prove a version of Conjecture \ref{conj:singularity} in Theorem \ref{thm:Borel-singularity-complex-geodesic} 
under the assumption of its WKB asymptotic expansion along analytic continuation.

We next consider the infinite-dimensional analogue of Picard–Lefschetz/Alien correspondence as explained in Section \ref{sec:intro-Alien}. We start with the Picard–Lefschetz side. Let us choose a hermitian metric $h$ on $M_\C$ which induces a $L^2$-hermitian inner product $<\cdot,\cdot>_\hbar$ on $\mathcal P_{a,b}(M_{\C})$ by 
$$
<\delta \gamma_1, \delta \gamma_2>_h:=\int_0^\tau h(\delta \gamma_1, \delta \gamma_2) dt, \qquad \delta \gamma_1, \delta \gamma_2 \in T_\gamma \mathcal P_{a,b}(M_{\C}).
$$
Let $\theta$ be a chosen phase. Consider the downward gradient flow of 
$$
\re(e^{-i\theta}\mathcal E_\C). 
$$
It can be described by a map 
\[
\gamma(t, s): [0,\tau] \times [0,+\infty) \to M_{\C},\qquad \gamma(0,s)=a, \ \gamma(\tau,s)=b
\]
satisfying the Morse flow equation  (see Section \ref{sec:PL-problem-path} for details)
\begin{equation}\label{eqn:morse-flow}
\partial_s\gamma=-\nabla_{L^2}\re(e^{-i\theta}\mathcal E_\C)(\gamma) \tag{$\dagger$}.
\end{equation}


Let us next consider the resurgence side. Assume $a,b$ are chosen close enough such that the heat kernel $K(\tau;a,b)$ has an asymptotic expansion in time $\tau$
$$
K(\tau;a,b)\sim \widetilde K_0= (4\pi \tau)^{-n/2}e^{-d_0^2/(4\tau)}\sum_{j\ge0}u_{0,j}\tau^j
$$
which is Borel summable by Theorem \ref{thm:local-borel-summability-introduction}. Here $d_0$ is the length of the minimal real geodesic $\gamma_0$ in $M$ from $a$ to $b$, and $u_0(a,a)=1$ by the standard heat kernel normalization. 

Let $\gamma_1$ be another holomorphic geodesic in $M_{\C}$ from $a$ to $b$, and the corresponding (complexified) length is $d_1$. By the above discussion, we expect a singularity of the Borel transform of $\tilde K_0$ at 
$$
\omega= \tau \mathcal E(\gamma_1)- \tau \mathcal E(\gamma_0)=\frac{d_1^2}{4}-\frac{d_0^2}{4}. 
$$
Let 
$$
\widetilde K_1= (4\pi \tau)^{-n/2}e^{-d_1^2/(4\tau)}\sum_{j\ge0}v_{1,j}\tau^j
$$
be the formal solution of the heat equation with $v_{1,0}$ normalized by the Jacobian factor along the holomorphic geodesic $\gamma_1$. This is obtained similarly to the standard heat kernel $\tilde K_0$ by solving the transport equation recursively but with the new phase factor $e^{-d_1^2/(4\tau)}$ (see Theorem \ref{thm:heat-kernel-expansion-complex-geodesic}). Such $\widetilde K_1$ is expected to be the asymptotic expansion of certain path integral over the Lefschetz thimble associated to $\gamma_1$. 

Based on the finite-dimensional Picard–Lefschetz/Alien correspondence, we conjecture the following result on the pointed Alien derivative of the heat kernel asymptotic expansion (see also Conjecture \ref{conj:alien-PL-integrality-heat}). 

\begin{conj}[Picard–Lefschetz/Alien correspondence for heat kernel]\label{conj:alien} The pointed Alien derivative of $\widetilde K_0$ at the singularity $\omega$ is 
$$
\dot\Delta_{\omega}^{+} \widetilde K_0=N \widetilde K_1
$$
where $N$ is the appropriate signed counting of solutions of 
$$
\gamma: [0,\tau]\times (-\infty,\infty)\to M_{\C}
$$
satisfying 
\begin{enumerate}
\item the downward gradient Morse flow equation \eqref{eqn:morse-flow} at the Stokes phase $\theta=\theta_*$ where
$$
\im(e^{-i\theta_*}\mathcal E_\C(\gamma_1))=\im(e^{-i\theta_*}\mathcal E_\C(\gamma_0)), \qquad \re(e^{-i\theta_*}\mathcal E_\C(\gamma_1))>\re(e^{-i\theta_*}\mathcal E_\C(\gamma_0))
$$
\item the boundary condition 
$$
\gamma(0,s)=a, \qquad \gamma(\tau,s)=b, \qquad \forall\,s\in\R
$$
\item the initial and final condition 
\[\lim_{s\rightarrow-\infty}\|\gamma(\cdot,s)-\gamma_0\|_{H_t^1}=0,\qquad \lim_{s\rightarrow+\infty}\|\gamma(\cdot,s)-\gamma_1\|_{H_t^1}=0.\]
\end{enumerate} 

\end{conj}

In the second part of this paper, we perform a concrete test of Conjecture \ref{conj:alien} on the hyperbolic plane $\HH_{\R}^2$. For each $k\in \Z$ we have a critical path $\gamma_{{k},d}$ and the corresponding Borel singularity 
$$
\omega_k=\frac{(d+2\pi i k)^2}{4}-\frac{d^2}{4}. 
$$
In Section \ref{sec:H^2-resurgence}, we compute explicitly each asymptotic series $\widetilde K_k$ and show that
\[
\dot\Delta_{\omega_k}^{+}\,\widetilde{K}_0
=2\,\widetilde{K}_k.
\]
By Conjecture \ref{conj:alien}, this predicts two connecting trajectories from $\gamma_k$ to $\gamma_0$ at the corresponding Stokes phase. In Section \ref{sec:H^2-flow}, we prove this prediction for $k=1$ and $d$ sufficiently small by solving the Morse flow equation via perturbation method. This furnishes a nontrivial test of Conjecture \ref{conj:alien}. 

\smallskip
\subsubsection*{Acknowledgments}

The authors would like to thank Maxim Kontsevich, Ryszard Nest, David Sauzin, Shanzhong Sun, Yifan Wu, Junrong Yan, Shing-Tung Yau for helpful discussions. S.~L. was supported by NSFC No.12571068.  Y.~L. was partially supported by BNSFC No.JR25001.  X.~T. was supported by BNSFC No.JR25001, BNSFC Youth Project Youth Project No.1254041 and NSFC Youth Project No.12501079. Part of this work was done while S.~L. was a visiting fellow of the International Centre for Mathematical Sciences  at Edinburgh in spring 2026, and he thanks the institute for its hospitality and provision of excellent working environment.

\part{Theory}

\section{Heat Kernel and Resurgence}\label{sec:heat-resurgence}

In this section, we establish the Borel-summability of heat kernel short-time asymptotic series on real analytic metric space. This allows us to do resurgence analysis in the plane of Borel transform of time, and set up the Picard–Lefschetz problem of Morse–Floer type on path spaces. 

\subsection{Borel-Laplace transform}\label{sec:Gevrey-Borel}
We will briefly review the Gevrey asymptotics and Borel-Laplace transform that will be used in our resurgence analysis. For more details, see \cite{Balser,Costin,Ecalle1981I,Ecalle1981II,Ecalle1985III,MitschiSauzin2016,SauzinSplitting}.
\begin{defn}\label{defn:series}
Let
$\widetilde f(\hbar)
=
\sum\limits_{m\ge0}a_m\hbar^m
\in \C[[\hbar]]$. Let $f(\hbar)$ be a holomorphic function defined in a sector
$$
\text{S}^{\theta_1,\theta_2,r}
=\left\{\hbar\in\C
\mid
0<|\hbar|<r,
\theta_1<\arg\hbar<\theta_2
\right\}. 
$$
\begin{itemize}
\item We say that $f$ admits $\widetilde f$ as an asymptotic
expansion in $\text{S}^{\theta_1,\theta_2,r}$, and write
\begin{equation*}
f(\hbar)\sim \widetilde f(\hbar),
\qquad
\hbar\to0,\quad \hbar\in\text{S}^{\theta_1,\theta_2,r},
\end{equation*}
if there exists constants $C_{N,\text{S}'}>0$ such that
\begin{equation}\label{eq:definitionasymptotic}
\left|
f(\hbar)
-
\sum_{m=0}^{N-1}a_m\hbar^m
\right|
\le
C_{N,\text{S}'}|\hbar|^N,
\qquad
\hbar\in \text{S}' =\overline{\text{S}^{\theta_1',\theta_2',r'}},
\end{equation}  for every
$N\ge0$ and for every closed subsector $\overline{\text{S}^{\theta_1',\theta_2',r'}}\Subset\text{S}^{\theta_1,\theta_2,r}$ ($0<r'<r,
\theta_1<\theta_1'<\theta_2'<\theta_2$). 

\item We say that $\widetilde f$ is a $1$--Gevrey series if
there are constants $C,A>0$ such that 
$$
|a_m|\le C A^m m!\quad \text{for} \quad m\ge0. 
$$
It is denoted by $\widetilde f\in\C[[\hbar]]_1$.
\item We say that $f$ admits $\widetilde f$ as a $1$--Gevrey asymptotic
expansion in $\text{S}$ if there are constants $C,A>0$ such that
$C_{N,\text{S}'}= C A^N N!$ in \eqref{eq:definitionasymptotic}.
\end{itemize}
\end{defn}

\begin{defn}
Let $\alpha\in\C$ with $\operatorname{Re}\alpha>0$. For a formal series $\widetilde g(\hbar)
=
\hbar^\alpha \widetilde f(\hbar)
=
\sum\limits_{m\ge0}a_m\hbar^{m+\alpha}
\in \hbar^\alpha\C[[\hbar]]$,
we define its Borel transform by
\begin{equation*}
\widehat g(\xi)
:=
\mathcal B\widetilde g(\xi)
=
\sum_{m\ge0}
\frac{a_m}{\Gamma(m+\alpha)}
\xi^{m+\alpha-1}
\in
\xi^{\alpha-1}\C[[\xi]].
\end{equation*}
If $\widetilde f \in\C[[\hbar]]_1$, then the above formal
series converges near origin, and the condition
$\operatorname{Re}\alpha>0$ ensures that $\widehat g$ is locally integrable
at $\xi=0$. Throughout
this definition, fractional powers of $\xi$ are defined using a principle branch
of $\log\xi$. We also use the convention
\begin{equation*}
\mathcal B(a_0)=a_0\delta_0,
\end{equation*}
where $\delta_0$ denotes the Dirac delta distribution at the origin in the
Borel plane. With this convention, the Borel transform extends to the case
$\alpha=0$ when the constant term is allowed: for $\widetilde f =\sum\limits_{n\geq 0} a_n \hbar^n\in\C[[\hbar]]_1$,
\[\widehat f:=\mathcal{B} \widetilde f = a_0\delta_0+\sum\limits_{n\geq0}\frac{a_n}{n!}\xi^n \in \C\delta_0 \oplus \C\{\xi\}.\]
\end{defn}

\begin{defn}\label{defn:Laplace}
Let $\theta\in\R$. Suppose that $\widehat f=\mathcal B\widetilde f$ is integrable at $0$ and it  admits
analytic continuation to a sectorial neighborhood of the ray
$e^{i\theta}\R_{\ge0}$ and has at most exponential growth there. The
Laplace transform of $\widehat f$ in direction $\theta$ is defined by
\begin{equation*}
\mathcal L^\theta\widehat f(\hbar)
:=
\int_{0}^{e^{i\theta}\infty}
e^{-\xi/\hbar}\widehat f(\xi)\,\dd\xi,
\end{equation*}
whenever the integral converges. We call
$\mathcal L^\theta\widehat f$ the Borel--Laplace sum of
$\widetilde f$ in the direction $\theta$. Here we use the convention that $\mathcal{L}^{\theta} a_0 \delta_0 := a_0$.
We denote the composition of Laplace transform and Borel transform by
\begin{equation*}
\mathcal S_\theta := \mathcal{L}^\theta \mathcal{B}.
\end{equation*}
\end{defn}




Let $\widetilde f(\hbar)=\sum\limits_{m\ge0}a_m\hbar^m\in\C[[\hbar]]_1$, and set
\begin{equation*}
\widetilde g(\hbar)=\hbar^\alpha\widetilde f(\hbar),
\end{equation*}
where either $\operatorname{Re}\alpha>0$, or $\alpha=0$ with the above
$\delta_0$ convention. It is well-known that if the regular part of the
Borel transform $\widehat g$ admits analytic continuation to a sectorial
neighborhood of the ray $e^{i\theta}\R_{\ge0}$, with at most
exponential growth there, then 
\begin{equation*}
\hbar^{-\alpha}\mathcal S_\theta\widetilde g(\hbar) \sim \sum_{m\ge0}a_m\hbar^m  \quad \quad \text{as } \ \hbar\rightarrow 0 \ \text{ in }\  \text{S}^{\theta-\frac{\pi}{2},\theta+\frac{\pi}2,r}
\end{equation*}
as a $1$--Gevrey asymptotic expansion in the sector for sufficiently small $r>0$.





\subsection{Borel summability of heat kernel asymptotics}\label{subsec:proof-borel-summability} We study the short-time asymptotic expansion of the heat kernel and its Borel summability. Throughout this subsection, the Gevrey parameter $\hbar$ used in previous subsection will be identified with the heat-time variable $\tau$\footnote{We will use $\tau$ as the time variable in the heat kernel expansion, while $t\in[0,\tau]$ denotes the parameter of path in Section \ref{sec:H^2-flow}.} (see Section \ref{sec:intro-MF}). 

Let $(M,g)$ be a smooth connected $n$--dimensional Riemannian manifold. We use the non--negative Laplace--Beltrami
operator $\Delta\ge 0$, acting on functions,
\[
\Delta f=-\mathrm{div}(\nabla f).
\]
When $(M,g)$ is complete with bounded geometry assumptions, there exists a unique heat
kernel $K(\tau;x,y)$ which is smooth in $(\tau,x,y)\in(0,\infty)\times
M\times M$. The fundamental local result is
\[
K(\tau;x,y)\sim(4\pi \tau)^{-n/2}\exp\Big(-\tfrac{d(x,y)^2}{4\tau}\Big)
\sum_{k\ge 0} u_k(x,y)\tau^k,\quad \tau\rightarrow0^+.
\]

\begin{rmk} The precise meaning of  this asymptotic expansion in the sense of Definition \ref{defn:series} is understood as
\begin{equation*}
(4\pi \tau)^{n/2}
\exp\!\left(\frac{d(x,y)^2}{4\tau}\right)
K(\tau;x,y)
\sim
\sum_{k\ge0}u_k(x,y)\tau^k \in \C[[\tau]].
\end{equation*}
\end{rmk}



On a local chart $\Omega$ with coordinate $(x^i)_{1\le i\le n}$, we can write the Laplace--Beltrami operator as
\[\Delta f=-\sum_{i,j}\partial_i(g^{ij}\partial_jf)+\sum_{i}a^i(x)\partial_if(x)+b(x)f(x),\]
where $a^i(x),b(x)$ are explicitly expressed in terms of the Riemannian metric $g$. If the metric $g$ is real analytic, then $a^i(x),b(x)$ are also real analytic. In this case, we have the following result.


\begin{thm}
\label{thm:local-borel-summability}
Let $(M,g)$ be real analytic. Let $V$ be a sufficiently small compact
convex neighborhood with the heat kernel asymptotic expansion 
\begin{equation}\label{eq:heatGevrey}
K(\tau;x,y)\sim (4\pi\tau)^{-n/2}e^{-d(x,y)^2/(4\tau)}\sum_{j\ge0}u_j(x,y)\tau^j, \quad (x,y)\in V\times V. 
\end{equation}
Then $\sum\limits_{j\ge0}u_j(x,y)\tau^{j}$
is $1$--Gevrey power series uniformly for $(x,y)\in V'\times V'$ with
$V'\Subset V$. % Hence its Borel transform is in $\C\delta_0 \oplus \C\{\xi\}$.  
\end{thm}

\begin{rmk}
The Gevrey estimate of Theorem \ref{thm:local-borel-summability} is a local parametrix statement. In some concrete analytic heat equations, it
is supplied by Borel summability theorems such as those of Lutz--Miyake--Sch\"afke, Costin--Park--Takei and Harg\'e \cite{CPT,Har,Har2,LMS}.
\end{rmk}

The rest of this subsection is devoted to the proof of Theorem
\ref{thm:local-borel-summability}. It is well known that the coefficients in
the short-time heat kernel expansion agree with the coefficients obtained
from the formal solution of the heat equation after extracting the exponential
factor and imposing the standard normalization $u_0(x,x)=1$. Therefore, our
goal is to prove that the coefficients $u_j(x,y)$ obtained recursively from
the transport equations satisfy the $1$--Gevrey estimate
\begin{equation*}
|u_j(x,y)|\leq AB^j j!
\end{equation*}
uniformly for $(x,y)\in V'\times V'$, where $V'\Subset V$ and $A,B>0$ are
constants independent of $j$. 

\vskip 1em

\noindent \textbf{Parametrix, recursion relation and complex extension}
\vskip 1em

First, as in the smooth case, we have the short-time asymptotic expansion of $K(\tau;x,y)$ 
\[K(\tau;x,y)\sim (4\pi\tau)^{-n/2}e^{-d(x,y)^2/(4\tau)}\widetilde{u}(\tau;x,y),\qquad \widetilde{u}(\tau;x,y)=\sum_{i=0}^{\infty}u_i(x,y)\tau^i.\]
For simplicity, denote $H(\tau;x,y)=(4\pi \tau)^{-n/2}e^{-d(x,y)^2/(4\tau)}$. Then one can compute
\begin{equation}\label{eq:heateq}
(\partial_\tau+\Delta_x)(H\widetilde{u})=H\left(\partial_\tau \widetilde{u}+\frac{T}{\tau}\widetilde{u}+\Delta_x \widetilde{u}\right)
\end{equation}
where $T$ is a first-order differential operator of the form
\[
T\widetilde{u}=
\sum_{i,j}\frac12 g^{ij}(x)\partial_i d^2(x,y)\partial_j\widetilde{u}
+B(x,y)\widetilde{u},
\]
with an analytic scalar coefficient $B(x,y)=-\Delta_x\left(\frac{d(x,y)^2}{4}\right)-\frac n2$ on $V\times V$ determined by $g^{ij}$, $a^i(x)$, $b(x)$ and $d(x,y)$. In particular, the coefficients of the first differential operator $T$ are real analytic.

Expanding Equation \eqref{eq:heateq}, we get the recursion relation for $u_j(x,y)$:
\[
Tu_0(x,y)=0,\qquad (T+j+1)u_{j+1}(x,y)=-\Delta u_j(x,y).   
\]
In particular, the normalization $u_0(x,x)=1$ is forced by the delta initial condition of the heat kernel. These coefficient equations admit analytic solutions on $V\times V$.

\smallskip

The analytic hypothesis allows us to extend all coefficients holomorphically to a small complex neighborhood. That is, after shrinking $V$ once and for all, we may choose $\rho>0$ so that every $u_j$ constructed by the recursion extends holomorphically to
\[
\mathcal P_{\rho}(V'):=\bigcup_{(x,y)\in V'\times V'}P_{\rho}(x)\times P_{\rho}(y)\subset\C^n\times\C^n.
\]
Here $P_\rho(x)$ is the complex coordinate polydisc of radius $\rho$ centered at $x$, so that $\mathcal P_\rho(V')$ is a small complex polydisc neighborhood of $V'\times V'$.

\begin{lem}\label{lem:recursionestimate}
After shrinking $V$ if necessary, and fixing $V'\Subset V$, there exist radii $0<\rho_1<\rho_2$ and a constant $C_T>0$ such that the following holds: every coefficient $u_j$ extends holomorphically to $\mathcal P_{\rho_2}(V')$, and we have the uniform estimate
\begin{equation}\label{eq:Tbound-polydisc}
\|u_{j+1}\|_{\rho_1}
\le \frac{C_T}{j+1}\,\|\Delta_x u_j\|_{\rho_2},
\qquad
\|f\|_{\rho}:=\sup_{\mathcal P_{\rho}(V')}|f|.
\end{equation}
In particular, restricting to the real locus gives
\[
\sup_{V'\times V'}|u_{j+1}|
\le \frac{C_T}{j+1}\sup_{V\times V}|\Delta u_j|.
\]
\end{lem}

\begin{proof}
We work in a fixed analytic coordinate chart and still write $\Delta_x$ for the holomorphic extension of the Laplace--Beltrami operator.  

After shrinking $V$, we may assume that any two points of $V$ are joined by a unique real analytic geodesic depending holomorphically on the endpoints. 
Then in the chosen analytic coordinates,
the real analytic dependence of the geodesic on its endpoints extends to a
holomorphic map on a small complex
neighborhood of $V'\times V'$, satisfying the holomorphic geodesic equation (see Equation \eqref{eq:holomorphic-geodesic} for an explicit explanation). Denote such a holomorphic extension by
\[
\Gamma_\sigma(x,y),
\qquad
0\le \sigma\le 1.
\]
So $\Gamma_0(x,y)=x$ and $\Gamma_1(x,y)=y$. 

After possibly shrinking $V$ further and choosing $0<\rho_1<\rho_2$
small enough, we may assume that
\[
(\Gamma_\sigma(x,y),y)\in \mathcal P_{\rho_2}(V'),\qquad\text{for all }(x,y)\in \mathcal P_{\rho_1}(V'),~0\le \sigma\le 1.
\]
In the holomorphic parametrix construction, the coefficients $u_j$
are still determined by the transport equations 
\[
u_{j+1}(x,y)=\pm u_0(x,y)
\int_0^1\sigma^j\,
u_0(\Gamma_\sigma(x,y),y)^{-1}
\bigl(\Delta_x u_j\bigr)(\Gamma_\sigma(x,y),y)
\,\dd\sigma .
\]
The sign is irrelevant for the estimate later. The holomorphicity of each $u_j$ follows from this formula:  
\begin{itemize}
%\setlength{\itemsep}{0pt}
    \item The function $u_0$ is
holomorphic and non-vanishing on a sufficiently small complex
neighborhood of the diagonal.
    \item By induction, suppose that $u_j$
extends holomorphically to the relevant complex tube.  Then
$\Delta_xu_j$ is holomorphic there as well, because $\Delta_x$ has holomorphic coefficients in our analytic coordinates.  
   \item The function $(x,y)\mapsto \Gamma_\sigma(x,y)$ is holomorphic and depends
continuously on $\sigma\in[0,1]$.
\end{itemize}
Then the integrand above is holomorphic
in $(x,y)$ and continuous in $\sigma$.  Therefore the integral
defines a holomorphic function of $(x,y)$.  This proves recursively
that every $u_j$ extends holomorphically to
$\mathcal P_{\rho_2}(V')$, after shrinking the tube if necessary.

\vskip 0.1cm
It remains to prove the estimate.  Since $u_0$ is holomorphic and
non-vanishing on the chosen tube, the ratio $\frac{u_0(x,y)}{u_0(\Gamma_\sigma(x,y),y)}$
is uniformly bounded for $(x,y)\in \mathcal P_{\rho_1}(V'),\
0\le \sigma\le 1$. Set
\[
C_T
:=
\sup_{\substack{(x,y)\in \mathcal P_{\rho_1}(V')\\0\le\sigma\le1}}
\left|
\frac{u_0(x,y)}{u_0(\Gamma_\sigma(x,y),y)}
\right| .
\]
Using the transport formula and the inclusion $(\Gamma_\sigma(x,y),y)\in \mathcal P_{\rho_2}(V')$, 
we obtain
\[
\begin{aligned}
|u_{j+1}(x,y)|\le C_T\int_0^1
\sigma^j\left|(\Delta_x u_j)(\Gamma_\sigma(x,y),y)\right|
\,\dd\sigma\le C_T\|\Delta_x u_j\|_{\rho_2}
\int_0^1 \sigma^j\,\dd\sigma
=\frac{C_T}{j+1}\|\Delta_x u_j\|_{\rho_2}.
\end{aligned}
\]
Taking the supremum over
$(x,y)\in \mathcal P_{\rho_1}(V')$ gives 
$\|u_{j+1}\|_{\rho_1}
\le
\frac{C_T}{j+1}
\|\Delta u_j\|_{\rho_2}$.
Restricting to the real locus immediately gives
\[
\sup_{V'\times V'} |u_{j+1}|
\le\frac{C_T}{j+1}\sup_{V\times V}|\Delta_x u_j|,
\]
where the right-hand side may be taken over $V\times V$, because the
real geodesic segments joining points of $V'$ remain in $V$ after the initial shrinking of $V$. 
\end{proof}

\vskip 1em
\noindent \textbf{Cauchy estimates and the estimates for $u_j$ and the remainder}
\vskip 1em

%The Cauchy estimate will play an important role in the proof.
\begin{lem}\label{lem:Cauchy}
Let $f$ be holomorphic on $P_\rho(z_0)\subset \C^n$ and suppose that $\sup_{P_\rho(z_0)}|f|\le M$.
Then for any multiindex $\alpha\in \mathbb{N}^n$,
\[
|\partial^\alpha f(z_0)|\le M~\frac{\alpha!}{\rho^{|\alpha|}}.
\]
More generally, for any $0<\rho'<\rho$, we have
\[\sup_{z\in P_{\rho'}(z_0)}|\partial^\alpha f(z)|\le M~\frac{\alpha!}{(\rho-\rho')^{|\alpha|}}\]
\end{lem}

\begin{proof}
First, let us prove the one-variable case. Assume $z_0=0$. Cauchy formula says, 
\[f^{(k)}(z)=\frac{k!}{2\pi i}\oint_{|\zeta-z|=r}\frac{f(\zeta)}{(\zeta-z)^{k+1}}\,\dd\zeta.\]
When $|z|\leq\rho'$, choose $r=\rho-\rho'$, then we have
\[|f^{(k)}(z)|\leq \frac{k!}{2\pi}\cdot(2\pi r)\cdot\frac{M}{r^{k+1}}=\frac{k!M}{r^k}=\frac{k!M}{(\rho-\rho')^k}.\]

For the multi-variable in the polydisc, just apply the one-variable Cauchy integral formula in each complex coordinate and iterate.
\end{proof}


The key lemma of the estimate for $u_j(x,y)$ is as follows:

\begin{lem}\label{lem:factorialU}
There exist constants $C_0,A_0>0$ (depending on $V'\subset V$) such that the coefficients satisfy
\begin{equation*}\label{eq:ujBound}
\|u_j\|_{\infty,V'\times V'} \le C_0~A_0^{j}~j!,
\quad j\ge0.
\end{equation*}
\end{lem}

\begin{proof}
First, for $0<\rho''<\rho'$, by Lemma \ref{lem:recursionestimate}, we have
\begin{equation}\label{eq:Tbound2}
\sup_{P_{\rho''}(x)\times P_{\rho''}(y)} |u_{j+1}|
 \le C_T\sup_{P_{\rho'}(x)\times P_{\rho'}(y)} \frac{|\Delta u_j|}{j+1}.
\end{equation}

\smallskip
Since $u_j$ is holomorphic on $P_\rho(x)\times P_\rho(y)$, by Cauchy estimate Lemma \ref{lem:Cauchy}, there exists a constant $C_L$ (depending on uniform bounds of $g^{ij},a^i,b$ on the polydisc) such that
\begin{equation}\label{eq:Lbound}
\sup_{P_{\rho'}(x)\times P_{\rho'}(y)} |\Delta_x u_j|
\le C_L(\rho-\rho')^{-2}
\sup_{P_{\rho}(x)\times P_{\rho}(y)} |u_j|.
\end{equation}

Choose $\rho\geq2\rho'-\rho''$, and combine inequalities \eqref{eq:Tbound2}--\eqref{eq:Lbound}, we get
\begin{align*}
\sup_{P_{\rho''}(x)\times P_{\rho''}(y)} |u_{j+1}|
 \le~&\frac{C_TC_L}{j+1}(\rho-\rho')^{-2}\sup_{P_{\rho}(x)\times P_{\rho}(y)} |u_j|\\
 \le~&\frac{4C_TC_L}{j+1}(\rho-\rho'')^{-2}\sup_{P_{\rho}(x)\times P_{\rho}(y)} |u_j|.
\end{align*}
\medskip
Now choose $N$ large enough, and choose a decreasing sequence of radii
\[\rho_0>\rho_1>\rho_2>\cdots>\rho_N\]
as follows
\[\rho_k=\rho_N+\frac{N-k}{N}\delta, \qquad \delta=\rho_0-\rho_N.\]
Then
\[\|u_N\|_{\rho_N}\leq\left(\frac{C}{\delta^2}\right)^N\frac{N^{2N}}{N!}\|u_0\|_{\rho_0}.\]
Using the weak version of Stirling formula
$N!\ge\left(\frac{N}{e}\right)^N$, we obtain
\[\|u_N\|_{\rho_N}\leq\left(\frac{Ce^2}{\delta^2}\right)^NN!\|u_0\|_{\rho_0}.\]
Then choosing $C_0=\|u_0\|_{\rho_0}$, $A_0=\frac{Ce^2}{\delta^2}$, we are done. 
\end{proof}




It remains to estimate the remainder.  Write
\[
K_N^{\rm loc}(\tau;x,y):=(4\pi\tau)^{-n/2}e^{-d(x,y)^2/(4\tau)}\sum_{j=0}^N u_j(x,y)\tau^j
=H(\tau;x,y)\sum_{j=0}^N u_j(x,y)\tau^j
\]
for the truncated local parametrix.  Then
\[
(\partial_\tau+\Delta_x)K_N^{\rm loc}(\tau;x,y)
=H(\tau;x,y)\tau^N E_N(x,y),
\]
where $E_N$ is an analytic function depending linearly on $\Delta_x u_N$ and lower-order terms.  Using Lemma \ref{lem:factorialU} and Cauchy
bounds for derivatives, we get
\[
\|E_N\|_{\infty,V_1\times V_1}\leq C_2 A_2^{N+1}(N+1)!,
\]
where $V'\Subset V_1\Subset V$.

We now localize the parametrix.  Choose a cutoff $\chi\in C_c^\infty(V\times V)$ which is identically one in a neighborhood of the diagonal over
$V_1\times V_1$, and in particular equals one on $V'\times V'$
after shrinking $V'$ if necessary.  Set
\[
\widehat K_N(\tau;x,y):=\chi(x,y)K_N^{\rm loc}(\tau;x,y).
\]
Then $\widehat K_N=K_N^{\rm loc}$ on $V'\times V'$.  Moreover,
\[
(\partial_\tau+\Delta)\widehat K_N=
\chi H\tau^N E_N+C_N,\qquad C_N:=[\Delta,\chi]K_N^{\rm loc}.
\]
Here $C_N$ is supported where derivatives of $\chi$ are nonzero;
hence it is supported away from the diagonal.

Let
\[
R_N(\tau;x,y):=K(\tau;x,y)-\widehat K_N(\tau;x,y).
\]
Since $\chi=1$ near the diagonal, $R_N$ has zero initial data.  Hence
\[
(\partial_\tau+\Delta)R_N=-\chi H\tau^N E_N
-C_N,\qquad\lim_{\tau\to0^+}R_N(\tau)=0 .
\]
By Duhamel's formula,
\[
\begin{aligned}
R_N(\tau;x,y)
=&-\int_0^\tau\int_M
K(\tau-t;x,z)\chi(z,y)H(t;z,y)t^N E_N(z,y)\,\dd z\,\dd t\\[0.1cm]
&-\int_0^\tau\int_MK(\tau-t;x,z)C_N(t;z,y)\,\dd z\,\dd t .
\end{aligned}
\]
The second integral is exponentially small as $\tau\to0^+$: %Indeed, $C_N$ is supported away from the diagonal, and therefore the Gaussian factor in $K_N^{\rm loc}$ gives an additional exponential decay.  More precisely, for $x,y\in V'$,
\[
\int_0^\tau\int_MK(\tau-t;x,z)C_N(t;z,y)\,\dd z\,\dd t
=O(e^{-c/\tau}),\quad x,y\in V'
\]
for some $c>0$. 

For the first integral, the cutoff restricts the $z$-integration to a
compact subset of $V$, we have
\[
\int_MK(\tau-t;x,z)\chi(z,y)H(t;z,y)\,\dd z
\leq C H(\tau;x,y),\qquad 0<t<\tau .
\]
Therefore, for $x,y\in V'$,
\begin{align*}
|R_N(\tau;x,y)|&\leq C A_2^{N+1}(N+1)!H(\tau;x,y)
\int_0^\tau t^N\,\dd t +O(e^{-c/\tau})\\
&\leq C' A'^{N+1}(N+1)!\tau^{N+1-n/2}e^{-d(x,y)^2/(4\tau)}.
\end{align*}
Since $\widehat K_N=K_N^{\rm loc}$ on $V'\times V'$, this proves the
stated Gevrey remainder estimate in Theorem \ref{thm:local-borel-summability}.\qed




\subsection{Complexification}\label{sec: complexification}

Theorem \ref{thm:local-borel-summability} suggests that we could perform resurgence analysis of heat kernel on the plane of Borel transform of time. As explained in the introduction, this phenomenon is expected to reflect a version of Picard–Lefschetz theory on the complexified path space. A natural way to introduce such complexification is to extend the real-analytic manifold $(M,g)$ into $(M_\C,g_{\C})$ where $M_{\C}$ is a complex manifold, $g_{\C}$ is a holomorphic metric such that $M$ is a connected component of the real locus of $M_{\C}$ and $g$ is the restriction of $g_{\C}$. 

%In this subsection, we will review several constructions of complexification of real-analytic manifolds and describe the induced complex structure on the path space. 

%\subsubsection{Complexification of real-analytic manifolds}

There are several closely related notions of complexification in the literature. Classical results of Bruhat--Whitney and Grauert imply that every paracompact real-analytic manifold admits a local/germ-level complexification, and one may often choose $M_{\C}$ Stein. There are two classical constructions for such complexifications, known as Grauert tubes. 
\begin{itemize}
%\setlength\itemsep{0pt}
\item[(1)] The first one, by Bruhat and Whitney \cite{BW}, is based on a gluing construction: complex analytic extensions of real analytic charts are glued together in a suitable way. 
\item[(2)] The second one was devised by Grauert \cite{Gr}, and developed by Akhiezzer and Gindikin \cite{AG}, Guillemin and Stenzel \cite{GS}, Lempert and Sz\"oke \cite{LS,Sz}, Patrizio and Wong \cite{PW}.
\end{itemize}



When $M$ is equipped with a real-analytic Riemannian metric $g$, and in local real-analytic coordinates $x=(x^1,\dots,x^n)$, write
\[
g=\sum_{i,j} g_{ij}(x)\,\dd x^i\dd x^j.
\]
Since the coefficients $g_{ij}$ are real analytic, they admit holomorphic extensions $g_{ij}^{\C}(z)$ to a small complex neighborhood. Hence one obtains a holomorphic symmetric $2$-tensor
\[
g_{\C}:=\sum_{i,j} g_{ij}^{\C}(z)\,\dd z^i\dd z^j ,
\]
which restricts to $g$ on the real locus. After shrinking the neighborhood if necessary, $\det(g_{ij}^{\C}(z))\neq 0$, so $g_{\C}$ is a holomorphic metric. This extension is canonical at the level of germs: different analytic charts or different local models give biholomorphically equivalent holomorphic metrics along $M$.

As in the real case, there is a unique torsion-free holomorphic connection $\nabla^{\C}$, also known as holomorphic Levi-Civita connection \cite{BiswasDumitrescuHRM},  satisfying
\[
\nabla^{\C} g_{\C}=0 .
\]
In local coordinates, its Christoffel symbols are the holomorphic extensions of the usual one:
\[
\Gamma^{k,\C}_{ij}
=
\frac12\sum_{\ell}(g_{\C})^{k\ell}
\bigl(
\partial_i g^{\C}_{j\ell}
+
\partial_j g^{\C}_{i\ell}
-
\partial_\ell g^{\C}_{ij}
\bigr).
\]
The geodesic equation also extends holomorphically:
\begin{equation}\label{eq:holomorphic-geodesic}
\ddot z^k+\sum_{i,j}\Gamma^{k,\C}_{ij}(z)\dot z^i\dot z^j=0.    
\end{equation}


Their solutions will be called holomorphic geodesics. 

%This is the basic reason why real geodesics, the world function, and the heat-kernel phase all admit local holomorphic continuations to the complexified target.





\begin{rmk} 
For Lie groups, one also has
the universal complexification, characterized by a universal property for
holomorphic homomorphisms into complex Lie groups; see, for example,
\cite{Hochschild1965,Hochschild1966,HilgertNeeb2012,Neeb2000}.  For Riemannian manifolds with negative sectional curvature, the structure of a Grauert domain is obstructed to extend to the whole tangent bundle. There is an upper bound on the radius of a Grauert tube. See Lempert-Sz\"oke \cite{LS,Sz}. 
\end{rmk}


%\sout{In our paper, we use the following elementary embedded version of complexification:
%locally along the given real manifold or real Lie group, the complexification
%looks like the standard pair $(\C^n,\R^n)$, and no global maximality or uniqueness condition is imposed unless explicitly stated.}

In this paper, we use a fixed admissible embedded complexification
\[
\iota:M\hookrightarrow M_{\mathbb C}.
\]
This means that $M$ is a connected component of the real locus of $M_{\C}$, locally modeled on
$(\mathbb C^n,\mathbb R^n)$, and that the metric $g$ extends to a
holomorphic metric $g_{\mathbb C}$ on $M_{\mathbb C}$.  


\begin{assumption}\label{assumption}
We assume this complexification is sufficiently large for the analysis below: the
holomorphic Levi-Civita connection and the holomorphic geodesic flow are
defined on a domain containing all real endpoints and complexified initial
directions considered in the paper.  In this
relative sense, the complexification is holomorphically geodesically
complete for the family of geodesics under study.  We do not impose any
global maximality or uniqueness condition on $M_{\mathbb C}$ unless this
is explicitly stated.    
\end{assumption}



%\subsubsection{Examples}

%\subsubsection*{Example 1. compact symmetric spaces of rank 1}
\begin{eg}[Compact symmetric spaces of rank 1]
The compact symmetric spaces of rank 1 are classified as follows:
\begin{itemize}
%\setlength{\itemsep}{-1pt}
    \item[$(1)$] the $n$-sphere $S^n$, $n\geq2$;
    \item[$(2)$] the real projective space $\mathbb{RP}^n$, $n\geq2$;
    \item[$(3)$] the complex projective space $\mathbb{CP}^n$, $n\geq1$;
    \item[$(4)$] the quaternionic projective space $\mathbb{HP}^n$, $n\geq1$;
    \item[$(5)$] the Cayley projective plane $\mathbb{P}^2(\mathrm{Cayley})$.
\end{itemize}
Denote them by $M_i$ respectively according to the order. Their (good) complexifications\footnote{A good complexification is a smooth affine algebraic variety $U$ defined over $\R$ such that $U(\R)\cong M$ and $U(\R)\hookrightarrow U(\C)$ is a homotopy equivalence. Complexification $M_{\C}$ of $M$ with the condition that
\[\iota:M\hookrightarrow M_{\C} \quad\text{is a homotopy equivalence},\]
is also known as a minimal complexification.} are well known as follows (Morimoto-Nagano \cite{MN}, Huckleberry-Snow \cite{HS}, Azad \cite{Az}, Patrizio-Wong \cite{PW}).
\begin{itemize}
    \item[(I)] The standard models $M_{i,\C}$ of their complex affine algebraic manifolds are: 
\begin{itemize}
%\setlength{\itemsep}{-1pt}
    \item[$(1)$] $M_{1,\C}=Q^n=$ complex affine hyperquadric, $n\geq2$;
    \item[$(2)$] $M_{2,\C}=\mathbb{P}^n(\C)-\bar{Q}^{n-1}=$ compact quadric, $n\geq2$;
    \item[$(3)$] $M_{3,\C}=\mathbb{P}^n(\C)\times\mathbb{P}^n(\C)-\mathbb{P}_{\infty}^{N-1}(\C)$, $n\geq1$, $N=(n+1)^2-1$;
    \item[$(4)$] $M_{4,\C}=\mathrm{Gr}(2,2n;\C)-\mathbb{P}_{\infty}^{N-1}(\C)$, $n\geq1$, $N=n(2n-1)-1$;
    \item[$(5)$] $M_{5,\C}$ is a 16 dimensional Stein manifold.
\end{itemize}
Here $\mathbb{P}_{\infty}^{N-1}(\C)$ is a hyperplane at infinity in $\mathbb{P}^N(\C)$.
   \item[(II)] There exists a real analytic strictly plurisubharmonic exhaustion $\rho_i:M_i^{\C}\rightarrow[1,\infty)$ such that 
\[M_i=\{\rho_i=1\},\qquad\dim_{\R}M_i=\dim_{\C}M_{i,\C}.\] Furthermore, the plurisubharmonic function $u_i=\cosh^{-1}\rho_i$ satisfies the homogeneous Monge-Amp\`ere equation on $M_{i,\C}-M_i$.  
\end{itemize}

\end{eg}

\begin{eg}[Kulkarni \cite{Ku}]
Let $G$ be a compact Lie group and $H$ be a closed subgroup of $G$. Then the homogeneous space $G/H$ admits a minimal affine algebraic complexification $G^{\C}/H^{\C}$. Furthermore, 
the hodge numbers of $G^{\C}/H^{\C}$ have the property
\[h^{p,q}(G^{\C}/H^{\C})=0,\quad p\neq q.\]
\end{eg}


%\subsubsection*{Example 2. complexification of the real hyperbolic plane $\HH_{\R}^2$\label{HR2}}

\begin{eg}[Complexification of the real hyperbolic plane $\HH_{\R}^2$] This example will play a major role in this paper, and we will discuss it in some detail. We consider the following three models.
\begin{itemize}
%\setlength\itemsep{0pt}
\item[$(1)$] The upper half plane model: 
\[\HH_{\R}^2=\{z=x+iy\in\C\mid y>0\}\] 
with the metric
\[g=\frac{\dd x^2+\dd y^2}{y^2}=-\frac{4\dd z\,\dd\bar{z}}{(z-\bar{z})^2}.\]
\item[$(2)$] The symmetric space model: \[\HH_{\R}^2=SL(2,\R)/SO(2,\R)\]
with the ($SL(2,\R)$-invariant) metric.
\item[$(3)$] The hyperboloid model: 
\[\HH_{\R}^2=\{-x_0^2+x_1^2+x_2^2=-1\mid x_i\in\R, x_0>0\}\] 
with the metric
\[\dd s^2=-\dd x_0^2+\dd x_1^2+\dd x_2^2\big|_{\HH_{\R}^2}.\]
\end{itemize}

There are two well-known complexifications of $\HH_{\R}^2$.
\vskip 0.1cm
\noindent$\bullet~$ \textbf{Minimal complexification}

\vskip 0.1cm
For the upper half plane model, we consider the space 
\[X:=\HH_{\R}^2\times\HH_{\R}^{2-}=\{(z,w)\in\C^2\mid \im z>0,\im w<0\}\]
with the conjugation
\[\tau:X\rightarrow X, \quad (z,w)\mapsto(\bar{w},\bar{z})\]
and the holomorphic metric
\[g_{\C}=-2(\dd z\otimes \dd w+\dd w\otimes \dd z)/(z-w)^2.\]
We have the embedding 
\[\iota:\HH_{\R}^2\hookrightarrow X,\quad z\mapsto(z,\bar{z}),\]
satisfying $\iota^*g_{\C}=g$. Such $X$ is a $SL(2,\R)$-invariant Stein domain in $\HH_{\C}^2$ (see the definition below).  $SL(2,\R)$ acts on $X$ properly, and the metric $g_{\C}$ is a (the unique) $SL(2,\C)$-invariant holomorphic metric. 


\vskip 0.2cm
\noindent$\bullet~$ \textbf{Affine complexification}
\vskip 0.1cm

Upon complexifying $SL(2,\R)$, $SO(2,\R)$, we obtain the affine complexification
\[\HH_{\C}^2=SL(2,\C)/SO(2,\C)\cong\{-z_0^2+z_1^2+z_2^2=-1\mid z_i\in\C\}\]
with the holomorphic metric
\[\dd s_{\C}^2=-\dd z_0^2+\dd z_1^2+\dd z_2^2\big|_{\HH_{\C}^2}.\]
The map
\[SL(2,\R)/SO(2,\R)\hookrightarrow SL(2,\C)/SO(2,\C),\qquad g\cdot SO(2,\R)\mapsto g\cdot SO(2,\C)\]
constitutes a $SL(2,\R)$-equivariant embedding which realizes $\HH_{\R}^2=SL(2,\R)/SO(2,\R)$ as one connected component of the totally real submanifold of $\HH_{\C}^2$. %Also note that $\HH_{\R}^2$ is a connected component of fixed point set of the complex conjugation.

The two complexifications play rather different roles. The domain
\[
X=\HH_{\R}^2\times \HH_{\R}^{2-}
\]
is the minimal complexification: it is a Stein neighborhood of the real slice and corresponds to choosing one single-valued branch of the complexified geodesic geometry. By contrast, the affine quadric $\HH_{\C}^2$ is larger and contains the extra complex branches responsible for monodromy and the infinitely many holomorphic geodesics with the same real endpoints (see Example \ref{eg:complex-geodesics-H2}). In particular, $\HH^2_{\C}$ satisfies Assumption \ref{assumption} while $X$ does not. We have
\begin{itemize}
    \item $\HH_{\R}^2$ is a connected component of fixed point set of the complex conjugation in $\HH_{\C}^2$.
    \item $\iota:\HH_{\R}^2\hookrightarrow\HH_{\C}^2$ is not homotopy equivalence.
\end{itemize}


More concretely, consider the map
\[
\Psi:X\longrightarrow \HH_{\C}^2,\qquad
\Psi(z,w)=\left(
i\frac{1+zw}{z-w},\,
i\frac{z+w}{z-w},\,
i\frac{1-zw}{z-w}
\right)
\]
One checks that
\[
\Psi^*(\dd s_{\C}^2)= -\,\frac{2(\dd z\otimes \dd w+\dd w\otimes \dd z)}{(z-w)^2}=g_{\C}.
\]
Hence the holomorphic metric on the minimal model is simply the pullback of the affine holomorphic metric; in particular the holomorphic Levi-Civita connection on $X$ is the pullback of that of $\dd s_{\C}^2$.    

\end{eg}



\subsection{Picard–Lefschetz problem on complexified path spaces}\label{sec:PL-problem-path}
We are ready to discuss the Picard–Lefschetz problem on the complexified path space and connect it to the resurgence analysis of heat kernel. 

%\subsubsection{Path space on complexified target}

Let $M_{\C}$ be a complexified manifold of $M$ with an anti-holomorphic involution $\tau$, complex structure $J$, and
let $a,b\in M\subset M_{\C}$ be real endpoints. For variational problems, it is natural to work on the Sobolev path space
\[
\mathcal P_{a,b}(M_{\C})
=
\left\{
\gamma\in H^1([0,\tau],M_{\C}) :
\gamma(0)=a,\ \gamma(\tau)=b
\right\}.
\]
Since $H^1([0,\tau])\hookrightarrow C^0([0,\tau])$, $\gamma$ is continuous and the endpoint conditions are
well-defined. The complex structure on $M_{\C}$ induces a pointwise almost complex
structure on the path space. For 
$$
\xi\in T_\gamma\mathcal P_{a,b}(M_{\C})
=
H^1_0([0,\tau],\gamma^*TM_{\C}),
$$
the induced complex structure is defined by 
\[
(\mathbb J_\gamma \xi)(t)
:=
J_{\gamma(t)}\xi(t).
\]
This satisfies $\mathbb J_\gamma^2=-\mathrm{id}$. 

Let $g_{\C}$ be a holomorphic complex-bilinear metric on
$M_{\C}$.  We define the complexified energy by
\[
\mathcal E_{\C}(\gamma)
:=
\frac{1}{4}\int_0^\tau
g_{\C,\gamma(t)}(\dot\gamma(t),\dot\gamma(t))\,\dd t,\quad \gamma\in H^1.
\]
We choose $\frac{1}{4}$ for convenience of resurgence analysis. In local holomorphic coordinates $z=(z^1,\ldots,z^n)$,
$$
\mathcal E_{\C}(z)
=
\frac{1}{4}\int_0^\tau \sum_{i,j}
g_{ij}(z(t))\,\dot z^i(t)\dot z^j(t)\,\dd t,
$$
where the functions $g_{ij}$ are holomorphic.  
$\mathcal E_{\C}$ is holomorphic in the Banach-manifold sense:
in $H^1$-charts it is a holomorphic map from an open subset of a complex
Hilbert space to $\C$.
Equivalently, its first variation is complex-linear in the variation.
Indeed, for
$
\xi\in T_\gamma\mathcal P^1_{a,b}(M_{\C})
$, one has
\[
\dd\mathcal E_{\C,\gamma}(\xi)
=
\frac{1}{4}\int_0^\tau\sum_{i,j,k}
(\partial_k g_{ij})(\gamma(t))\,\xi^k(t)\dot\gamma^i(t)\dot\gamma^j(t)\,\dd t
+\frac{1}{2}
\int_0^\tau \sum_{i,j}
g_{ij}(\gamma(t))\,\dot\xi^i(t)\dot\gamma^j(t)\,\dd t.
\]
Moreover,
\[
\dd\mathcal E_{\C,\gamma}(\mathbb J_\gamma\xi)
=i\, \dd \mathcal E_{\C,\gamma}(\xi),
\]
since the coefficients $g_{ij}$ are holomorphic and the metric is
complex-bilinear. Thus we have arrived at the infinite-dimensional Picard–Lefschetz problem associated to the action
$$\mathcal E_{\C}: \PP_{a,b}(M_{\C})\to \C. 
$$
\begin{rmk}
As explained in the introduction, we can rescale the time by $\tau$ to normalize it to the unit interval $[0,1]$. Then $\tau$ plays the role of $\hbar$, and it is nature to use $\tau \mathcal E_{\C}$ for the holomorphic Morse function in order to match our convention in the finite dimensional case. Nevertheless, they lead to the same critical paths and Lefschetz thimbles. In the following discussions, we will use $\mathcal E_{\C}$ to derive the Morse flow equation and use the critical values of $\tau \mathcal E_{\C}$ to identify the Borel singularities. 
\end{rmk}

%\subsubsection{Critical path and holomorphic geodesic}
Let
$$
\gamma: [0,\tau]\to M_{\C}, \qquad \gamma \in \PP_{a,b}(M_{\C})
$$
be a critical path of the energy functional $\mathcal E_{\C}$. By the standard calculus of variation, $\gamma$ satisfies 
\[
\ddot z^k+\sum_{i,j}\Gamma^{k,\C}_{ij}(z)\dot z^i\dot z^j=0.
\]
where $z^i(t)$ represents $\gamma$ in local holomorphic coordinates. This is the complex analogue of geodesic equation, and solutions are holomorphic geodesics. Thus critical paths of $\mathcal E_{\C}$ corresponds precisely to holomorphic geodesics from $a$ to $b$. 

%\subsubsection{Morse flow} 

Let us now consider the downward gradient Morse flow associated to $\mathcal E_{\C}$. First we choose a hermitian metric $h$ on $M_\C$. It induces a $L^2$ hermitian inner product $<\cdot,\cdot>_\hbar$ on $\mathcal P_{a,b}(M_{\C})$ by 
$$
<\delta \gamma_1, \delta \gamma_2>_h:=\int_0^\tau h(\delta \gamma_1, \delta \gamma_2) dt, \quad  \text{for}\quad \delta \gamma_i\in T_\gamma\mathcal P_{a,b}(M_{\C}).
$$
Write
$
h=\sum\limits_{i,j}h_{i\bar j}\dd z^i \dd\bar z^{\bar j}
$ in a local holomorphic coordinates $z=(z^1,\ldots,z^n)$.  Then 
$$
<\delta \gamma_1, \delta \gamma_2>_h=\int_0^\tau \sum_{i,j} h_{i\bar j}(\gamma(t))\delta \gamma_1^i(t) \overline{\delta \gamma^{j}_2(t)}\,\dd t. 
$$

Let $\theta$ be a chosen phase. Consider the the downward gradient flow of 
$$
\re(e^{-i\theta}\mathcal E_\C). 
$$
It can be described by a map $\gamma(t, s): [0,\tau] \times [0,+\infty) \to M_{\C}$ satisfying the flow equation 
\[\partial_s\gamma=-\nabla_{L^2}\re(e^{-i\theta}\mathcal E_\C)(\gamma),\]
with the fixed endpoint conditions $\gamma(0,s)=a,\,\gamma(\tau,s)=b$. Here $\nabla_{L^2}$ is the gradient with respect to the above $L^2$-hermitian metric. 


In a local holomorphic coordinate $\{z^i\}$ of $M_{\C}$, we have 
$
g_{\C}=\sum\limits_{i,j}g_{ij}\dd z^i \dd z^j,  \ h=\sum\limits_{i,j}h_{i\bar j}\dd z^i \dd\bar z^j. 
$
We also write the flow locally as $z^i(t,s)$. Then the above downward gradient Morse flow is explicitly 
\begin{equation*}
\frac{\partial z^j}{\partial s}=\frac{1}{2}\sum_{k,m}e^{i\theta}h^{j\bar k}\overline{g_{mk}\nabla^{\C}_{\dot z}\dot z^m}.
\end{equation*}
Here $\dot z^m=\frac{\partial z^m}{\partial t}$ and $\nabla^{\C}$ is the holomorphic Levi-Civita connection of the holomorphic metric $g_\C$. Both $h^{j\bar k}$ and $g_{mk}$ are evaluated at $\gamma(t,s)$.


\begin{prop}\label{prop:FG}
Let $z(s,\cdot)$ be a smooth solution of the flow equation. Then
\begin{equation}\label{eq:Fdecrease}
\frac{\dd}{\dd s}F_{\theta}\bigl(z(s,\cdot)\bigr)
=-\norm{\partial_s z(s,\cdot)}_{L^2_t}^2\le 0,
\qquad
\frac{\dd}{\dd s}G_{\theta}\bigl(z(s,\cdot)\bigr)=0.
\end{equation}
Here $\norm{- }_{L^2_t}^2=<-,->_{h}$. In particular, $G_{\theta}$ is constant along every finite-energy half-trajectory.
\end{prop}

\begin{proof}
The first identity follows from the basic property of downward gradient flow. Since $\mathcal E_{\C}$ is holomorphic on $\PP_d^{\C}$, the Cauchy--Riemann equations imply that $\nabla G_{\theta(d)}=\mathbb J\nabla F_{\theta(d)}$. Therefore,
\begin{align*}
\frac{\dd}{\dd s}G_{\theta(d)}(z(s))
=&~\dd\,G_{\theta(d)}(z(s))(\partial_s z)
=\re<\nabla G_{\theta(d)},\partial_s z>_h\\
=&-\re<\mathbb J\nabla F_{\theta(d)},\nabla F_{\theta(d)}>_h=0.\qedhere
\end{align*}
\end{proof}


\begin{eg}
Take $M=\R^n$ with Euclidean metric $g$, and $M_{\C}=\C^n$ with holomorphic coordinate $z=(z^1,\ldots,z^n)$. The holomorphic metric and the standard hermitian metric are 
$g_{\C}=\sum\limits_{i=1}^n\dd z^i \dd z^i,\ h=\sum\limits_{i=1}^n\dd z^i \dd \bar{z}^i$. The holomorphic energy function is 
\[\mathcal E_{\C}(z)=\frac{1}{4}\int_0^{\tau}\sum_{k=1}^n\left(\dot{z}^k\right)^2\dd t.\]
In the flat case, the critical path (geodesic) is unique and given by the affine segment
$z(t)=a+\frac{t}{\tau}(b-a)$. The downward gradient Morse flow of $ \re(e^{-i\theta}\mathcal E_\C)$ is 
\[\partial_s z(s,t)=\frac{1}{2}e^{i\theta}\partial_t^2 \bar{z}(s,t),\qquad z(s,0)=a,~ z(s,\tau)=b.\]
\end{eg}

\smallskip
\begin{eg}\label{eg:H^2-PL} Consider the hyperboloid model: 
\[\HH_{\R}^2=\{-x_0^2+x_1^2+x_2^2=-1\mid x_i\in\R, x_0>0\}\]
with the complexification 
\[\HH_{\C}^2=\{-z_0^2+z_1^2+z_2^2=-1\mid z_i\in\C\}\]
and the holomorphic metric
\[\dd s_{\C}^2=-\dd z_0^2+\dd z_1^2+\dd z_2^2\big|_{\HH_{\C}^2}.\]

Let us denote\footnote{We regard $z,w\in\mathbb C^3$ as column vectors in matrix formulas.  Later,
when writing points or paths explicitly, we often use the horizontal notation
$z=(z_0,z_1,z_2)$ only for readability.  This should not cause confusion:
all products involving $\eta$, projections, and tangent equations are computed with column vectors.}
\[\langle z,w\rangle_{\eta}=z^T\eta w, \qquad z,w\in \C^3, \quad \eta:=\diag(-1,1,1)
\]
for the complex bilinear extension of the Lorentz form. Then
\[\HH_{\C}^2=\{z\in\C^3: \langle z, z\rangle_\eta=-1\}.\]
The holomorphic Levi-Civita connection on the quadric $\HH_{\C}^2$ is obtained by tangential projection from that on  $\C^3$. If $\gamma(t)$ is a holomorphic curve in $\HH_{\C}^2$, then
\[
\nabla^{\C}_{\dot\gamma}\dot\gamma=0
\quad\Longleftrightarrow\quad
\ddot\gamma \ \text{is normal to } T_{\gamma}\HH_{\C}^2 .
\]
Since
\[
T_{\gamma}\HH_{\C}^2=\{\xi\in \C^3:\ \langle\gamma,\xi\rangle_{\eta}=0\},
\]
the normal line is $\C\gamma$, and then the geodesic equation is
\[
\ddot\gamma=\lambda(t)\gamma .
\]
Differentiating $\langle\gamma,\gamma\rangle_{\eta}=-1$ twice gives
\[
\langle\gamma,\dot\gamma\rangle_{\eta}=0,
\qquad
\langle\gamma,\ddot\gamma\rangle_{\eta}=-\langle\dot\gamma,\dot\gamma\rangle_{\eta},
\]
hence $\lambda(t)=\langle\dot\gamma,\dot\gamma\rangle_{\eta}$. Along a geodesic this quantity is constant:
\[
\frac{d}{dt}\langle \dot\gamma,\dot\gamma\rangle_\eta
=2\langle \ddot\gamma,\dot\gamma\rangle_\eta
=2\lambda \langle \gamma,\dot\gamma\rangle_\eta=0.
\]
Denote this constant by $c$, then every holomorphic geodesic satisfies the linear ODE
\[
\ddot\gamma=c\,\gamma .
\]
Therefore, after choosing a branch $\mu$ with $\mu^2=c$, the holomorphic geodesics are
\[
\gamma(t)=\cosh(\mu t)\,a+\frac{\sinh(\mu t)}{\mu}\,v,
\qquad
v=\dot\gamma(0)\in T_a\HH_{\C}^2,\quad \langle v,v\rangle_{\eta}=\mu^2 .
\]

Now let $a,b\in \HH_{\R}^2$ and let $d=d(a,b)>0$ be their real hyperbolic distance, so that
\[
\langle a,b\rangle_{\eta}=-\cosh d.
\]
Define the unit tangent vector at $a$ pointing toward $b$ by
\[
u=\frac{b-\cosh d\, a}{\sinh d},
\qquad
\langle a,u\rangle_{\eta}=0,\quad \langle u,u\rangle_{\eta}=1.
\]
Then every holomorphic geodesic in the affine complexification joining $a$ to $b$ is contained in the complex $2$-plane $\C a\oplus \C u$ and is given by
\begin{align*}
\gamma_k(t)
=~&\cosh\!\Big(\frac{d+2\pi i k}{\tau}t\Big)\,a
+\sinh\!\Big(\frac{d+2\pi i k}{\tau}t\Big)\,u,
\qquad k\in\Z \\
=~&\cosh\!\Big(\frac{d+2\pi i k}{\tau}t\Big)\,a
+\frac{\sinh\!\big(\frac{d+2\pi i k}{\tau}t\big)}{\sinh d}\,
\bigl(b-\cosh d\,a\bigr).
\end{align*}
The corresponding energy is
\[
\mathcal E_{\C}(\gamma_{k})
=\frac{1}{4}\int_0^{\tau}g_{\HH_{\C}^2}(\gamma_k'(t),\gamma_k'(t))\,\dd t
=\frac{(d+2\pi i k)^2}{4\tau}.
\]

Let us now describe the Morse flow. Let $h$ denote the restriction of the standard ambient Hermitian metric on $\C^3$. Thus for any $\xi,\zeta\in T_{z}\HH_{\C}^2\,\subset\,\C^3$, their hermitian inner product is
\[
h(\xi,\zeta)=\xi\cdot \overline{\zeta}.
\]
The induced $L^2$-hermitian metric on the path space is
\[<\xi,\zeta>_h:=\int_0^{\tau} \xi(t)\cdot \overline{\zeta(t)}\,\dd t.
\]

Let $P_z$ be the $h$-orthogonal projection onto $T_z\HH_{\C}^2$, given by
\begin{equation*}\label{eq:projection}
P_z(v)=v-\frac{z^T\eta v}{|z|^2}\,\overline{\eta z},\qquad v\in\C^3.
\end{equation*}
Then the downward gradient Morse flow of $ \re(e^{-i\theta}\mathcal E_\C)$ becomes 
\[
\partial_s z=P_z\bigl(e^{i\theta}\eta\,\overline{z_{tt}}\bigr),
\qquad
z(s,0)=a,
\qquad
z(s,\tau)=b,
\]
with the constraint
\[\langle z(s,t), z(s,t)\rangle_{\eta}=-1.\]
\end{eg}


\subsection{Holomorphic geodesics and Borel singularities}
Now we consider the resurgence analysis of heat kernel. For $x,y$ in a sufficiently small compact convex neighborhood $V\subset M$, the expansion in Theorem \ref{thm:local-borel-summability} is the usual short-time expansion on the real
manifold. It is attached to the real minimal branch
\[
\Phi_0(x,y)=\frac{d(x,y)^2}{4},
\]
coming from the unique minimal real geodesic in a sufficiently small convex
neighborhood.



More generally, in the complexification $M_{\mathbb C}$, one may also consider asymptotic solutions of the heat equation attached to other holomorphic geodesic branches.

Throughout this subsection, unless explicitly stated otherwise, the variables $x$ and $y$ are allowed to be complex points of
$M_{\mathbb C}$. We write $\Delta_x^{\mathbb C}$ for the holomorphic extension of the Laplace operator in the $x$-variable. When $x,y\in M$, the notation $\Delta_x$ refers to the original real Laplacian.

Let
\[
E:TM_{\mathbb C}\longrightarrow M_{\mathbb C}\times M_{\mathbb C},
\qquad E(x,v)=\bigl(x,\exp_x^{\mathbb C}(v)\bigr),
\]
be the holomorphic endpoint map associated with the holomorphic
Levi-Civita connection of $g_{\mathbb C}$. For later use, we introduce the following definitions.






\begin{defn}
Let $x_0\in M_{\mathbb C}$ and $v_0\in T_{x_0}M_{\mathbb C}$.  Set
\[
y_0:=\exp_{x_0}^{\mathbb C}(v_0),\qquad
\gamma_0(t):=\exp_{x_0}^{\mathbb C}(t v_0),\qquad 0\leq t\leq 1.
\]

We say that $v_0$, or equivalently the geodesic segment $\gamma_0$, is non-conjugate at the endpoint if
\[
\dd_v\exp_{x_0}^{\mathbb C}\big|_{v=v_0}:
T_{x_0}M_{\mathbb C}\longrightarrow T_{y_0}M_{\mathbb C}
\]
is an isomorphism.  We say that $\gamma_0$ is non-conjugate (or non-caustic) along the segment, if
\[
\dd_v\exp_{x_0}^{\mathbb C}\big|_{v=t v_0}:
T_{x_0}M_{\mathbb C}\longrightarrow T_{\gamma_0(t)}M_{\mathbb C}
\]
is an isomorphism for every $0<t\leq 1$.
\end{defn}

By the holomorphic inverse function theorem, if
$(x_0,v_0)$ is no-conjugate, then there exists a connected open set $U_\gamma\subset TM_{\mathbb C}$ such that
\[
E_\gamma:=E|_{U_\gamma}:U_\gamma\longrightarrow \Omega_\gamma
\]
is biholomorphic onto its image
\[
\Omega_\gamma:=E_\gamma(U_\gamma)\subset M_{\mathbb C}\times M_{\mathbb C}.
\]
{If, moreover, the central segment $\gamma_0$ is non-conjugate along the
whole interval $0<t\leq 1$, then we shall shrink $U_\gamma$, if
necessary, so that for every $(x,v')\in U_\gamma$,
\[
\dd_v\exp_x^{\mathbb C}\big|_{v=t v'}
\]
is an isomorphism for all $0<t\leq 1$.  In this case we call $\Omega_\gamma$ a no-conjugate (or non-caustic) geodesic branch chart.}

\begin{defn}
With the above notations. We denote the inverse branch in the second variable by
\[
v_\gamma:\Omega_\gamma\longrightarrow TM_{\mathbb C},
\qquad v_\gamma(x,y)\in T_xM_{\mathbb C},
\]
so that
\[
y=\exp_x^{\mathbb C}\bigl(v_\gamma(x,y)\bigr).
\]
The associated holomorphic geodesic branch is defined to be
\[
\gamma_{x,y}(t):=\exp_x^{\mathbb C}\bigl(t\,v_\gamma(x,y)\bigr),\qquad
0\leq t\leq 1.
\]
Thus
\[
\gamma_{x,y}(0)=x,\qquad\gamma_{x,y}(1)=y,\qquad
\dot\gamma_{x,y}(0)=v_\gamma(x,y).
\]
\end{defn}


Furthermore, we define its squared complex length to be
\[
\ell_\gamma(x,y)^2:=\int_0^1g_{\mathbb C,\gamma_{x,y}(t)}
\bigl(\dot\gamma_{x,y}(t),\dot\gamma_{x,y}(t)\bigr)\,\dd t.
\]
By construction, the integrand is constant
in $t$:
\[
\frac{\dd}{\dd t}g_{\mathbb C,\gamma_{x,y}(t)}\bigl(\dot\gamma_{x,y}(t),\dot\gamma_{x,y}(t)\bigr)
=2g_{\mathbb C,\gamma_{x,y}(t)}\bigl(\nabla^{\mathbb C}_{\dot\gamma_{x,y}}\dot\gamma_{x,y},\dot\gamma_{x,y}\bigr)=0.
\]
Hence
\[
\ell_\gamma(x,y)^2=g_{\C,x}\bigl(v_\gamma(x,y),v_\gamma(x,y)\bigr).
\]


\begin{lem}\label{lem:eikonal-equation}
Let $\Omega_{\gamma}$, $v_{\gamma}$ be defined as above and let $\Phi_\gamma=\frac{\ell_{\gamma}(x,y)^2}{4}$, then it satisfies the holomorphic eikonal equation
\[
\Phi_\gamma=g_{\mathbb C}\bigl(\nabla_x^{\mathbb C}\Phi_\gamma,\nabla_x^{\mathbb C}\Phi_\gamma\bigr)
\qquad\text{on }~ \Omega_\gamma.
\]
Here $\nabla_x^{\mathbb C}\Phi_\gamma$ denotes the holomorphic gradient
with respect to $g_{\mathbb C}$ in the
$x$-variable. %More precisely, for each $(x,y)\in\Omega_\gamma$,
%\[\nabla_x^{\mathbb C}\Psi_\gamma(x,y)\in T_xM_{\mathbb C}\]
%is uniquely determined by
%\[g_{\mathbb C,x}\bigl(\nabla_x^{\mathbb C}\Psi_\gamma(x,y),\xi\bigr)=d_x\Psi_\gamma(x,y)(\xi),
%\qquad\xi\in T_xM_{\mathbb C}.
%\]
%Equivalently,
%\[g_{\mathbb C}\bigl(\nabla_x^{\mathbb C}\Psi_\gamma,
%        \nabla_x^{\mathbb C}\Psi_\gamma\bigr)=
%  g_{\mathbb C}^{-1}
%  \bigl(d_x\Psi_\gamma,d_x\Psi_\gamma\bigr).
%\]
If $x,y\in M$ and
the branch is the real geodesic branch, then
$\nabla_x^{\mathbb C}$ restricts to the usual real gradient with respect to $g=g_{\mathbb C}|_M$.
\end{lem}



\begin{proof}
Fix $y$, and vary the initial endpoint $x$.  Let
$x_\varepsilon$ be a holomorphic curve with
\[
x_0=x, \qquad \left.\frac{\dd x_\varepsilon}{\dd\varepsilon}\right|_{\varepsilon=0}=\xi \in T_xM_{\C}.
\]
After shrinking the $\varepsilon$-disc, we may assume
\[
(x_\varepsilon,y)\in \Omega_\gamma,
\] 
then the chosen geodesic branch $v_{\gamma_{x_{\varepsilon},y}}$ depends
analytically on $\varepsilon$.  Hence there is an analytic $\varepsilon$--family of
geodesics
\[
\gamma_\varepsilon:[0,1]\to M_{\C},\qquad
\gamma_\varepsilon(0)=x_\varepsilon,\quad\gamma_\varepsilon(1)=y .
\]
Set
\[
J(t):=\left.\frac{\partial \gamma_\varepsilon(t)}
{\partial \varepsilon}\right|_{\varepsilon=0}.
\]
Then
\[
J(0)=\xi, \qquad J(1)=0.
\]
Differentiating the energy
\[
\ell_{\gamma_\varepsilon}(x_\varepsilon,y)^2=\int_0^1
g_{\C}(\dot\gamma_\varepsilon(t),
\dot\gamma_\varepsilon(t))\,\dd t
\]
at $\varepsilon=0$, we obtain
\begin{align*}
  \dd_x(\ell_\gamma^2)(\xi)
&=2\int_0^1g_{\C}(\nabla_t^{\C} J,\dot\gamma)\,\dd t  \\
&=2\Big[g_{\C}(J,\dot\gamma)
\Big]_{t=0}^{t=1}-2\int_0^1g_{\C}(J,\nabla_t^{\C}\dot\gamma)\,\dd t =-2\,g_{\C}(\xi,\dot\gamma(0)).
\end{align*}

It follows that
\[
\dd_x\Phi_\gamma(\xi)
=-\frac12 g_{\C}(\xi,\dot\gamma(0)),\quad\text{
equivalently,}\quad
\nabla_x^{\C}\Phi_\gamma=-\frac12 \dot\gamma(0).
\]
Thus,
\[
g_{\mathbb C}\bigl(\nabla_x^{\mathbb C}\Phi_\gamma,
\nabla_x^{\mathbb C}\Phi_\gamma\bigr)
=\frac14 g_{\mathbb C,x}\bigl(v_\gamma(x,y),v_\gamma(x,y)\bigr)
=\frac{\ell_\gamma(x,y)^2}{4}=\Phi_\gamma(x,y).\qedhere
\]
\end{proof}



\begin{thm}\label{thm:heat-kernel-expansion-complex-geodesic}
With the above notations and assumptions, {on a no-conjugate geodesic branch chart} $\Omega_{\gamma}$, there exists a unique $\C$-valued formal solution $\widetilde{K}_\gamma(\tau;x,y)$ to the heat equation 
\[(\partial_{\tau}+\Delta_x^{\C})\widetilde{K}_{\gamma}=0\]
of the form
\begin{equation}\label{eq:formal-expansion-gamma}
\widetilde{K}_\gamma(\tau;x,y)=(4\pi \tau)^{-n/2}
e^{-\Phi_\gamma(x,y)/\tau}\sum_{j\geq 0}u_{\gamma,j}(x,y)\tau^j.    
\end{equation}
normalized by $u_{\gamma,0}=J_\gamma^{-1/2}$ where $J_\gamma$ is the Jacobian of the complex exponential
map along the chosen holomorphic geodesic branch. 
Here $\Delta_x^{\mathbb C}$ denotes
the holomorphic extension of the non-negative Laplacian in the $x$-variable, and $\sum\limits_{j\geq 0}u_{\gamma,j}(x,y)\tau^j$ is a formal series in $\tau$. 
\end{thm}

\begin{proof}
We look for a formal solution of the form
\begin{equation*}
\widetilde K_\gamma(\tau;x,y)
=
(4\pi\tau)^{-n/2}
e^{-\Phi_\gamma(x,y)/\tau}
\widetilde u_\gamma(\tau;x,y),
\qquad
\widetilde u_\gamma(\tau;x,y)
=
\sum_{j\ge0}u_{\gamma,j}(x,y)\tau^j .
\end{equation*}
Substituting this expression into the complexified heat operator gives
\begin{align*}
(\partial_\tau+\Delta_x^{\C})\widetilde K_\gamma
=
(4\pi\tau)^{-n/2}
e^{-\Phi_\gamma/\tau}
\Bigg[
\frac{\Phi_\gamma-g_{\mathbb C}\bigl(\nabla_x^{\mathbb C}\Phi_\gamma,\nabla_x^{\mathbb C}\Phi_\gamma\bigr)}{\tau^2}\widetilde u_\gamma  
+&\frac1\tau\left(
2\nabla_x^{\C}\Phi_\gamma\cdot\nabla_x^{\C}\widetilde u_\gamma-(\Delta_x^{\C}\Phi_\gamma)\widetilde u_\gamma
-\frac n2\widetilde u_\gamma\right) \\
+&\partial_\tau\widetilde u_\gamma
+\Delta_x^{\C}\widetilde u_\gamma
\Bigg].
\end{align*}
By Lemma \ref{lem:eikonal-equation}, the coefficient of the leading $\tau^{-2}$-order term vanishes. Define the transport operator
\begin{equation*}
T_\gamma:=2\nabla_x^{\C}\Phi_\gamma\cdot\nabla_x^{\C}-
\Delta_x^{\C}\Phi_\gamma-\frac n2.
\end{equation*}
Then the remaining expression becomes
\begin{equation*}
(\partial_\tau+\Delta_x^{\C})\widetilde K_\gamma
=(4\pi\tau)^{-n/2}
e^{-\Phi_\gamma/\tau}
\left[\frac1\tau T_\gamma\widetilde u_\gamma
+\partial_\tau\widetilde u_\gamma
+\Delta_x^{\C}\widetilde u_\gamma\right].
\end{equation*}

Now substitute
\begin{equation*}
\widetilde u_\gamma(\tau;x,y)
=\sum_{j\ge0}u_{\gamma,j}(x,y)\tau^j .
\end{equation*}
Comparing powers of $\tau$ gives the transport equations
\begin{equation*}
T_\gamma u_{\gamma,0}=0,\qquad 
\left(T_\gamma+j\right)u_{\gamma,j}
=-\Delta_x^{\C}u_{\gamma,j-1},~~\text{for}~~j\ge1.
\end{equation*}

It remains to see that these equations can be solved locally along the chosen holomorphic geodesic branch. Since the branch is assumed to be no-conjugate, every
$y$ near the chosen branch can be written uniquely in the form
\[
y=\exp_x^{\C}(v(x,y)),\qquad v(x,y)\in T_xM_{\C},
\]
with $v(x,y)$ depending analytically on $x$ and $y$. Near the chosen initial vector and away from the isotropic
cone $g_{\C}(v,v)=0$, choose a branch of the square root and set
\[
r:=\bigl(g_{\C}(v,v)\bigr)^{1/2},\qquad
\omega:=\frac{v}{r}.
\]
Then
\[
y=\exp_x^{\C}(r\omega).
\]
Here $r$ is a complex-valued radial variable and $\omega$ lies on the complex unit quadric
\[
Q_y^{\C}:=
\{\omega\in T_yM_{\C}: g_{\C}(\omega,\omega)=1\}.
\]
In these coordinates, we have
\begin{equation*}
\Phi_\gamma=\frac{r^2}{4},\qquad
2\nabla_x^{\C}\Phi_\gamma\cdot\nabla_x^{\C}
=r\partial_r.
\end{equation*}

Let $J_\gamma(r,\omega)$ be the reduced Jacobian of the complex exponential
map, defined by
\begin{equation*}
\dd\operatorname{vol}_{g_\C}
=r^{n-1}J_\gamma(r,\omega)\,\dd r\,\dd\omega.
\end{equation*}
Then, with our convention $\Delta^{\C}=-\operatorname{div}\nabla^{\C}$,
\begin{equation*}
\Delta_x^{\C}\Phi_\gamma
=-\frac n2-\frac r2\partial_r\log J_\gamma.
\end{equation*}
Therefore
\begin{equation*}
T_\gamma=r\partial_r+\frac r2\partial_r\log J_\gamma.
\end{equation*}
The leading transport equation becomes
\begin{equation*}
\left(r\partial_r+\frac r2\partial_r\log J_\gamma
\right)u_{\gamma,0}=0.
\end{equation*}
Hence locally
\begin{equation*}
u_{\gamma,0}(r,\omega)
=C_\gamma(\omega)J_\gamma(r,\omega)^{-1/2},
\end{equation*}
We fix the normalization by choosing the analytic function $C_\gamma(\omega)\equiv1$.

For $j\ge1$, the transport equation becomes
\begin{equation*}
\left(r\partial_r+\frac r2\partial_r\log J_\gamma+j
\right)u_{\gamma,j}=-\Delta_x^{\C} u_{\gamma,j-1}.
\end{equation*}
Since $J_\gamma$ is non-vanishing on the chosen branch, this is a first-order
linear ordinary differential equation in the radial variable with analytic 
coefficients. Once the normalization of $u_{\gamma,j}$ is fixed on a transverse section, it has a unique analytic solution along the branch.

By induction, all coefficients $u_{\gamma,j}$ are locally constructed and
are determined by the transport recursion. Therefore the formal expression
\begin{equation*}
\widetilde K_\gamma(\tau;x,y)
=(4\pi\tau)^{-n/2}e^{-\Phi_\gamma(x,y)/\tau}
\sum_{j\ge0}u_{\gamma,j}(x,y)\tau^j
\end{equation*}
satisfies
\[
(\partial_\tau+\Delta_x^{\C})\widetilde K_\gamma=0.\qedhere
\]
\end{proof}



\begin{rmk}
The natural normalization in Theorem \ref{thm:heat-kernel-expansion-complex-geodesic} agrees with the usual WKB normalization, and we shall call the formal expansion as the WKB-type formal solution or a WKB branch.    
\end{rmk}


To state a converse rigid statement, we introduce the following definition.
\begin{defn}\label{def:geometric-WKB}
Let $\Omega_{\gamma}$ be a no-conjugate geodesic branch chart, assume that there exist holomorphic functions
\[
\Psi_\kappa,\ u_{\kappa,j}\in \mathcal O(\Omega_{\gamma}),
\qquad j\geq 0,
\]
with $u_{\kappa,0}$ nowhere vanishing, such that
\[
\widetilde K_\kappa(\tau;x,y)=\frac{1}{(4\pi\tau)^{n/2}}
\exp\left(-\frac{\Psi_\kappa(x,y)}{\tau}\right)
\sum_{j\geq 0}u_{\kappa,j}(x,y)\tau^j
\]
solves
\[
\left(\partial_\tau+\Delta_x^{\mathbb C}\right)\widetilde K_\kappa=0
\]
as a formal series in $\tau$. We say that it is a geometric branch if the complex Lagrangian generated by
$\Psi_\omega$,
\[
\Lambda_{\Psi_\kappa}:=\left\{\bigl(x,\dd_x\Psi_\kappa(x,y);y,-\dd_y\Psi_\kappa(x,y)\bigr)
:(x,y)\in\Omega_{\gamma}\right\},
\]
is contained in the holomorphic geodesic relation. That is, on the no-conjugate chart $\Omega_\gamma$, there exists a holomorphic inverse branch $v_\gamma(x,y)$ of the endpoint map $E(x,v)=(x,\exp_x^{\mathbb C}v)$ such that
\[
y=\exp_x^{\mathbb C}v_\gamma(x,y),\quad\text{and}\quad
\nabla_x^{\mathbb C}\Psi_\kappa(x,y)=-\frac12 v_\gamma(x,y).
\]
\end{defn}


\begin{prop}\label{prop:phase-rigidity-geodesic-branch}
Let $\widetilde{K}_{\kappa}$ be a geometric branch in Definition \ref{def:geometric-WKB}. Then
\[
\Psi_\kappa(x,y)=\frac{\ell_\gamma(x,y)^2}{4}\qquad
\text{on }~ \Omega_\gamma.
\]
And we can normalize it as Theorem \ref{thm:heat-kernel-expansion-complex-geodesic}.
\end{prop}

\begin{proof}
Substituting
\[
\widetilde K_\kappa=(4\pi\tau)^{-n/2}e^{-\Psi_\kappa/\tau}
\sum_{j\geq 0}u_{\kappa,j}\tau^j
\]
into $(\partial_\tau+\Delta_x^{\mathbb C})\widetilde K_\kappa=0$ and
taking the coefficient of $\tau^{-2}$ gives
\[
\left(\Psi_\kappa-g_{\mathbb C}(\nabla_x^{\C}\Psi_\kappa,\nabla_x^{\C}\Psi_\kappa)\right)u_{\kappa,0}=0.
\]
Since $u_{\kappa,0}$ is nowhere vanishing, $\Psi_\kappa$ satisfies the holomorphic eikonal equation
\[
\Psi_\kappa=g_{\mathbb C}(\nabla_x^{\C}\Psi_\kappa,\nabla_x^{\C}\Psi_\kappa).
\]
Combined with the geometric condition, in particular,
\[\nabla_x^{\mathbb C}\Psi_\kappa(x,y)=-\frac12 v_\gamma(x,y),\]
we get
\[
\Psi_{\kappa}=g_{\mathbb C}(\nabla_x^{\C}\Psi_\kappa,\nabla_x^{\C}\Psi_\kappa)=\frac14 g_{\mathbb C,x}(v_\gamma(x,y),v_\gamma(x,y))
=\frac{1}{4}\ell_{\gamma}(x,y)^2.\qedhere\]
\end{proof}





Now we fix $x,y\in M\subset M_{\C}$ and fix a minimal real geodesic branch $\gamma_0$ from $x$ to $y$, with phase
\[
\Phi_0(x,y)=\frac{\ell_0(x,y)^2}{4}.
\]
We have proven in Theorem \ref{thm:local-borel-summability} that the heat kernel has a $1$--Gevrey formal expansion
\[
\widetilde K_0(\tau;x,y)
=\frac{1}{(4\pi\tau)^{n/2}}
e^{-\Phi_0(x,y)/\tau}\widetilde{u}_0(\tau;x,y),\qquad \widetilde{u}_0(\tau;x,y)=
\sum_{k\geq 0}u_{0,k}(x,y)\tau^k.
\]
Let
\[
\widehat u_0(\zeta;x,y)
=\sum_{k\geq 0}u_{0,k}(x,y)\frac{\zeta^k}{k!}
\]
be the Borel transform of $\widetilde{u}_0$. 




\begin{thm}\label{thm:Borel-singularity-complex-geodesic}
With the above notations, suppose that $\widehat u_0$ admits analytic continuation in the Borel
plane, and suppose that $\zeta=\omega(x,y)$ is a non-zero {integrable} singularity.
After taking a small Hankel contour
$\Gamma_\omega$ around $\zeta=\omega(x,y)$, assume that the expression
\[
\widetilde K(\tau;x,y)
=
\frac{1}{(4\pi \tau)^{n/2}}
e^{-\Phi_0(x,y)/\tau}
\frac{1}{\tau}
\int_{\Gamma_\omega}
e^{-\zeta/\tau}
\widehat u_0(\zeta;x,y)\,\dd\zeta
\]
has a non-zero geometric branch of the complexified heat equation attached to a holomorphic geodesic $\gamma$
\[
\widetilde K(\tau;x,y)
\sim
\frac{1}{(4\pi\tau)^{n/2}}
e^{-\Psi_\kappa(x,y)/\tau}
\sum_{k\geq0}u_{\kappa,j}(x,y)\tau^j,
\]
%which formally solves the complexified heat equation and satisfies the initial condition attached to $\gamma$. Then
Then
\[
\omega(x,y)=\Psi_\kappa(x,y)-\Phi_0(x,y)=\frac{1}{4}\bigl(l_{\gamma}^2(x,y)-l_0^2(x,y)\bigr).
\]
\end{thm}



\begin{proof}
The exponential factor in the Hankel contribution is
\[
e^{-\Phi_0(x,y)/\tau}e^{-\zeta/\tau}.
\]
Since the Hankel contour $\Gamma_\omega$ is taken around the singularity
$\zeta=\omega(x,y)$, write locally $\zeta=\omega(x,y)+\xi$, 
\[
e^{-\Phi_0(x,y)/\tau}e^{-\zeta/\tau}=e^{-(\Phi_0(x,y)+\omega(x,y))/\tau}e^{-\xi/\tau}.
\]
Thus the WKB phase carried by the Hankel contribution is
\[
\Psi_\kappa(x,y)=\Phi_0(x,y)+\omega(x,y).
\]
By Proposition \ref{prop:phase-rigidity-geodesic-branch},
\[
\omega(x,y)=\Psi_\kappa(x,y)-\Phi_0(x,y)=\frac{1}{4}\bigl(l_{\gamma}^2(x,y)-l_0^2(x,y)\bigr).\qedhere
\]
\end{proof}


We next turn to several examples to clarify how the above correspondence between Borel-plane singularities and complexified geodesic branches is realized in concrete situations.
\smallskip
\begin{eg}[The hyperbolic plane $\HH_{\R}^2$]
\label{eg:complex-geodesics-H2}
The exact heat kernel formula on $\HH^2_{\R}$ is known
\begin{equation*}
K_{\HH_{\R}^{2}}(\tau;a,b)
=
\frac{\sqrt{2}e^{-\tau/4}}{(4\pi \tau)^{3/2}}\int_{d^2/4}^{\infty}\frac{e^{-w/\tau}}{\sqrt{\cosh\sqrt{4w}-\cosh d}}~\dd w
\end{equation*}
which has already the form of Laplace-transform. The Borel coordinate is related to $w$ by 
$$
\xi=w-d^2/4. 
$$
The Borel singularities appear at
$$
{\cosh\sqrt{4w}}=\cosh d, \quad \text{i.e.}\quad \xi=\frac{(d+2\pi i k)^2}{4}-\frac{d^2}{4}, \quad k\in \Z.
$$
They correspond precisely to the critical values of $\tau\mathcal E_{\C}$ relative to the real geodesic $\gamma_{{0},d}$
$$
\tau\mathcal E_{\C}(\gamma_{{k},d})-\tau\mathcal E_{\C}(\gamma_{{0},d})
$$
as computed in Example \ref{eg:H^2-PL}. A detailed resurgence and Picard–Lefschetz analysis of this example will be given in Section \ref{sec:H^2-resurgence} and Section \ref{sec:H^2-flow}.


For general even-dimensional hyperbolic spaces, explicit Borel summation / analytic continuation statements were analyzed in \cite{Du}.  A resurgence analysis of Borel singularities of hyperbolic metric on Riemann surfaces, within the Picard–Lefschetz framework,  is also studied in \cite{Wu-thesis}.  

\end{eg}




\begin{eg}[The sphere $S^2$, $0<d(x,y)<\pi$]
The heat kernel of $S^2$ is radial, and has the spectral expansion
\[
K_{S^2}(\tau;r)
=
\frac{1}{4\pi}
\sum_{\ell=0}^{\infty}
(2\ell+1)e^{-\tau\ell(\ell+1)}P_\ell(\cos r),\qquad r:=d_{S^2}(x,y)\in(0,\pi).
\]
To read off the Borel-plane singularities, it is more transparent to use the following Abel--Poisson integral representation:
\[
K_{S^2}(\tau;r)=\frac{\sqrt{2}\,e^{\tau/4}}{(4\pi\tau)^{3/2}}\sum_{k\in\mathbb Z}(-1)^k
\int_r^\pi \frac{(\phi+2\pi k)
\exp\!\left[-\frac{(\phi+2\pi k)^2}{4\tau}\right]}
{\sqrt{\cos r-\cos\phi}}
\,\dd\phi.
\]
Indeed, the exponential weights in the Abel--Poisson formula are of the form
\[
\exp\!\left[-\frac{(\phi+2\pi k)^2}{4\tau}\right].
\]
Near the lower endpoint $\phi=r$, the $k$-th summand has action
\[
S_k(r)=\frac{(r+2\pi k)^2}{4}.
\]
Thus the shortest geodesic sector corresponds to $k=0$, and the other terms correspond to the other geodesic branches on the same great circle, with lengths
\[
\ell_k=|r+2\pi k|,\qquad k\in\mathbb Z.
\]
Therefore the action difference between the $k$-th geodesic branch and the shortest one is
\[
\omega_k(r)=S_k(r)-S_0(r)=
\frac{(r+2\pi k)^2-r^2}{4}
=\pi k r+\pi^2 k^2.
\]
These are precisely the expected locations of the Borel singularities of the perturbative sector associated with the shortest geodesic.

\begin{rmk}
Unlike the affine complexification of $\HH_{\R}^2$, the affine complexification of $S^2$ does not produce additional holomorphic geodesics between real endpoints. The branches indexed by $k$ are simply the real winding geodesics on $S^2$, so this example records the real multi-geodesic contribution rather than a new complexified one.    
\end{rmk}


\begin{rmk}
The restriction $0<r<\pi$ avoids the diagonal $r=0$ and the cut locus $r=\pi$. At $r=0$, closed geodesic degeneracies occur and the explicit computation relating geodesic information to Borel singularities will be computed in Example \ref{eg:r=0}. At $r=\pi$, the shortest geodesic is no longer isolated, since $x$ and $y$ are antipodal. 
\end{rmk}
\end{eg}



\begin{eg}[The sphere $S^2$, $d(x,y)=0$]\label{eg:r=0}
For $r=0$, we can compute explicitly as follows. The heat trace of $S^2$ is known by
\[
Z_{S^2}(\tau)
:=
\operatorname{Tr}(e^{-\tau\Delta_{S^2}})
=
\sum_{\ell=0}^{\infty}(2\ell+1)e^{-\ell(\ell+1)\tau}.
\]
It is convenient to work with the shifted heat trace
\[
\Phi_{S^2}(\tau)
:=e^{-\tau/4}Z_{S^2}(\tau)
=\sum_{\ell=0}^{\infty}
2\left(\ell+\frac12\right)e^{-(\ell+\frac12)^2\tau}.
\]
Set $f_\tau(x):=2x e^{-\tau x^2}$, then
\[
\Phi_{S^2}(\tau)=\sum_{\ell=0}^{\infty}
f_\tau\left(\ell+\frac12\right).
\]
We apply the midpoint Euler--Maclaurin formula
\[
\sum_{\ell=0}^{\infty} f\left(\ell+\frac12\right)
\sim
\int_0^\infty f(x)\,\dd x
-\sum_{m\geq 1}
\frac{B_{2m}(\frac12)}{(2m)!}
f^{(2m-1)}(0),
\]
where $B_{2m}(x)$ denotes the Bernoulli polynomial.  In the present case,
\[
\int_0^\infty f_\tau(x)\,\dd x=\frac1\tau,\qquad 
f_\tau^{(2m-1)}(0)
=2\frac{(-\tau)^{m-1}}{(m-1)!}(2m-1)!.
\]
Substituting this into the Euler--Maclaurin formula gives
\[
\Phi_{S^2}(\tau)
\sim\frac1\tau+\sum_{m\geq 1}
\frac{(-1)^m B_{2m}(\frac12)}{m!}\tau^{m-1}=\frac{1}{\tau}\left(1+\frac{(-1)^m B_{2m}(\frac12)}{m!}
\tau^{m}\right).
\]
Set $\phi_{S^2}(\tau):=\tau\Phi_{S^2}(\tau)$, then
\[
\widehat{\phi}_{S^2}(\xi)
=1+\sum_{m\geq 1}
\frac{(-1)^m B_{2m}(\frac12)}{(m!)^2}\xi^m.
\]

Using the identities of Bernoulli polynomials
\[
B_{2m}\left(\frac12\right)
=\left(2^{1-2m}-1\right)B_{2m},\qquad
B_{2m}
=(-1)^{m+1}
\frac{2(2m)!}{(2\pi)^{2m}}\zeta(2m).
\]
and the formula of the Dirichlet eta function
\[
\eta(s):=(1-2^{1-s})\zeta(s)
=\sum_{k\geq 1}\frac{(-1)^{k-1}}{k^s},
\]
we obtain
\[
\widehat{\phi}_{S^2}(\xi)=1+2\sum_{m\geq 1}
\frac{\binom{2m}{m}}{4^m}
\left(\sum_{k\geq 1}\frac{(-1)^{k-1}}{(k^2\pi^2)^m}\right)\xi^m.
\]
Interchanging the two sums for $|\xi|<\pi^2$,
gives
\[
\widehat{\phi}_{S^2}(\xi)
=1+2\sum_{k\geq 1}(-1)^{k-1}
\sum_{m\geq 1}\frac{\binom{2m}{m}}{4^m}\left(\frac{\xi}{k^2\pi^2}\right)^m.
\]
Using the generating function
\[
\sum_{m\geq 0}\frac{\binom{2m}{m}}{4^m}x^m
=(1-x)^{-1/2},
\]
we arrive at the explicit Borel expression
\[
\widehat{\phi}_{S^2}(\xi)=1+2\sum_{k\geq 1}(-1)^{k-1}\left[\left(1-\frac{\xi}{k^2\pi^2}\right)^{-1/2}
-1\right].
\]
It shows that $\widehat{\phi}_{S^2}(\xi)$ has algebraic branch points at $k^2\pi^2$, $k\in\Z_{\ge1}$. Furthermore, 
\[
\operatorname{Sing}\bigl(\widehat{\phi}_{S^2}\bigr)
=\{\,k^2\pi^2:\ k\in\Z_{\ge1}\,\}.
\]

With the Laplace kernel $e^{-\xi/\tau}$, the
corresponding exponentially small terms have the form
\[
e^{-k^2\pi^2/\tau}=e^{-(2k\pi)^2/(4\tau)}.
\]
where $(2k\pi)^2/(4\tau)$ is the geodesic action associated with the $k$-fold iterate of a
closed geodesic of length $2\pi$ on the unit sphere.  Thus in the Borel plane, we can see the correspondence
\[
\text{closed geodesic length }2\pi k
\quad\Longleftrightarrow\quad
\text{Borel singularity } \xi=k^2\pi^2
\quad\Longleftrightarrow\quad
e^{-k^2\pi^2/\tau}.
\]
\end{eg}


\begin{rmk} For the compact rank-one symmetric spaces, the method relating heat-kernel
exponential sectors to Borel singularities is similar to the above on $S^2$, i.e the basic action-difference formula for the singularity locations. The difference may appear in the amplitudes and in the local type of the Borel singularities, such as the transverse Jacobi field mulitiplicities and the Stokes coefficients.
\end{rmk}




\begin{eg}
The higher-rank analogue of the hyperbolic plane
$SL(2,\R)/SO(2,\R)$ is
\[
M_n=SL(n,\R)/SO(n,\R).
\]
We use the standard real model
\[
M_n\cong \mathcal P_n^{\R}=
\{Q\in \operatorname{Sym}_n(\R): Q>0,\ \det Q=1\},
\qquad gSO(n,\R)\longmapsto gg^T.
\]
Its affine complexification is
\[
SL(n,\C)/SO(n,\C)
\simeq\mathcal P_n^{\C}
:=\{Q\in \operatorname{Sym}_n(\C):\det Q=1\}.
\]
Let
\[
\mathfrak a=
\left\{H=\operatorname{diag}(h_1,\ldots,h_n):\sum_{j=1}^n h_j=0\right\}.
\]
Up to the Weyl group, a point of $M_n$ may be written as
\[
Q=\exp(2H_0),\qquad H_0\in\mathfrak a^+ .
\]
With the normalization
\[
\langle A,B\rangle=\operatorname{tr}(AB),\qquad A,B\in\mathfrak a,
\]
the principal real geodesic from $I$ to $\exp(2H_0)$ is
\[
Q_0(t)=\exp(2tH_0),\qquad 0\leq t\leq 1,
\]
and its squared length is, up to the fixed global normalization of the metric,
\[
L_0^2=\operatorname{tr}(H_0^2).
\]

The affine complexification remembers the multi-valuedness of the matrix
logarithm.  If
\[
m=(m_1,\ldots,m_n)\in\mathbb Z^n,\qquad\sum_{j=1}^n m_j=0,
\]
then
\[
Q_m(t)=\exp(2t H_m),
\qquad H_m:=H_0+\pi i\,\operatorname{diag}(m_1,\ldots,m_n),\quad 0\leq t\leq 1,
\]
is a holomorphic geodesic branch joining the same endpoints in
$\mathcal P_n^{\C}$.  Its complex squared length is
\[
L_m^2=\operatorname{tr}(H_m^2).
\]
Consequently the affine model contains the full lattice of holomorphic geodesic
actions
\[
\left\{\frac14\,\operatorname{tr}(H_m^2):m\in\mathbb Z^n,\ \sum_j m_j=0\right\}.
\]
Relative to the principal real branch, the corresponding action
differences are
\[
\omega_m=\frac14\left(\operatorname{tr}(H_m^2)-\operatorname{tr}(H_0^2)\right).
\]

The same lattice is visible from the heat kernel.  Let $K_\tau(Q)$ be the
heat kernel based at the identity, and write its radial part as
$K_\tau(\exp(2H))$.  For a noncompact symmetric space $G/K$, the Abel
transform sends the radial heat equation to a Euclidean heat equation on
$\mathfrak a$, with the usual $\rho$-shift:
\[
\mathcal A K_\tau(H)
=e^{-|\rho|^2\tau}
(4\pi\tau)^{-r/2}
\exp\left(-\frac{|H|^2}{4\tau}\right),
\qquad r=\dim\mathfrak a=n-1 .
\]
Hence the heat kernel is recovered by applying the inverse Abel transform:
\[
K_\tau(\exp(2H))
=\mathcal A^{-1}
\left[e^{-|\rho|^2\tau}
(4\pi\tau)^{-r/2}
\exp\left(-\frac{|\cdot|^2}{4\tau}\right)
\right](H).
\]
For $SL(n,\R)/SO(n,\R)$, the inverse Abel transform has
Abel-type kernels with square-root branch factors
\[
\bigl(\sinh(\xi_a-h_j)\bigr)^{-1/2},
\]
so its complex singular hyperplanes are
\[
\xi_a-h_j\in \pi i\mathbb Z.
\]
These are exactly the hyperplanes produced by the same logarithmic
multi-valuedness $H_0\mapsto H_m$.
Equivalently, the possible Borel singular locations, relative to
the principal branch, are governed by the same action differences
\[
\omega_m=\frac14\left(\operatorname{tr}(H_m^2)-\operatorname{tr}(H_0^2)\right),\qquad
H_m=H_0+\pi i\,\operatorname{diag}(m_1,\ldots,m_n).
\]
Similarly, this lattice of extra branches is not visible in the minimal or germ-level complexification.  
\end{eg}



\subsection{Heat asymptotics and Alien operator}\label{sec:heat-alien}

We briefly recall the alien-operator convention used below, following
\cite{LiLiTang2026}.  For systematic accounts of alien calculus and resurgence
theory, we refer the reader to Écalle's foundational trilogy
\cite{Ecalle1981I,Ecalle1981II,Ecalle1985III} and to the monograph of
Mitschi and Sauzin \cite{MitschiSauzin2016}. 

Suppose that $\widehat f(\xi)$ is a holomorphic germ near $0$ obtained by the Borel transform of a $1$--Gevrey series. Let $\gamma_\omega^+$ be a continuation path approaching a small
neighborhood of $\omega$ from the positive side of the relevant Stokes ray,
and let $\gamma_\omega^-$ be obtained from $\gamma_\omega^+$ by adjoining a
small clockwise loop around $\omega$. Suppose that $\widehat{\phi}$ admits analytic
continuation along $\gamma_\omega^{\pm}$ to a singular point $\omega\in\C^*$ and denote the resulting holomorphic germ near $\omega$ by $\cont_{\gamma_\omega^{\pm}} f$. The local variation is
\[
\var_{\omega}^+(\widehat f)
:=\cont_{\gamma_\omega^{+}}\widehat f-
\cont_{\gamma_\omega^{-}}\widehat f .
\]
It measures the local monodromy of the analytically continued Borel germ at $\omega$. See Figure \ref{fig:alien-operator-schematic} below.

\begin{figure}[H]
\centering
\begin{tikzpicture}[>=Latex, scale=0.62]

\coordinate (O) at (0,0);
\coordinate (W) at (5.2,3.5);
\coordinate (WW) at (7.8,5.27);
\coordinate (M1) at (2.6,1.75);
\coordinate (M2) at (1.71,1.13);

\fill[green!20] (5.6,3.7) circle (0.38);
\draw[green!50!black] (5.6,3.7) circle (0.38);
\node[green!40!black] at (6,4.5) {$D$};

\filldraw[black] (O) circle (1.3pt);
\node[below left] at (O) {$0$};

\draw[dashed, thick] (O) -- (WW);
\draw[->,thick]
    (0.38,0.00)
    .. controls (1.2,0.05) and (3.1,0.85) ..
    (4,1.95)
    .. controls (4.62,2.55) and (4.70,5.00) .. (5.4,3.9)
    ;
\draw[->,thick]
    (0.38,0.00)
    .. controls (1.2,0.05) and (3.2,0.75) ..
    (5.35,2.7)
    .. controls (5.62,3) and (5.8,3.2) .. (5.6,3.7)
    ;
\draw[red, line width=1.1pt]
    ($(W)+(-0.12,-0.12)$) -- ($(W)+(0.12,0.12)$);
\draw[red, line width=1.1pt]
    ($(W)+(-0.12,0.12)$) -- ($(W)+(0.12,-0.12)$);
\draw[red, line width=1.1pt]
    ($(M1)+(-0.12,-0.12)$) -- ($(M1)+(0.12,0.12)$);
\draw[red, line width=1.1pt]
    ($(M1)+(-0.12,0.12)$) -- ($(M1)+(0.12,-0.12)$);
\draw[red, line width=1.1pt]
    ($(M2)+(-0.12,-0.12)$) -- ($(M2)+(0.12,0.12)$);
\draw[red, line width=1.1pt]
    ($(M2)+(-0.12,0.12)$) -- ($(M2)+(0.12,-0.12)$);
\node[below,red] at (W) {$w$};
\node[above] at (6.05,1.58) {$\gamma_w^+$};
\node[above] at (4.65,4.58) {$\gamma_w^-$};
\end{tikzpicture}
\caption{The path $\gamma_\omega^+$ analytically continues the Borel germ
from the point near origin to a small disc $D$ behind the singular point $\omega$,
passing to the right of the preceding singularities on the same ray.  The
path $\gamma_\omega^-$ is obtained from $\gamma_\omega^+$ by adjoining a
small clockwise loop around $\omega$ inside $D$.}
\label{fig:alien-operator-schematic}
\end{figure}

\begin{defn}
\label{def:positive-alien-operator}
With the above notation and assumption, we define
\[
\Delta_{\omega}^{+}
=
\mathcal B^{-1}\circ
\tau_{-\omega}\circ
\operatorname{var}_{\omega}^{+}\circ
\mathcal B : \ \hbar^\alpha\C[[\hbar]]
\longrightarrow
\hbar^\beta\C[[\hbar]]
\]
if $\tau_{-\omega}\operatorname{var}_\omega^+(\widehat f)
\in
\xi^{\beta-1}\C\{\xi\}$ where $(\tau_{-\omega}\widehat\varphi)(\xi)=\widehat\varphi(\xi+\omega)$. Here $\beta\in\C$ satisfies that $\re\beta>0$.
\end{defn}

The corresponding pointed alien operator is
\begin{equation}\label{eq:pointed-positive-alien-operator}
\dot\Delta_{\omega}^{+}
:=e^{-\omega/\hbar}\Delta_{\omega}^{+}.
\end{equation}

For formal objects such as short-time heat-kernel WKB expansions, we use the
following convention.  The alien calculation is performed on the 
$1$--Gevrey power series, while the exponential factor is treated as the action
label and the common algebraic heat-kernel prefactor is carried through the
calculation unchanged: \footnote{In \'Ecalle's general theory
of resurgent functions, the restriction to integrable singularities is
removed by working with a larger space of singularities, equipped with a
suitable convolution product; see, for example,
\cite{Ecalle1981I,SauzinSplitting}.  In this larger setting, the exponential
factor is regarded as a formal action label in the transseries grading, and
is not part of the local Borel singularity on which the alien calculation is
performed.  Algebraic factors such as $\hbar^{-n/2}$ do have Borel
transforms, but these are generally non-integrable singularities supported at
the Borel origin, for instance derivatives of the convolution unit $\delta$
in the case of negative integral powers, or more general 
singularities for fractional powers.  Multiplication by such algebraic
factors on the $\hbar$-side corresponds to convolution with these origin
singularities on the Borel side.  This convolution is compatible with the
operations used in the alien calculation: analytic continuation, taking the
local variation at a non-zero point $\omega$, translation back to the
origin, and the inverse Borel transform.  Since the algebraic factor carries
no additional monodromy at $\omega\neq0$, it is recovered unchanged after
these operations. In the applications in this paper, this full singularity formalism gives the same
result as the simplified convention.
We therefore do not spell out the full formalism here.}

\[
\dot\Delta_\omega^+\big(
(4\pi\hbar)^{-n/2}
e^{-\Phi_0/\hbar}
\widetilde u(\hbar)\big)
=
(4\pi\hbar)^{-n/2}
e^{-(\Phi_0+\omega)/\hbar}
\Delta_\omega^+\widetilde u(\hbar),
\]
whenever the right-hand side is defined. The exponential factor restores the action shift carried by the Borel
singularity.  For instance, if a formal solution with exponential weight
$e^{-\Phi_0/\hbar}$ has a Borel singularity at
\[
\omega=\Phi_1-\Phi_0,
\]
then $\dot\Delta_{\omega}^{+}$ produces a contribution with weight $e^{-\Phi_1/\hbar}$.

\begin{lem}
\label{lem:pointed-alien-commutation}
Let $\re\alpha,\re\beta>0$. Let
$\widehat\phi(\xi,x)\in \xi^{\alpha-1}\C\{\xi\}$ and
assume that its Borel inverse $\widetilde\phi(\hbar,x)$ satisfies $\mathcal{B}\Delta^+_{\omega(x)}\widetilde\phi
\in
\xi^{\beta-1}\C\{\xi\}$ for fixed $x$. Assume that $\widehat\phi$, $\omega(x)$, and the chosen continuation paths
$\gamma_x^\pm$ depend analytically on $x$, in the sense that
\begin{equation}\label{eq:commuteswithpartialx}
\partial_x\circ \operatorname{cont}_{\gamma_x^\pm}
=
\operatorname{cont}_{\gamma_x^\pm}\circ \partial_x,
\qquad
\partial_x\circ \mathcal{B}
=
\mathcal{B} \circ \partial_x .
\end{equation}
Then
\begin{equation}
\bigl[\hbar^2\partial_\hbar,\dot\Delta_{\omega(x)}^{+}\bigr]\widetilde\phi=0,
\qquad
\bigl[\partial_x,\dot\Delta_{\omega(x)}^{+}\bigr]\widetilde\phi=0.
\end{equation}
\end{lem}


\begin{proof}
The first commutation is immediate from the definition of the alien operator and
\[
\mathcal B\bigl(\hbar^2\partial_\hbar\widetilde\phi\bigr)
=\xi\,\widehat\phi.
\]
We prove the second commutation.  Here we use the enlarged space of
Borel singularities discussed in the preceding footnote, so that for any $\widetilde f \in \hbar^{\beta}\C[[\hbar]]_1$,
\[
\partial_\xi\mathcal B\widetilde f
=
\mathcal{B}\hbar^{-1} \widetilde f.
\]
By the assumptions \eqref{eq:commuteswithpartialx} and the chain rule, 
\[
\partial_x
\mathcal B(\Delta_{\omega(x)}^+\widetilde\phi)
=\mathcal B\bigl(\Delta_{\omega(x)}^+(\partial_x\widetilde\phi)\bigr)
+(\partial_x\omega(x))\,\partial_\xi
\mathcal B(\Delta_{\omega(x)}^+\widetilde\phi).
\]
Here the second term on the right comes from the shift $\xi \to \xi+ \omega(x)$. Applying $\mathcal B^{-1}$, we obtain
\[
\partial_x\Delta_{\omega(x)}^+\widetilde\phi
=\Delta_{\omega(x)}^+(\partial_x\widetilde\phi)
+\frac{\partial_x\omega(x)}{\hbar}\Delta_{\omega(x)}^+\widetilde\phi .
\]
Therefore
\begin{align*}
\partial_x
\bigl(\dot\Delta_{\omega(x)}^+\widetilde\phi\bigr)
&=\partial_x\bigl(e^{-\omega(x)/\hbar}\Delta_{\omega(x)}^+\widetilde\phi\bigr)  \\
&=-\frac{\partial_x\omega(x)}{\hbar}
e^{-\omega(x)/\hbar}
\Delta_{\omega(x)}^+\widetilde\phi
+e^{-\omega(x)/\hbar}
\partial_x\Delta_{\omega(x)}^+\widetilde\phi  \\
&=-\frac{\partial_x\omega(x)}{\hbar}
e^{-\omega(x)/\hbar}
\Delta_{\omega(x)}^+\widetilde\phi
+e^{-\omega(x)/\hbar}\left(\Delta_{\omega(x)}^+(\partial_x\widetilde\phi)
+\frac{\partial_x\omega(x)}{\hbar}
\Delta_{\omega(x)}^+\widetilde\phi
\right)  \\
&=e^{-\omega(x)/\hbar}
\Delta_{\omega(x)}^+(\partial_x\widetilde\phi)
=\dot\Delta_{\omega(x)}^+(\partial_x\widetilde\phi). \qedhere
\end{align*} 
\end{proof}

\begin{cor}
\label{cor:alien-heat-solution-action-gap}
Let
\[
\widetilde K_0(\hbar;x,y)
=
(4\pi\hbar)^{-n/2}
e^{-\Phi_0(x,y)/\hbar}
\widetilde u_0(\hbar;x,y)
\]
be the formal heat solution associated with the minimal real geodesic branch,
as in Theorem~\ref{thm:local-borel-summability}. 
Assume that $\dot\Delta_\omega^+$ is defined on
$\widetilde K_0$ in the sense in Lemma~\ref{lem:pointed-alien-commutation}, and that the hypotheses \eqref{eq:commuteswithpartialx} hold for every local coordinate
direction $\partial_{x_i}$ appearing in $\Delta_x^{\C}$.  Then
\[
\left(\partial_\hbar+\Delta_x^{\C}\right)
\bigl(\dot\Delta_\omega^+\widetilde K_0\bigr)
=0.
\]
Moreover, let $\gamma$ be a no-conjugate holomorphic geodesic and assume that $\omega=\Phi_\gamma-\Phi_0 $. Then there exists a coefficient $S_{0\gamma}^+\in\C$ such that
\[
\dot\Delta_{\Phi_\gamma-\Phi_0}^{+}\widetilde K_0
=
S_{0\gamma}^+\,\widetilde K_\gamma ,
\]
where $\widetilde K_\gamma(\hbar;x,y)
=
(4\pi\hbar)^{-n/2}
e^{-\Phi_\gamma(x,y)/\hbar}
\widetilde u_\gamma(\hbar;x,y)$ is the normalized formal heat solution with phase $\Phi_\gamma$, uniquely
determined by the transport recursion and the chosen normalization as in Theorem~\ref{thm:heat-kernel-expansion-complex-geodesic}.
\end{cor}

\begin{proof}
This follows directly from Lemma~\ref{lem:pointed-alien-commutation}, which shows that the alien operator commutes with the scaled heat operator, together with Theorem \ref{thm:Borel-singularity-complex-geodesic}, which identifies the relevant singularity as the action gap $\omega=\Phi_\gamma-\Phi_0$, and Theorem~\ref{thm:heat-kernel-expansion-complex-geodesic}, which gives the normalized formal heat solution $\widetilde K_\gamma$.
\end{proof}

%{\color{red} (it is better to explain some basic idea on how to check "the hypotheses of
%Lemma~\ref{lem:pointed-alien-commutation} hold for every local coordinate
%direction $\partial_{x_i}"$)}


Whenever the proportionality relation
\[
\dot\Delta_{\Phi_\gamma-\Phi_0}^{+}\widetilde K_0
=
S_{0\gamma}^{+}\,\widetilde K_\gamma
\]
holds, we call $S_{0\gamma}^{+}$ the Stokes coefficient from $\widetilde K_0$ to 
$\widetilde K_\gamma$.  By
Corollary~\ref{cor:alien-heat-solution-action-gap}, once the sector
$\widetilde K_\gamma$ is fixed by the chosen normalization, this coefficient
is a well-defined complex number.

The expected integrality of $S_{0\gamma}^{+}$ is a jump number of Lefschetz thimble rather than a purely formal consequence of alien calculus.  In
finite-dimensional exponential integrals, the number computeded from alien culculus agrees with the Picard--Lefschetz wall-crossing coefficient; after
choosing orientations, this coefficient is an intersection number, or
equivalently the signed count of direct connecting trajectories at the
corresponding Stokes phase \cite{LiLiTang2026}. 

\begin{conj}
\label{conj:alien-PL-integrality-heat}
Let $(M,g)$ be real analytic, and let $(M_{\C},g_{\C})$ be a
complexification satisfying our assumption \ref{assumption}.  Fix two endpoints $x,y$, and let
$\{\gamma_i\}_{i\in I}$ be the relevant no-conjugate holomorphic geodesic from $x$ to $y$.  Denote their phases by $\Phi_i(x,y):=\frac{\ell_{\gamma_i}(x,y)^2}{4}$. We conjecture that each $\gamma_i$ determines a normalized formal heat
solution of the form
\[
\widetilde K_i(\hbar;x,y)
=
(4\pi\hbar)^{-n/2}
e^{-\Phi_i(x,y)/\hbar}
\widetilde u_i(\hbar;x,y),
\qquad
\widetilde u_i(\hbar;x,y)\in \C[[\hbar]],
\]
where $\widetilde u_i$ is a $1$--Gevrey formal series. Moreover, there exists coefficient $S_{ij}^{+}\in\Z$ such that
\[
\dot\Delta_{\Phi_j-\Phi_i}^{+}\widetilde K_i
=
S_{ij}^{+}\,\widetilde K_j.
\]
\end{conj}

The hyperbolic heat-kernel model studied in the next section provides the
first test case for Conjecture~\ref{conj:alien-PL-integrality-heat}.  Starting
from the principal real WKB sector $\widetilde K_0$, we will compute the
pointed alien contribution at the first nontrivial action gap.  The
resurgent calculation gives the coefficient (see Section \ref{sec:H^2-resurgence})
\[
S_{01}^{+}=2.
\]
On the Picard--Lefschetz side, for two sufficiently close real endpoints, the
path-space Morse analysis shows that the corresponding wall crossing is
governed by two direct connecting trajectories.  In this sense, the
hyperbolic model gives a concrete local verification of the expected
integrality mechanism.




\part{A case study of $\mathbb{H}_{\R}^2$ }

The rest of this paper is devoted to a case study of  the hyperbolic space $\mathbb{H}_{\R}^2$ in the context of the Picard–Lefschetz problem and heat kernel resurgence outlined in Section \ref{sec:heat-resurgence}. We start with the resurgence analysis in Section \ref{sec:H^2-resurgence}. In Section \ref{sec:H^2-flow}, we perform the Picard–Lefschetz analysis and compare with the resurgence computations. 



\section{Resurgence analysis}\label{sec:H^2-resurgence}


In this section, we analyze the resurgent structure of the hyperbolic heat kernel
through the complex length branches of the hyperbolic distance.  We first
construct, directly from the radial heat equation, a normalized formal solution
$\widetilde{K}_k$ associated with each branch
$\ell_k=d+2\pi i k$, as described in Theorem \ref{thm:heat-kernel-expansion-complex-geodesic}.  %\sout{After extracting the standard leading factor $\left(\frac{\ell_k}{\sinh\ell_k}\right)^{1/2}$, the remaining power series is uniquely fixed by the transport recursion and the regularity condition at the complex-length origin.  This normalization agrees with the usual WKB normalization, and we shall therefore refer to these objects as WKB-type formal objects.}

We then return to the principal real-geodesic expansion and study its Borel
singularities.  At the $k$-th singular point $\omega_k(d)$, the local Borel-plane
variation produces the same normalized formal object
$\widetilde{K}_k$, with an overall coefficient $2$.  Equivalently,
\[
\dot\Delta_{\omega_k(d)}^{+}\,
\widetilde{K}_0=2\,\widetilde{K}_k.
\]
Thus the hyperbolic Stokes constant is $2$, in direct analogy with the
finite-dimensional alien/Picard--Lefschetz correspondence discussed in
\cite{LiLiTang2026}.


\subsection{Short time asymptotic expansion}

Let $a,b\in\HH_{\R}^{2}$, and write $d$ for their geodesic distance $d=d(a,b)$. By homogeneity and isotropy, the heat kernel depends only on the time $\tau$ and $d$. %we denote it by
%\[K_{\HH_{\R}^{2}}(\tau,d).\]

\smallskip
As we explained in the introduction, $\tau$ plays the role of $\hbar$ in the saddle problem of exponential integral. In the remainder of this section, we shall use $\hbar$ to denote the time $\tau$  to align with our finite-dimensional discussion in \cite{LiLiTang2026}. Thus the heat kernel will be denoted instead by 
\[K_{\HH_{\R}^{2}}(\hbar,d).\]

A classical integral representation is
\begin{equation*}\label{eq:H2-heat-kernel-s}
K_{\HH_{\R}^{2}}(\hbar,d)
=\frac{\sqrt{2}\,e^{-\hbar/4}}{(4\pi\hbar)^{3/2}}
\int_{d}^{\infty}
\frac{s\,e^{-s^{2}/4\hbar}}
{\sqrt{\cosh s-\cosh d}}
\,\dd s.
\end{equation*}
After the change of variables
\[
w=\frac{s^{2}}{2},
\qquad
\dd w=s\,\dd s,
\]
this becomes
\begin{equation}\label{eq:H2-heat-kernel-w}
K_{\HH_{\R}^{2}}(\hbar,d)
=\frac{\sqrt{2}\,e^{-\hbar/4}}{(4\pi\hbar)^{3/2}}
\int_{d^{2}/2}^{\infty}
\frac{e^{-w/2\hbar}}
{\sqrt{\cosh\sqrt{2w}-\cosh d}}
\,\dd w.
\end{equation}
We refer to \cite{Du} for a resurgence analysis of this representation.

\smallskip

The general $1$--Gevrey discussion of Section \ref{subsec:proof-borel-summability} provides the
formal framework.  In the present hyperbolic model, the exact representation
\eqref{eq:H2-heat-kernel-w} makes the relevant Borel transform explicit.

We use the half-integer Borel convention
\[
\widetilde\phi(\hbar)
=\sum_{n\ge0}c_n\,\hbar^{n+\frac12}
\quad \mapsto\quad
\widehat\phi(\xi)
=\mathcal B\widetilde\phi(\xi)
:=\sum_{n\ge0}
\frac{c_n}{\Gamma\!\left(n+\frac12\right)}
\,\xi^{\,n-\frac12}.
\]
If $\widehat\phi$ admits analytic continuation along the ray 
$e^{i\theta}\R_+$, without singularity on that ray, and has at most
exponential growth there, then we set
\begin{equation*}\label{eq:directional-laplace}
\bigl(\mathcal L^\theta\widehat\phi\bigr)(\hbar)
:=\int_0^{e^{i\theta}\infty}
e^{-\xi/\hbar}\widehat\phi(\xi)\,\dd\xi.
\end{equation*}
Then, on
\[
S_\theta=\left\{\hbar\in\C^*:\left|\arg\hbar-\theta\right|<\frac{\pi}{2}\right\},
\]
we have uniformly on every proper subsector,
\begin{equation}\label{eq:asymptoticBL}
\bigl(\mathcal L^\theta\widehat\phi\bigr)(\hbar)\sim\sum_{n\ge0}c_n\,\hbar^{n+\frac12},
\qquad\hbar\to0.
\end{equation}

\smallskip

Returning to \eqref{eq:H2-heat-kernel-w}, set
\[w=\frac{d^2}{2}+2\xi.\]
Then
\begin{equation}\label{eq:H2-heat-kernel-borel-laplace}
K_{\HH_{\R}^{2}}(\hbar,d)
=\frac{2\sqrt{2}\,e^{-\hbar/4}}{(4\pi\hbar)^{3/2}}
\exp\!\left(-\frac{d^2}{4\hbar}\right)
\bigl(\mathcal L^0\widehat\phi_d\bigr)(\hbar),
\end{equation}
where
\begin{equation}\label{eq:H2-borel-germ}
\widehat\phi_d(\xi)
=\left(\cosh\sqrt{d^2+4\xi}-\cosh d\right)^{-1/2}.
\end{equation}
Near $\xi=0$,
\[
\widehat\phi_d(\xi)
=\left(\frac{d}{2\sinh d}\right)^{1/2}
\xi^{-1/2}\sum_{n\ge0}c_n(d)\,\xi^n,
\qquad c_0(d)=1.
\]
Applying \eqref{eq:asymptoticBL} gives
\begin{equation*}\label{eq:H2-short-time-expansion}
K_{\HH_{\R}^{2}}(\hbar,d)
\sim
\frac{e^{-\hbar/4}}{4\pi\hbar}
\exp\!\left(-\frac{d^{2}}{4\hbar}\right)
\left(\frac{d}{\sinh d}\right)^{1/2}
\sum_{n\ge0}a_n(d)\,\hbar^n,
\qquad
\hbar\to0^+,
\end{equation*}
with
\begin{equation*}\label{eq:H2-an-cn-relation}
a_n(d)=\frac{\Gamma\left(n+\frac12\right)}{\sqrt{\pi}}\,c_n(d).
\end{equation*}

The coefficients $a_n(d)$ admit a closed generating formula.  Define
\[
q(d):=\frac{\sinh d}{d},
\qquad\mathcal D:=\frac1d\frac{\dd}{\dd d}.
\]
Since
\[
\cosh\sqrt{d^2+4\xi}=e^{2\xi\mathcal D}\cosh d,
\]
we obtain
\[
\cosh\sqrt{d^2+4\xi}-\cosh d
=2\xi\,q(d)\left(1+\sum_{m\ge1}
\frac{(2\xi)^m}{(m+1)!}
\frac{\mathcal D^m q(d)}{q(d)}\right).
\]
Consequently,
\begin{equation*}\label{eq:H2-an-closed-formula}
a_n(d)=\frac{\Gamma\!\left(n+\frac12\right)}{\sqrt{\pi}}\,[\xi^n]\,\left(1+\sum_{m\ge1}
\frac{(2\xi)^m}{(m+1)!}\frac{\mathcal D^m q(d)}{q(d)}\right)^{-1/2},\qquad n\ge0.
\end{equation*}
In particular,
\[a_0(d)=1.\]
For $d=0$, the formula is understood by regular continuation from $d>0$.

\subsection{Formal objects from the radial heat equation}

We now define the formal objects associated with the complex length branches
directly from the radial heat equation.

Let
\[
\ell_k:=d+2\pi ik,\qquad\Phi_k(d):=\frac{\ell_k^2}{4},
\qquad k\in\Z.
\]
The branch $k=0$ is the real branch, for which $\ell_0=d$ and
$\Phi_0(d)=d^2/4$.

In accordance with Theorem \ref{thm:heat-kernel-expansion-complex-geodesic}, we have the following explicit construction. 

\begin{prop}[Formal heat objects and transport recursion]
\label{prop:H2-formal-heat-recursion}

For each branch $\ell=\ell_k$, there is a formal object of the form
\begin{equation}\label{eq:H2-formal-heat-object-k}
\widetilde{K}_k(\hbar,d)
=
\frac{e^{-\hbar/4}}{4\pi\hbar}
\exp\!\left(-\frac{\Phi_k(d)}{\hbar}\right)
\left(\frac{\ell_k}{\sinh\ell_k}\right)^{1/2}
\sum_{n\ge0}A_n(\ell_k)\hbar^n,
\end{equation}
obtained by substituting it 
into the radial heat equation 
\[
(\partial_\hbar+\Delta_{\mathrm{rad},\ell})
\widetilde{K}=0,\qquad 
\Delta_{\mathrm{rad},\ell}=-\partial_\ell^2-\coth \ell\,\partial_\ell.
\]
The coefficients are uniquely determined by
\begin{equation}\label{eq:H2-transport-recursion-An}
A_0(\ell)=1,\qquad \left(\ell\partial_\ell+n\right)A_n(\ell) 
=\mathcal R_\ell A_{n-1}(\ell),\ \text{for} \  n\geq1,
\end{equation}
where
\[
\mathcal R_\ell:=\partial_\ell^2+\frac1\ell\partial_\ell
+\frac{1}{4\sinh^2\ell}-\frac{1}{4\ell^2}.
\]
\end{prop}

\begin{proof}
Work first with the auxiliary variable $\ell$, and write
\[
\widetilde{K}(\hbar,\ell)
=\frac{e^{-\hbar/4}}{4\pi\hbar}
e^{-\ell^2/4\hbar}\left(\frac{\ell}{\sinh\ell}\right)^{1/2}A(\hbar,\ell).
\]
Substitution into
\[
(\partial_\hbar+\Delta_{\rad,\ell})\widetilde{K}=0
\]
gives, after cancellation of the exponential and leading amplitude factors,
\[
\left(\hbar\partial_\hbar+\ell\partial_\ell\right)A=\hbar\,\mathcal R_\ell A.
\]
Writing
\[
A(\hbar,\ell)=\sum_{n\ge0}A_n(\ell)\hbar^n
\]
and comparing powers of $\hbar$ gives
\[
\ell\partial_\ell A_0=0,\qquad
(\ell\partial_\ell+n)A_n=\mathcal R_\ell A_{n-1}\quad(n\ge1).
\]
We fix the normalization by taking $A_0=1$.

For uniqueness, suppose $A_{n-1}$ is already fixed.  Then
\[
(\ell\partial_\ell+n)A_n=F_n(\ell)
\]
has the general solution
\[
A_n(\ell)=\ell^{-n}
\left(C_n+\int_0^\ell u^{n-1}F_n(u)\,du\right).
\]
Regularity at $\ell=0$ forces $C_n=0$.  Hence $A_n$ is uniquely
determined at each order.  Finally, evaluate at $\ell=\ell_k$.
\end{proof}

\begin{rmk}\label{rmk:Normalizationconvetion}
The factor $\frac{\ell_k}{\sinh\ell_k}$ is the Jacobian of the complexified exponential map. This agrees with the normalization convention in Theorem \ref{thm:heat-kernel-expansion-complex-geodesic}. It also agrees with the usual WKB normalization.  In
the WKB expansion at the holomorphic geodesic branch $\gamma_k$, one first
extracts the exponential factor $e^{-\Phi_k(d)/\hbar}$ and the one-loop
determinant
\[
\left[\frac{\det_D(-\partial_t^2)}
{\det_D(-\partial_t^2+\ell_k^2)}\right]^{1/2}
=\left(\frac{\ell_k}{\sinh\ell_k}\right)^{1/2}.
\]
The remaining WKB series is then normalized to be $1+O(\hbar)$.  %This is exactly the convention used in \eqref{eq:H2-formal-heat-object-k}: after extracting the same leading factor, we impose $A_0(\ell_k)=1$.  With this common normalization, the higher coefficients are fixed by the same radial heat-equation recursion, so the two formal objects agree.
\end{rmk}


\subsection{Borel singularities and pointed alien operators in the hyperbolic case}

We now apply the alien-operator convention above to the principal Borel germ
$\widehat\phi_d$ in \eqref{eq:H2-borel-germ}.  
We have seen in Example \ref{eg:complex-geodesics-H2} that the nonzero Borel singularities are
\[\omega_k(d)=\frac{(d+2\pi i k)^2-d^2}{4},\qquad k\in\Z\setminus\{0\}.\]
Using the notation introduced in Section \ref{sec:heat-alien},
\[\omega_k(d)=\Phi_k(d)-\Phi_0(d).\]
Thus the Borel singularities occur precisely at the phase differences between
the principal real geodesic and the non-principal holomorphic geodesic branches.

Writing
\[\xi=\omega_k(d)+\zeta,\]
for any $k\in\mathbb{Z}\setminus\{0\}$, the translated local germ of $\widehat\phi_d$ at $\omega_k(d)$ is therefore
\[\widehat\phi_{d,k}(\zeta)
:=\left(\cosh\sqrt{\ell_k^2+4\zeta}
-\cosh\ell_k\right)^{-1/2}.\]
The branch is fixed by analytic continuation of $\widehat\phi_d$ along the
positive path $\gamma_{\omega_k}^{+}$. (This notation should not be confused with our notation for geodesics.
) See Figure \ref{fig:hyperbolicsingularity} below. 



\begin{figure}[H]
\centering
\resizebox{0.67\textwidth}{!}{%
\begin{tikzpicture}[scale=0.1,>=stealth]
% parameter d (only affects the vertical scaling of the parabola)
\def\dpar{1.2}

% singular points: omega_k = -pi^2 k^2 + i pi d k
\coordinate (w0)  at (0,0);
\coordinate (w1)  at ({-pi*pi*1^2},{ pi*\dpar*1});
\coordinate (w2)  at ({-pi*pi*2^2},{ pi*\dpar*2});
\coordinate (w3)  at ({-pi*pi*3^2},{ pi*\dpar*3});
\coordinate (wm1) at ({-pi*pi*1^2},{-pi*\dpar*1});
\coordinate (wm2) at ({-pi*pi*2^2},{-pi*\dpar*2});
\coordinate (wm3) at ({-pi*pi*3^2},{-pi*\dpar*3});

% endpoint of gamma_2^+ : on the parabola, just beyond omega_2
% choose a nearby point corresponding to t = 2.12
\def\tend{2.12}
\coordinate (zend) at ({-pi*pi*\tend*\tend},{pi*\dpar*\tend});

% axes

% dashed parabola x = -y^2/d^2
\draw[gray,dashed,thick,domain=-11.5:11.5,samples=200,smooth]
  plot ({-(\x*\x)/(\dpar*\dpar)}, {\x});

% singularities (red crosses)
\draw[red,thick] ($(w0)+(-0.35,-0.35)$) -- ($(w0)+(0.35,0.35)$);
\draw[red,thick] ($(w0)+(-0.35,0.35)$) -- ($(w0)+(0.35,-0.35)$);

\foreach \P in {w1,w2,w3,wm1,wm2,wm3}{
  \draw[red,thick] ($(\P)+(-0.45,-0.45)$) -- ($(\P)+(0.45,0.45)$);
  \draw[red,thick] ($(\P)+(-0.45,0.45)$) -- ($(\P)+(0.45,-0.45)$);
}

% labels for singularities
\node[red,above right] at (w0) {$0$};

\node[red,above right] at (w1) {$\omega_1(d)$};
\node[red,below right] at (w2) {$\omega_2(d)$};
\node[red,above right] at (w3) {$\omega_3(d)$};

\node[red,below right] at (wm1) {$\omega_{-1}(d)$};
\node[red,below right] at (wm2) {$\omega_{-2}(d)$};
\node[red,below right] at (wm3) {$\omega_{-3}(d)$};

% gamma_2^+ path: starts at 0, passes above omega_1 and omega_2,
% and ends near omega_2 on the parabola
\draw[black,thick,->]
  (w0)
  .. controls (-3.5,1.8) and (-7.5,4.8) ..
  (-10.5,5.7)
  .. controls (-20,8.8) and (-33,10.5) ..
  (-43.0,8.9)
  .. controls (-44.0,8.45) and (-44.25,8.18) ..
  (zend);

\node[black,above] at (-18,8.9) {$\gamma_{\omega_2}^+$};

% green circle centered at the endpoint, radius 2 epsilon / 3
\draw[green!60!black,thick] (zend) circle [radius=1.9];

% optional endpoint marker
\fill[black] (zend) circle (1.1pt);

\end{tikzpicture}}
\caption{Borel singularities of $\widehat{\phi}_d$ on the parabola
$\omega_k(d)=\frac{(d+2\pi i k)^2-d^2}{4}$, together with the path
$\gamma_{\omega_2}^+$.}
\label{fig:hyperbolicsingularity}
\end{figure}


Since $\widehat\phi_{d,k}$ has a square-root singularity at
$\zeta=0$, a small clockwise loop around $\omega_k(d)$ changes its sign.
With the positive variation convention,
\begin{equation}\label{eq:H2-local-variation}
\tau_{-\omega_k(d)}\,
\var_{\omega_k(d)}^{+}\widehat\phi_d
=
2\,\widehat\phi_{d,k}.
\end{equation}
Applying the inverse Borel transform gives
\[
\Delta_{\omega_k(d)}^{+}\widetilde\Psi_0
=2\,\widetilde\Psi_{d,k},\qquad
\widetilde\Psi_{d,k}
:=\mathcal B^{-1}\widehat\phi_{d,k},
\qquad \widetilde\Psi_{0}
:=\mathcal B^{-1}\widehat\phi_{d}.
\]

The leading term of $\widehat\phi_{d,k}$ is
\[
\widehat\phi_{d,k}(\zeta)=\left(\frac{\ell_k}{2\sinh\ell_k}\right)^{1/2}\zeta^{-1/2}\bigl(1+O(\zeta)\bigr),
\]
and hence
\[
\widetilde\Psi_{d,k}(\hbar)
=\sqrt{\pi}\left(\frac{\ell_k}{2\sinh\ell_k}
\right)^{1/2}\hbar^{1/2}\bigl(1+O(\hbar)\bigr).
\]
After restoring the common prefactor in
\eqref{eq:H2-heat-kernel-borel-laplace}, this reproduces the one-loop
normalization of $\widetilde{K}_k$ in
Remark \ref{rmk:Normalizationconvetion}.  To compare the full WKB formal
solutions, we shall use the following commutation property.

\begin{prop}
\label{prop:H2-pointed-alien-stokes-constant}
Fix $k\in\Z\setminus\{0\}$.  Let
$\widetilde{K}_k$ be the WKB-type formal solution attached to
$\gamma_k$, normalized as in
\eqref{eq:H2-formal-heat-object-k} and \eqref{eq:H2-transport-recursion-An}.  Then
\[
\dot\Delta_{\omega_k(d)}^{+}\,
\widetilde{K}_0=2\,\widetilde{K}_k.
\]
Equivalently, the positive Stokes constant from the principal real branch to
the $k$-th holomorphic geodesic branch is
$
S_{0k}^{+}=2.
$ In particular,
\[
\dot\Delta_{\omega_1(d)}^{+}\,
\widetilde{K}_0=2\,\widetilde{K}_1.
\]
\end{prop}

\begin{proof}
By \eqref{eq:H2-local-variation},
\[
\tau_{-\omega_k(d)}\,
\var_{\omega_k(d)}^{+}\widehat\phi_d
=2\,\widehat\phi_{d,k}.
\]
Using the leading expansion of $\widehat\phi_{d,k}$ and restoring the
prefactor in \eqref{eq:H2-heat-kernel-borel-laplace}, we obtain
\begin{equation}\label{eq:H2-pointed-alien-leading-solution}
\dot\Delta_{\omega_k(d)}^{+}\,
\widetilde{K}_0
=2\,\frac{e^{-\hbar/4}}{4\pi\hbar}
\exp\!\left(-\frac{\Phi_k(d)}{\hbar}
\right)\left(\frac{\ell_k}{\sinh\ell_k}
\right)^{1/2}\bigl(1+O(\hbar)\bigr).
\end{equation}
Hence, after division by $2$, the pointed alien contribution has the same
phase and the same one-loop leading factor as
$\widetilde{K}_k$.

It remains to identify the full formal power series.  The principal WKB
solution $\widetilde{K}_0$ satisfies the radial heat equation, or
equivalently
\[
\mathfrak H\widetilde{K}_0=0,\qquad
\mathfrak H:=\hbar^2\partial_\hbar
-\hbar^2\left(\partial_d^2+\coth d\,\partial_d
\right).
\]
By Lemma \ref{lem:pointed-alien-commutation},
$\dot\Delta_{\omega_k(d)}^{+}$ commutes with both $\hbar^2\partial_\hbar$ and $\partial_d$.  Therefore
\[
\mathfrak H
\left(\dot\Delta_{\omega_k(d)}^{+}\widetilde{K}_0\right)=0.
\]
Thus the pointed alien contribution is again a WKB-type formal solution of
the radial heat equation, now with phase $\Phi_k(d)$.

Moreover, after dividing by $2$, its reduced formal power series is obtained
from the local germ $\widehat\phi_{d,k}$.  After extracting the universal
factor
\[
\left(\frac{\ell_k}{2\sinh\ell_k}
\right)^{1/2}\zeta^{-1/2},
\]
the remaining Taylor coefficients are regular at $\ell=0$.  Hence this
formal solution satisfies the same regularity normalization used to define
$\widetilde{K}_k$.

The phase, the one-loop leading factor, and the regularity condition at
$\ell=0$ uniquely determine the WKB-type formal solution attached to
$\gamma_k$.  Comparing with
\eqref{eq:H2-pointed-alien-leading-solution} proves
\[
\dot\Delta_{\omega_k(d)}^{+}\,
\widetilde{K}_0=2\,\widetilde{K}_k.\qedhere
\]
\end{proof}

\begin{cor}
\label{cor:H2-pointed-alien-general-transition}
Let $j,m\in\Z$ with $j\neq m$, and set
\[
\omega_{j\to m}(d):=\Phi_m(d)-\Phi_j(d).
\]
Then
\begin{equation}\label{eq:H2-general-pointed-alien-transition}
\dot\Delta_{\omega_{j\to m}(d)}^{+}\,
\widetilde{K}_j=2\,\widetilde{K}_m.
\end{equation}
Thus every direct pointed alien transition between two distinct complex
geodesic branches carries the same Stokes constant $2$.
\end{cor}


\begin{proof}
Let $\widehat\phi_{d,j}$ be the reduced Borel germ associated with
$\widetilde K_j$.  We use Proposition \ref{prop:H2-pointed-alien-stokes-constant} after re-centering the
construction at the branch $\gamma_j$. Then the branch $\gamma_m$ appears at the Borel-plane
displacement
\[
\omega_{j\to m}(d)=\Phi_m(d)-\Phi_j(d).
\]
The radial heat equation, the transport recursion, and the normalization
$A_0=1$ are unchanged under this re-centering.  Therefore
Proposition \ref{prop:H2-pointed-alien-stokes-constant} gives
\[
\tau_{-\omega_{j\to m}(d)}\operatorname{var}^{+}_{\omega_{j\to m}(d)}\widehat\phi_{d,j}
=2\,\widehat\phi_{d,m}.
\]
In particular, the singularity at $\omega_{j\to m}(d)$ is of square-root type by
Proposition \ref{prop:H2-pointed-alien-stokes-constant}, with $\gamma_j$ as the reference branch.

After inverse Borel transform and insertion of the pointed factor, the exponential weight changes from $e^{-\Phi_j(d)/\hbar}$ to
$e^{-\Phi_m(d)/\hbar}$.
Now the translated reduced germ is $2\widehat\phi_{d,m}$, and the one-loop
factor and the transport normalization are the same as those used in the
definition of $\widetilde K_m$.  Hence
\[
\dot\Delta^{+}_{\omega_{j\to m}(d)}\widetilde K_j
=2\,\widetilde K_m.\qedhere
\]
\end{proof}


The coefficient $2$ is the resurgent quantity that will later be compared
with the connecting-trajectory count for the corresponding pair of critical
paths in the complexified path-space Morse theory.



\section{Picard–Lefschetz analysis}\label{sec:H^2-flow}
In this section, we will study the $\HH_{\R}^2$ model from Picard-Leschetz side on the complexified path space. Let us first recall necessary notations from Example \ref{eg:H^2-PL}. Let
\[
\HH_{\C}^2:=\bigl\{z=(z_0,z_1,z_2)\in \C^3:\ -z_0^2+z_1^2+z_2^2=-1\bigr\},
\qquad
\eta:=\diag(-1,1,1).
\]
After a real Lorentz isometry, we normalize the two real endpoints as
\[
a=(1,0,0),
\qquad
b_d=(\cosh d,\sinh d,0).
\]
We consider the endpoint-fixed Sobolev path space
\[
\PP_d^{\C}:=\PP_{a,b_d}(\mathbb{H}_{\C}^2)=\bigl\{z\in H^1([0,\tau],\HH_{\C}^2)\mid z(0)=a,\ z(\tau)=b_d\bigr\}
\]
equipped with the action by the holomorphic energy functional
$$
S(z)=\frac12\int_0^{\tau} z_t^T\eta z_t\,\dd t=\frac{1}{2}\int_0^{\tau}\langle z_t,z_t\rangle_{\eta}\,\dd t.
$$
\begin{rmk} The factor $1/2$ is the traditional normalization for the energy. In Section \ref{sec:PL-problem-path}, we have used the factor $1/4$ to define $\mathcal E_{\C}$ so that the exponential $\exp(-\mathcal E_{\C})$ leads to the heat kernel.  Since
$$
S=2\mathcal E_{\C},
$$
they give rise to equivalent Morse flows. We choose to work with $S$ instead in this section for Picard–Lefschetz and Morse flow analysis. It simplifies expressions in various computations. 
\end{rmk}

Let $\theta$ be a chosen phase.  Define the twisted real part and imaginary part of $S$ by
\begin{align*}
F_{\theta}(z)=&\re\bigl(e^{-i\theta}S(z)\bigr),
\qquad\quad
G_{\theta}(z)=\im\bigl(e^{-i\theta}S(z)\bigr).    
\end{align*}
The downward Morse flow system of $F_{\theta}$ is computed in Example \ref{eg:H^2-PL} and given by
\begin{equation}\label{eq:flow}
\begin{cases}
\partial_sz=P_z(e^{i\theta}\eta \bar{z}_{tt})=e^{i\theta}\left(\eta\bar{z}_{tt}-\frac{z^T\bar{z}_{tt}}{|z|^2}\eta\bar{z}\right)\\[0.1cm]
\langle z(s,t),z(s,t)\rangle_{\eta}=z(s,t)^T\eta z(s,t)=-1\\[0.1cm]
z(s,0)=a,~z(s,\tau)=b_d
\end{cases}.
\end{equation}
%Note that since we work with the complexified path space, the equation \eqref{eq:flow} cannot be treated as a standard parabolic initial-value problem in the $s$-variable. {\color{red}(what does this mean, in real components?)}









\subsection{The holomorphic geodesics and the complex Jacobi operator}

For each $j\in\Z$, set
\[
c_{j,d}:=\frac{d+2\pi i j}{\tau}.
\]
As shown in Example \ref{eg:H^2-PL}, for each $j$, there corresponds a holomorphic geodesics joining $a$ to $b_d$ by  
\begin{equation}\label{eq:gammajd}
\gamma_{j,d}(t)=\bigl(\cosh(c_{j,d}t),\,\sinh(c_{j,d}t),\,0\bigr),
\qquad 0\le t\le \tau.
\end{equation}
Moreover,
\[
\frac{\dd^2}{\dd t^2}\,\gamma_{j,d}=c_{j,d}^2\gamma_{j,d}.
\]

Along $\gamma_{j,d}$, introduce the tangent and normal fields
\[
T_{j,d}(t):=\frac1{c_{j,d}}\,\frac{\dd}{\dd t}\gamma_{j,d}(t)
=\bigl(\sinh(c_{j,d}t),\,\cosh(c_{j,d}t),\,0\bigr),
\qquad
N:=(0,0,1).
\]
We have
\[
\ip{\gamma_{j,d}}{T_{j,d}}_{\eta}=0,
\qquad
\ip{\gamma_{j,d}}{N}_{\eta}=0,
\qquad
\ip{T_{j,d}}{T_{j,d}}_{\eta}=1,
\qquad
\ip{N}{N}_{\eta}=1.
\]
Therefore every tangent variation along $\gamma_{j,d}$ can be written as
\[
\xi(t)=\alpha(t)T_{j,d}(t)+\beta(t)N,
\qquad
\alpha(0)=\alpha(\tau)=\beta(0)=\beta(\tau)=0.
\]


\begin{prop}\label{prop:nondegenerate}
The $\C$-linear second-variation operator of the holomorphic action at $\gamma_{j,d}$ is
\begin{equation*}\label{eq:Ajd}
L_{j,d}^{\C}(\alpha,\beta)=\bigl(-\alpha_{tt},\,-\beta_{tt}+c_{j,d}^2\beta\bigr)
\end{equation*}
on the Dirichlet domain
\[
\mathrm{Dom}(L_{j,d}^{\C})=\bigl(H^2\cap H^1_0([0,\tau],\C)\bigr)\oplus \bigl(H^2\cap H^1_0([0,\tau],\C)\bigr).
\]
Equivalently,
\[
L_{j,d}^{\C}(\alpha T_{j,d}+\beta N)=(-\alpha_{tt})T_{j,d}+(-\beta_{tt}+c_{j,d}^2\beta)N.
\]
Furthermore, for $d>0$, $\ker L_{j,d}^{\C}=\{0\}$, hence every critical path $\gamma_{j,d}$ is nondegenerate. 
\end{prop}


\begin{proof} 
Consider an arbitrary endpoint-fixed variation of $\gamma_{j,d}$ parametrized by small $\varepsilon$
\[u(\varepsilon,t)\in \HH_{\C}^2,\qquad u(0,t)=\gamma_{j,d}(t).\]
Let
\[\xi(t)=\partial_{\varepsilon}u(0,t),\qquad \zeta(t)=\partial_{\varepsilon}^2u(0,t)\]
be the first and second order variations, which satisfy the boundary condition
\[\xi(0)=\xi(\tau)=0,\qquad \zeta(0)=\zeta(\tau)=0.\]
Then we can compute
\begin{align*}
    \frac{\dd^2}{\dd\varepsilon^2}S(u(\varepsilon))\big|_{\varepsilon=0}=~&\int_0^{\tau}\left(\langle\xi_t,\xi_t\rangle_{\eta}+\langle(\gamma_{j,d})_t,\zeta_t\rangle_{\eta}\right)\,\dd t\\
    =~&\int_0^{\tau}\langle\xi_t,\xi_t\rangle_{\eta}\dd t-\int_0^{\tau}\langle(\gamma_{i,d})_{tt},\zeta\rangle_{\eta}\,\dd t\\
    =~&\int_0^{\tau}\langle\xi_t,\xi_t\rangle_{\eta}\dd t-c_{j,d}^2\int_0^{\tau}\langle\gamma_{j,d},\zeta\rangle_{\eta}\,\dd t\\
       =~&\int_0^{\tau}(\alpha_t^2+\beta_t^2-c_{j,d}^2 \alpha^2)\ \dd t-c_{j,d}^2\int_0^{\tau}\langle\gamma_{j,d},\zeta\rangle_{\eta}\,\dd t.
\end{align*}

Since $u(\varepsilon,t)\in\HH_{\C}^2$, 
\[\langle u(\varepsilon,t),u(\varepsilon,t)\rangle_{\eta}=-1.\]
Taking the derivatives with respect to $\varepsilon$ twice, we obtain
\[\langle\gamma_{j,d},\zeta\rangle_{\eta}+\langle\xi,\xi\rangle_{\eta}=0.\]
Thus for $\xi(t)=\alpha(t)T_{j,d}(t)+\beta(t)N$, we arrive at
\[\delta^2S_{\gamma_{j,d}}(\xi,\xi)=\int_0^{\tau}\left(\alpha_t^2+\beta_t^2+c_{j,d}^2\beta^2\right)\,\dd t=\int_0^{\tau}[\alpha(-\alpha_{tt})+\beta(-\beta_{tt}+c_{j,d}^2\beta)]\,\dd t.\]
This shows the formula for $L_{j,d}^{\C}$. 

We next show $\ker L_{j,d}^{\C}=\{0\}$. The tangential equation $-\alpha_{tt}=0$ with Dirichlet boundary conditions has only the trivial solution. Solutions of the normal equation $\beta_{tt}=c_{j,d}^2\beta$
take the form 
\[\beta(t)=c_1 e^{c_{j,d}t}+c_2 e^{-c_{j,d}t}.\] 
Imposing $\beta(0)=0$ gives $c_2=-c_1$, hence
\[
\beta(t)=2c_1\sinh(c_{j,d}t).
\]
Now $\beta(\tau)=0$ implies
\[
2c_1\sinh(c_{j,d}\tau)=2c_1\sinh(d+2\pi ij)=(-1)^j\times 2c_1\sinh d=0.
\]
Since $d>0$, $\sinh d\ne0$, we get $c_1=0$, that is, $\beta\equiv 0$.
\end{proof}





\begin{rmk}
The normal Dirichlet spectra of $L_{j,d}^{\C}$ are
\[\left(\frac{n\pi}{\tau}\right)^2+c_{j,d}^2,\quad n\ge1.\]
At $d=0$, the normal block becomes degenerate exactly when
\[
\Bigl(\frac{n\pi}{\tau}\Bigr)^2+\Bigl(\frac{2\pi ij}{\tau}\Bigr)^2=0,
\]
that is, when $n=2|j|$. In this case 
$$
\gamma_{j,d}(t)=\bigl(\cos(\frac{jt}{\tau}),\,i \sin(\frac{jt}{\tau}),\,0\bigr)
$$
lies on the sphere $S^2\subset \HH_{\C}^2$, and this is the familiar Morse--Bott phenomenon on $S^2$. This observation will play an essential role in our perturbation analysis in Section \ref{sec:case d=0}. 
\end{rmk}





\subsection{Linearized flow equation}

In this subsection, we will analyze linearization of the downward gradient flow equation of $F_{\theta}$ at the critical path $\gamma_{j,d}$, which plays an important role in the study of local structure of the Lefschetz thimble in the following subsections. 

\begin{prop}\label{prop:linearization}
At the critical path $\gamma_{j,d}$, the linearization of the flow equation \eqref{eq:flow} is
\begin{equation}\label{eq:linflow}
\partial_s \xi + A_{j,d,\theta}\xi=0,
\end{equation}
with
\[
A_{j,d,\theta}:=e^{i\theta}\,\mathcal{C}\,M_{\gamma_{j,d}}^{-1}\,L_{j,d}^{\C},\qquad M_{\gamma_{j,d}}=\begin{pmatrix}
    |\gamma_{j,d}(t)|^2  & 0\\
    0 & 1
\end{pmatrix}.\]
Here $\mathcal C$ denotes coefficientwise complex conjugation in the frame $(T_{j,d},N)$. In particular, for $d>0$
\[
\ker A_{j,d,\theta}=\ker L_{j,d}^{\C}=0.
\]
\end{prop}

\begin{proof}
Set
\[
V(z):=P(\eta\bar{z}_{tt})=\eta\bar{z}_{tt}-\lambda(z)\eta\bar{z},\qquad \lambda(z):=\frac{z^T\bar{z}_{tt}}{|z|^2}.\]
To simplify notations, we also denote
\[\gamma(t)=\gamma_{j,d}(t),\quad~ c=c_{j,d},\quad~ T=T_{j,d}(t),\quad~ \xi=\alpha T+\beta N.\]

Linearizing \eqref{eq:flow} at $\gamma=\gamma_{j,d}$ along $\xi$, we obtain
\[DV_{\gamma}(\xi)=\eta\,\bar{\xi}_{tt}-D\lambda_{\gamma}(\xi)\eta\,\bar\gamma-\bar{c}^2\,\eta\,\bar{\xi}.\]
Using $\gamma_{tt}=c^2\gamma$, we compute $D\lambda_{\gamma}(\xi)$ as follows:
\begin{align*}
    D\lambda_{\gamma}(\xi)=~&\frac{\gamma^T\bar{\xi}_{tt}+\bar{c}^2\,\xi^T \bar \gamma}{|\gamma|^2}-\frac{\bar c^2(\xi^T \bar \gamma+ \gamma^T \bar \xi)}{|\gamma|^2}=~\frac{\gamma^T\bar{\xi}_{tt}-\bar{c}^2\,\gamma^T\,\bar{\xi}}{|\gamma|^2}\\
    =~&\frac{\gamma^T(\bar{\alpha}_{tt}+\bar{c}^2\bar{\alpha})\bar{T}+2\gamma^T\bar{c}\bar{\alpha}_t\bar{\gamma}+\gamma^T\bar{\beta}_{tt}N-\bar{c}^2\bar \alpha\,\gamma^T\,\bar{T}}{|\gamma|^2}\\
    =~&\frac{\gamma^T\bar{T}}{|\gamma|^2} \bar{\alpha}_{tt}+2\bar{c}\,\bar{\alpha}_t.
\end{align*}
Using the identity
\[\cosh^2(ct)-\sinh^2(ct)=1,\]
we can compute
\[\eta\bar{T}-\frac{\gamma^T\bar{T}}{|\gamma|^2} \eta \bar \gamma=\frac{1}{|\gamma|^2}(\sinh(ct),\cosh(ct),0)=\frac{1}{|\gamma|^2}T.\]
Substituting these results into $DV_{\gamma}(\xi)$, 
we obtain
\[DV_{\gamma}(\xi)=\frac{\bar{\alpha}_{tt}}{|\gamma|^2}T+\left(\bar{\beta}_{tt}-\bar{c}^2\bar{\beta}\right)N.\qedhere\]
\end{proof}

The following standard notion of exponential dichotomy for a linear evolution equation and the  Lyapunov--Perron stable-manifold theorem for semilinear evolution equations with an exponential dichotomy are useful in our later discussion. We state them as follows. For more details, see \cite{Henry1981, Lin1986, PSS1997, Schnaubelt1999}.


\begin{defn}\label{def:exponential-dichotomy}
Let $A:D(A)\subset H\to H$ be a $\R$-linear closed operator on a real
Hilbert space $H$.  We say that the linear equation
\[
 \partial_s\xi+A\xi=0
\]
has an exponential dichotomy with splitting
\[
 H=E^s\oplus E^u
\]
and gap $\mu_*>0$ if there exist bounded projections
\[
 \Pi^s:H\to E^s,\qquad \Pi^u:H\to E^u,
 \qquad \Pi^s+\Pi^u=\mathrm{Id},
\]
such that $E^s$ and $E^u$ are invariant for the linear flow, and the
following estimates hold:
\[
 \|e^{-sA}\Pi^s\|_{H\to H}\le C e^{-\mu_*s},
 \qquad s\ge0,
\]
and
\[
 \|e^{sA}\Pi^u\|_{H\to H}\le C e^{-\mu_*s},
 \qquad s\ge0.
\]
Here $e^{-sA}\Pi^s$ denotes the forward evolution on the stable subspace,
whereas $e^{sA}\Pi^u$ denotes the backward evolution on the unstable
subspace.
\end{defn}


\begin{lem}[\cite{Henry1981,PSS1997}]\label{lem:Lyapunov--Perron}Let $A$ be a $\R$-linear closed operator on a Hilbert space $H$, with
domain continuously embedded in a stronger space $H^1$, and assume that
the linear equation
\[
 \partial_s\xi+A\xi=0
\]
admits an exponential dichotomy as in Definition \ref{def:exponential-dichotomy}. Consider
\begin{equation}\label{eq:abstract-hyperbolic-end}
 \partial_s\xi+A\xi=N(\xi),
 \qquad
 N(0)=0,\qquad DN(0)=0,
\end{equation}
where $N$ is $C^1$ from a small $H^1$-ball to $H$, and satisfies
\[
 \|N(\xi_1)-N(\xi_2)\|_{H}
 \le
 C_N\bigl(\|\xi_1\|_{H^1}+\|\xi_2\|_{H^1}\bigr)
 \|\xi_1-\xi_2\|_{H^1}.
\]
Then every solution of \eqref{eq:abstract-hyperbolic-end} satisfying
\[
 \lim_{s\to+\infty}\|\xi(s)\|_{H^1}=0
\]
actually converges exponentially: for every $0<\mu<\mu_*$, after increasing
$s_0$, there exists $C_\mu>0$ such that
\[
 \|\xi(s)\|_{H^1}\le C_\mu e^{-\mu(s-s_0)}
 \qquad (s\ge s_0).
\]  
Moreover the small initial data of all such decaying
half-trajectories form a $C^1$ graph over $E^s$, tangent to $E^s$ at
the origin.
\end{lem}


\subsection{Lefschetz thimble for generic $\theta$} We now describe Lefschetz thimbles of the energy functional on $\mathcal{P}_d^{\C}$ for generic phase $\theta$ to prepare for the Picard–Lefschetz analysis. 

\begin{defn}\label{def:thimble}
Fix a critical path $\gamma_{j,d}$ and a phase $\theta$. Let $\mathcal M_j^{\theta,+}$ be the set of maps
\[
z:[0,+\infty)\longrightarrow \PP_d^{\C}
\]
such that, in any local chart along the image,
\[
z\in C^0([0,+\infty);H^1_t)\cap C^1((0,+\infty);L^2_t)\cap C^0((0,+\infty);H^2_t\cap H^1_{0,t}),
\]
and $z$ satisfies
\begin{equation}
\begin{cases}
\partial_s z=P_z\bigl(e^{i\theta}\eta\,\bar{z}_{tt}\bigr)
\quad\text{ in\, }L^2([0,\tau])~\text{ for every }s>0,
\\[0.1cm]
z(s,0)=a,
\qquad
z(s,\tau)=b_d,
\\[0.1cm]
\lim\limits_{s\to+\infty}\Norm{z(s,\cdot)-\gamma_{j,d}(\cdot)}_{H^1([0,\tau])}=0,
\\[0.1cm]
\int_0^{+\infty}\norm{\partial_s z(s,\cdot)}_{L^2_t}^2\,\dd s<\infty.
\end{cases}
\end{equation}
The Lefschetz thimble attached to $\gamma_{j,d}$ is
\begin{equation*}\label{eq:Jjdef}
\mathcal{J}_j^{\theta}:=\ev(\mathcal M_j^{\theta,+})
=\{z(0,\cdot): z\in \mathcal M_j^+\}\subset \PP_d^{\C}.
\end{equation*}
\end{defn}

\begin{rmk}
One can also formulate the half-cylinder problem first in a weak class with
\[
z\in L^\infty_{\mathrm{loc}}([0,+\infty);H^1_t)\cap H^1_{\mathrm{loc}}([0,+\infty);L^2_t).
\]
The equation is tested against compactly supported tangent test fields, together with the same endpoint condition, finite-energy condition, and strong $H^1_t$ asymptotic condition. Whenever the image stays in a compact subset of $\HH_{\C}^2$, the local coordinate systems used later in the no-escape analysis provide the required bootstrap on every finite strip, so such a weak solution upgrades to the strong class above. Therefore, in the following we will work directly with strong solutions and strong $H^1_t$ asymptotics.
\end{rmk}

From Definition \ref{def:thimble}, we have
\begin{itemize}
    \item For every $z_0\in\mathcal{J}_j^{\theta}$,
    \[G_{\theta}(z_0)=G_{\theta}(\gamma_{j,d}),\qquad F_{\theta}(z_0)\ge F_{\theta}(\gamma_{j,d});\]
    \item We impose the genericity assumption that the twisted phase separates the imaginary parts of the critical values:
    \[G_{\theta}(\gamma_{j,d})\neq G_{\theta}(\gamma_{k,d}),\qquad\text{whenever}~j\ne k.\]
    With this genericity assumption, for any $j\ne k$, there is no finite-energy solution
    \[z:\R\times[0,\tau]\rightarrow\HH_{\C}^2\]
    of \eqref{eq:flow} such that 
    \[\lim_{s\rightarrow-\infty}\Norm{z(s,\cdot)-\gamma_{j,d}}_{H^1([0,\tau])}=0,
    \qquad
    \lim_{s\rightarrow+\infty}\Norm{z(s,\cdot)-\gamma_{k,d}}_{H^1([0,\tau])}=0.\]
    It follows that the thimbles $\mathcal{J}_j^{\theta}$ are pairwise disjoint, and each critical path determines exactly one thimble.
\end{itemize}



\begin{prop}\label{prop:expconv}
Let $d>0$, and let $z\in\mathcal M_j^{\theta,+}$. Then there exist constants $C,\mu>0$ such that
\[
\norm{z(s,\cdot)-\gamma_{j,d}}_{H^1_t}
\le Ce^{-\mu s}
\qquad s\gg0.
\]
Moreover, one can take $\mu$ smaller than
\[
\mu_{j,d,*}:=\min\left\{
\inf_n \phi^{\tan}_{n,j,d},\,
\inf_n |\phi^{\mathrm{nor}}_{n,j,d}|\right\},
\]
where $\phi^{\tan}_{n,j,d}$ are the positive eigenvalues of
\[
-\alpha_{tt}=\phi\,|\gamma_{j,d}(t)|^2\alpha,
\qquad \alpha(0)=\alpha(\tau)=0,
\]
and
\[
\phi^{\mathrm{nor}}_{n,j,d}=
\left(\frac{n\pi}{\tau}\right)^2+c_{j,d}^2 .
\]
\end{prop}



\begin{proof}
Choose exponential coordinates near $\gamma_{j,d}$ and write
\[
z(s,\cdot)=\exp_{\gamma_{j,d}}^{\C}\xi(s,\cdot).
\]
Since $z(s,\cdot)\to\gamma_{j,d}$ in $H^1_t$, for all sufficiently large
$s$ the trajectory lies in this chart. In these coordinates the equation
has the form
\[
\partial_s\xi+A_{j,d,\theta}\xi=N_{\theta}(\xi),\qquad
N_{\theta}(0)=0,\qquad DN_{\theta}(0)=0,
\]
where $A_{j,d,\theta}$ is from Proposition \ref{prop:linearization}, and $N_{\theta}$ satisfies the standard quadratic estimate in $H^1_t$.

Set
\[B_{j,d}=M_{\gamma_{j,d}}^{-1}L^{\C}_{j,d},\qquad B_{j,d}^{\#}=\mathcal{C}B_{j,d}\mathcal{C}.\]
Then as $\R$-linear operators,
\[A_{j,d,\theta}^2=B_{j,d}^{\#}B_{j,d},\]
and in particular, the real spectrum of $A_{j,d,\theta}$ is obtained from the complex spectrum of $B_{j,d}$ as follows:
\[\phi\in \sigma_{\C}(B_{j,d})~~\Rightarrow~~ \pm|\phi|\in \sigma_{\R}(A_{j,d,\theta}).
\]

\begin{itemize}
    \item The tangential eigenvalues of $B_{j,d}$ are the eigenvalues
$\phi^{\mathrm{tan}}_{n,j,d}$ of the weighted Dirichlet problem
\[
-\alpha_{tt}(t)=\phi^{\mathrm{tan}}_{n,j,d}
|\gamma_{j,d}(t)|^2\alpha(t)=\phi_{n,j,d}^{\tan}\cosh\left(\frac{2dt}{\tan}\right)\alpha(t),
\qquad\alpha(0)=\alpha(\tau)=0.
\]
They are independent of $j$ and positive. In particular, the minimal eigenvalue satisfies

\[\phi_{1,d}^{\tan}:=\phi_{1,j,d}^{\tau}\ge\frac{1}{\cosh(2d)}\left(\frac{\pi}{\tau}\right)^2>0.\]
\item The normal eigenvalues are
\[
\phi^{\mathrm{nor}}_{n,j,d}=
\left(\frac{n\pi}{\tau}\right)^2+c_{j,d}^2,
\qquad n\ge1.
\]
In particular, $|\phi^{\mathrm{nor}}_{n,j,d}|\rightarrow+\infty$ as $n\rightarrow\infty$. By Proposition \ref{prop:nondegenerate}, none of the normal eigenvalues vanishes when $d>0$. 
\end{itemize}

Then
\[
\operatorname{dist}\bigl(0,\sigma_{\R}(A_{j,d,\theta})\bigr)\ge \mu_{j,d,*}:=\min\left\{\inf_{n\ge1}\phi^{\mathrm{tan}}_{n,j,d},
\inf_{n\ge1}\left|\phi^{\mathrm{nor}}_{n,j,d}\right|\right\}>0.
\]
Thus $A_{j,d,\theta}$, regarded as an operator on the underlying real
Hilbert space, has an exponential dichotomy with gap bounded from
below by any number smaller than $\mu_{j,d,*}$.

By Lemma \ref{lem:Lyapunov--Perron}, for every
$0<\mu<\mu_{j,d,*}$, after increasing $s_0$ there exists
$C_\mu>0$ such that
\[
\|\xi(s,\cdot)\|_{H^1_t}\le C_\mu e^{-\mu(s-s_0)}\qquad(s\ge s_0).
\]
Finally, the equation and elliptic estimates for the $\R$-linear elliptic
operator $A_{j,d,\theta}$ bootstrap the decay to higher Sobolev norms on the end. Since the
exponential chart is smooth and has identity differential at the origin, the
same $H^1_t$-estimate holds for
$z(s,\cdot)-\gamma_{j,d}$. 
\end{proof}

Immediately, we have the local graph result for the thimble germ.

\begin{cor}\label{cor:localthimble}
Let $d>0$.  There exists a neighborhood $\mathcal{U}$ of $\gamma_{j,d}$ in
$\PP^{\C}_d$, a real Hilbert splitting
\[
T_{\gamma_{j,d}}\PP^{\C}_d
=E^s_{j,d,\theta}\oplus E^u_{j,d,\theta},
\]
where $E^s_{j,d,\theta}$ (or $E^u_{j,d,\theta}$ respectively) is the Hilbert direct sum of the decaying (or growing respectively) real
lines of the scalar mode equations, and a $C^1$ map
\[
\Psi_{j,d,\theta}:B_\rho(E^s_{j,d,\theta})\longrightarrow E^u_{j,d,\theta},
\qquad\Psi_{j,d,\theta}(0)=0,\qquad D\Psi_{j,d,\theta}(0)=0,
\]
such that
\[
\mathcal{J}_j^{\theta}\cap \mathcal{U}=
\left\{\exp_{\gamma_{j,d}}^{\C}\bigl(\xi^s+\Psi_{j,d,\theta}(\xi^s)\bigr):\xi^s\in B_\rho(E^s_{j,d,\theta})\right\}.
\]
where $B_\rho(E^s_{j,d,\theta})$ denotes the open ball of radius $\rho$
centered at the origin in $E^s_{j,d,\theta}$ with respect to the Hilbert norm. In particular,
\[
T_{\gamma_{j,d}}\mathcal{J}_j^{\theta}=E^s_{j,d,\theta}.
\]
\end{cor} 



\begin{rmk} We can regard this local graph description as the infinite-dimensional analogue
of the holomorphic Morse lemma. In the path-space setting, the complex Hessian is replaced by the
Jacobi operator $L^{\C}_{j,d}$, and the $\R$-linear linearized flow is governed
by
\[
A_{j,d,\theta}=e^{i\theta(d)}\,\mathcal{C}\,M_{\gamma_{j,d}}^{-1}\,L^{\C}_{j,d}.
\]
Its spectral decomposition plays the role of
the diagonal quadratic normal form. 
\end{rmk}


\subsection{Real path space as a thimble}
In this subsection, we study the Lefschetz thimble attached to the real geodesic $\gamma_{0,d}$ for the  $\theta=0$ flow, that is,
\begin{equation}\label{eq:flow0}
\partial_s z=-\nabla F_0(z)=P_z(\eta\bar{z}_{tt}). 
\end{equation}
For simplicity, we will write $\mathcal J_0$ instead of $\mathcal J_0^0$ and prove (Theorem \ref{thm:global-real-thimble})
\[\mathcal{J}_0:=\{z(0)\mid z\in\mathcal{M}_0^{0,+}\}=\PP_d^{\R}\subset\PP_d^{\C}.\]
Thus the starting real path space $\PP_d^{\R}$ for the heat kernel can be viewed as one Lefschetz thimble on $\PP_d^{\C}$. It is interesting to see whether other Lefschetz thimbles $\mathcal{J}_j$ have heat type interpretations. 

\subsubsection{The real slice and the real thimble}
Recall that $a,b_d\in\HH_{\R}^2$, and we can work on the real path space with endpoints $a,b_d$, that is,
\[\PP_d^{\R}:=\PP_{a,b_d}(\HH_{\R}^2)=\{x\in H^1([0,\tau],\HH_{\R}^2)\mid x(0)=a,x(\tau)=b_d\}.\]

\begin{prop}\label{prop:real-invariant}
At $\theta=0$, for every real path $x\in \PP^{\R}_d$, we have
\[
P_x(\eta x_{tt})\in T_x \HH^2_{\R}
\]
pointwise in $t$.  The restriction of the flow equation \eqref{eq:flow} at $\theta=0$ to
$\PP^{\R}_d$ is the endpoint-fixed heat flow on
$\HH^2_{\R}$.
\end{prop}

\begin{proof}
Let $x\in \PP^{\R}_d$.  Then the vector
$\eta x_{tt}$ is real.  Moreover, in the projection formula
\[
P_x(\eta x_{tt})=\eta x_{tt}-\frac{x^T x_{tt}}{|x|^2}\eta x,
\]
all coefficients are real.  Hence $P_x(\eta x_{tt})$ is a real tangent
vector to $\HH^2_{\R}$.  Therefore the full
$\theta=0$ equation \eqref{eq:flow0} restricts to an evolution equation on the real slice
$\PP^{\R}_d$. 

On $\PR$ the functional $S$ is real, and $F_0$ is the real energy
\[
F_0(x):=\frac12\int_0^\tau g_{\HH_{\R}^2}(x_t,x_t)\,\dd t.
\]
Therefore, the restricted equation is the downward $L^2$-gradient flow of $F_0$. Here the $L^2$ metric on $\PP^{\R}_d$ is induced from the standard metric on $\R^3$ restricting to $\HH_{\R}^2$. 
\end{proof}


To understand this real flow, let us recall the following necessary existence and compactness theorem.


\begin{thm}[{\cite{ES,Ham}}]
\label{thm:eells-sampson-closed}
Let $(M,g)$ be a compact Riemannian manifold possibly with boundary and let $(N,h)$ be a compact
Riemannian manifold with non-positive sectional curvature. If $\partial M\neq\varnothing$, fix a
smooth boundary value $\varphi:\partial M\to N$ and an initial map
$u_0\in C^\infty(M,N)$ with $u_0|_{\partial M}=\varphi$.

Consider the Dirichlet harmonic-map heat flow
\[
\partial_s u=\operatorname{tr}_g\nabla du,
\qquad u(s,\cdot)|_{\partial M}=\varphi,\qquad 
u(0,\cdot)=u_0.
\]
Then the following hold.

\begin{enumerate}
\item[\textup{(1)}]
\textup{(Eells--Sampson--Hamilton theorem)}
The heat flow has a unique smooth solution
\[
u:[0,+\infty)\times M\to N .
\]
Moreover, the energy is non-increasing along the flow.

\item[\textup{(2)}]
\textup{(Eells--Sampson compactness theorem)} There exist a sequence $s_j\to+\infty$ and a smooth harmonic map
$u_\infty:M\to N$ such that, after passing to a subsequence,
\[
u(s_j,\cdot)\longrightarrow u_\infty
\qquad\text{strongly in }H^1(M,N).
\]
Moreover, the convergence is smooth on $M$. 
\end{enumerate}
\end{thm}


We are going to apply Eells-Sampson-Hamilton Theorem \ref{thm:eells-sampson-closed} to $M=[0,\tau]$ and $N=\HH_{\R}^2$ to describe the thimble in the real slice. Even though Theorem \ref{thm:eells-sampson-closed} is stated for compact targets, our trajectory has a no-escape bound:
\[
d_{H^2_{\R}}(x(s,t),a)\le
\int_0^\tau |x_t(s,r)|\,\dd r\le
\sqrt{\tau}\|x_t(s,\cdot)\|_{L^2_t}\le
\sqrt{2\tau F(x_0)}.
\]
Our result is
 
\begin{prop}\label{prop:real-thimble}
At $\theta=0$, we have
\[\mathcal{J}_{0}^{\R}:=\mathcal{J}_0\,\cap\,\mathcal{P}_d^{\R}=\mathcal{P}_d^{\R}.\]
\end{prop}

\begin{proof}
Fix arbitrary initial path $x_0\in \PP^{\R}_d$. 
By Proposition~\ref{prop:real-invariant}, the real heat-flow trajectory is also a solution of the full $\theta=0$ equation, viewed inside $\PP^{\C}_d$.

Now, $\HH^2_{\R}$ is complete, simply connected and has non-positive
sectional curvature, by Theorem \ref{thm:eells-sampson-closed}, the endpoint-fixed heat flow exists smoothly for all $s\ge0$, and there exists a sequence $s_n\to+\infty$ such that $x(s_n,\cdot)$ converges,
after passing to a subsequence, in $H^1$ and smoothly on compact subintervals
of $(0,\tau)$, to a harmonic map $x_\infty$.  In our one-dimensional
source case, 
$x_\infty$ is a geodesic in $\HH^2_{\R}$ from $a$ to $b_d$. Since $\HH^2_{\R}$ is a Hadamard manifold, by Cartan--Hadamard theorem, the geodesic with fixed
endpoints $a$ and $b_d$ is unique. Hence
\[
 x_\infty=\gamma_{0,d}\in \mathcal{J}^{\R}_0.\qedhere
\]
%Finally the estimate in Proposition~\ref{prop:hessian-positive} at $\gamma_{0,d}$ upgrades
%subsequential convergence to exponential convergence of the forward tail. Thus the real heat-flow trajectory starting from $x_0$ converges to
%$\gamma_{0,d}$.  Hence $x_0\in \mathcal{J}^{\R}_0$. 
\end{proof}



\subsubsection{Conjugation symmetry and the local stable manifold}




Now the linearization of the full $\theta=0$ flow  \eqref{eq:linflow} at $\gamma_{0,d}$, becomes
\[\partial_s\xi+B_d\bar{\xi}=0,\qquad B_d:=M_{\gamma_{0,d}}^{-1}L_{0,d}^{\C}.\]
In the real frame
\[T_{0,d}(t)=\left(\sinh\frac{dt}{\tau},\cosh\frac{dt}{\tau},0\right),
\qquad
N=(0,0,1),
\]
if $\xi=\alpha T_{0,d}+\beta N$, then
\[
B_d(\alpha,\beta)
=
\left(
-\frac{1}{|\gamma_{0,d}(t)|^2}\alpha_{tt},
\,
-\beta_{tt}+\frac{d^2}{\tau^2}\beta
\right).
\]


Furthermore write $\xi=u+iv$, with $u,v\in T_{\gamma_{0,d}}\PP_d^{\R}$, then
\[u_s+i v_s=-B_d u+iB_dv.\]
Since $B_d$ is a differential operator with real-valued coefficients, taking real and imaginary parts yields
\begin{equation}\label{eq:uvlineaer}
u_s=-B_
du,\qquad v_s=B_dv.
\end{equation}
In particular, the forward stable and unstable spaces are give by
\[E^s=T_{\gamma_{0,d}}\PP_d^{\R}, \qquad E^u=iT_{\gamma_{0,d}}\PP_d^{\R}.\]


Applying Corollary \ref{cor:localthimble} to the Lefchetz $\mathcal{J}_0$, i.e. to the $\theta=0$ full flow near $\gamma_{0,d}$, there exists a neighborhood
$\mathcal{U}\subset \PP^{\C}_d$ of $\gamma_{0,d}$ such that the set of points whose forward
flow remains in $\mathcal{U}$ and converges to $\gamma_{0,d}$ is a $C^1$ local stable
manifold
\[
\mathcal{J}_0\cap \mathcal{U}=
\left\{
\exp_{\gamma_{0,d}}^{\C}\bigl(u+f(u)\bigr):u\in B_\rho(T_{\gamma_{0,d}}\PP_d^{\R})\right\},
\]
where
\[
f:B_\rho(T_{\gamma_{0,d}}\PP_d^{\R})\longrightarrow iT_{\gamma_{0,d}}\PP_d^{\R},
\qquad
f(0)=0,\qquad Df(0)=0.
\]



\medskip
Define the anti-holomorphic involution
\[
\tau:\PP_d^{\C}\to\PP_d^{\C},
\qquad
\tau(z)=\overline z.
\]
Its fixed-point set is exactly $\PR$.
Because $F_0(\overline z)=F_0(z)$, the vector field $X(z):=-\nabla F_0(z)$ satisfies
\begin{equation}\label{eq:tau-equiv}
X(\tau z)=D\tau\,X(z).
\end{equation}
Thus the full flow is $\tau$-equivariant. With the above analysis, we have the following local result.

\begin{prop}\label{prop:equivariant-local-stable}
There exists a neighborhood $\mathcal U$ of $\gamma_{0,d}$ in $\PP_{d}^{\C}$ such that
\[
\mathcal{J}_0\cap\mathcal U=\PR\cap\mathcal U.
\]
Equivalently, the local stable manifold of $\gamma_{0,d}$ for the full complexified flow is exactly the real slice. Furthermore, if $z\in\mathcal{M}_0^+$, then there exists $S\ge 0$ such that
\[
z(s)\in\PR\qquad\forall~ s\ge S.
\]
\end{prop}

\begin{proof}
By the local stable-manifold statement as above, the local
stable set near $\gamma_{0,d}$ is a $C^1$ graph over the stable space:
\[\mathcal{J}_0\cap\mathcal{U}
=\{(u,v)\mid v=f(u)\},
\]
where the coordinates are taken with respect to
\[
T_{\gamma_{0,d}}\PP^{\C}_d
=E^s_{0,d,0}\oplus E^u_{0,d,0}
=T_{\gamma_{0,d}}\PP^{\R}_d\oplus iT_{\gamma_{0,d}}\PP^{\R}_d,
\]
and
\[
f(0)=0,\qquad Df(0)=0.
\]
Now $\tau$ acts in these coordinates by
\[
\tau(u,v)=(u,-v).
\]
The local stable manifold is unique among $C^1$ invariant graphs tangent to $E^s$, and the flow is $\tau$-equivariant by \eqref{eq:tau-equiv}. Therefore $\mathcal{J}_0\cap\mathcal{U}$ is $\tau$-invariant. If $(u,f(u))$ lies on the graph, then so does
\[
\tau(u,f(u))=(u,-f(u)).
\]
Since the graph is single-valued over $u$, it follows that $f(u)=-f(u)$, hence $f\equiv 0$.
Therefore
\[
\mathcal{J}_0\cap\mathcal{U}=\PR\cap\mathcal U. \qedhere
\]
\end{proof}

The following standard backward uniqueness will help us analyze the flow for $0\le s\le S$.

\begin{prop}[{\cite{LM1960}}]
\label{prop:backward-uniqueness-parabolic-system}
Let
\[w : [s_0,S]\times[0,1]\longrightarrow \R^N
\]
be a sufficiently regular solution of the uniformly parabolic linear system
\[\partial_s w-A(s,t)\partial_t^2 w
=B(s,t)\partial_t w+C(s,t)w
\]
on the compact cylinder $[s_0,S]\times[0,1]$, where the coefficients are bounded and $A(s,t)$ is uniformly positive definite. Assume that $w$ satisfies homogeneous boundary conditions at $t=0,1$. If
\[w(S,\cdot)=0,\]
then
\[w\equiv 0,\quad \text{on }~ [s_0,S]\times[0,1].\]
\end{prop}



Now for any $z\in \mathcal{M}_0^+$, since the full $\theta=0$ flow \eqref{eq:flow0} is $\tau$-equivariant,  $\tau(z)$ is also a solution of the full system. We apply the backward uniqueness result Proposition  \ref{prop:backward-uniqueness-parabolic-system} for parabolic systems 
to the difference $w=z-\tau(z)$, regarded as a real-valued system. Combined with Proposition \ref{prop:equivariant-local-stable}, we have $w\equiv0$, that is, taking $s_0=0$,
\[\tau(z)(s)=z(s),\quad \forall\,s\in[0,+\infty).\]
In particular, $z(0)\in\PP_d^{\R}$. Then we arrive at the following theorem.

\begin{thm}\label{thm:global-real-thimble}
At $\theta=0$, we have
\[
\mathcal{J}_0=\PP_d^{\R}.
\]
Equivalently, the stable Lefschetz thimble attached to the real geodesic
$\gamma_{0,d}$ consists precisely of all real paths. 
\end{thm}




\subsection{Prediction on two connecting trajectories}
We now start to consider the prediction (Conjecture \ref{conj:alien}) on the connecting trajectories from resurgence analysis. By Proposition \ref{prop:H2-pointed-alien-stokes-constant}, the Alien coefficient $2$ predicts that there should be exactly two solutions of downward Morse flow connecting $\gamma_{k,d}$ to $\gamma_{0,d}$ ($k\neq 0$) at the Stokes phase. We pick $k=1$ and verify this prediction for $d$ sufficiently small. This provides a nontrivial test of Conjecture \ref{conj:alien}.


The Stokes phase $\theta(d)$ for flows from $\gamma_{1,d}$ to $\gamma_{0,d}$ satisfies 
\[G_{\theta(d)}(\gamma_{1,d})=G_{\theta(d)}(\gamma_{0,d}),\qquad F_{\theta(d)}(\gamma_{1,d})\ge F_{\theta(d)}(\gamma_{0,d}).\]
Recall that
\[\gamma_{j,d}=\left(\cosh \frac{(d+2\pi ij)t}{\tau},\,\sinh \frac{(d+2\pi ij)t}{\tau},\,0 \right), \quad j=0,1.\]
Thus the above imaginary and real part constraints solve $\theta(d)$ to be
\begin{equation}\label{eq:thetad}
\theta(d)=\pi-\arctan\frac{d}{\pi}.    
\end{equation}
Then we will study the possible morse flow with the asymptotic conditions $\gamma_{1,d}$ and $\gamma_{0,d}$. In the following, we will call the case $j=0$ the lower end and the case $j=1$ the upper end. 


\vskip 0.2cm

We write
\[
\mathcal{M}_d(a,b_d;\gamma_{0,d},\gamma_{1,d})
\]
for the moduli space of finite-energy classical solutions $z:\R\times[0,\tau]\to \HH_{\C}^2$ of \eqref{eq:flow} such that
\begin{align*}
\Norm{z(s,\cdot)-\gamma_{1,d}}_{H^1([0,\tau])}\to 0\quad\text{as }s\to -\infty,
\\
\Norm{z(s,\cdot)-\gamma_{0,d}}_{H^1([0,\tau])}\to 0\quad\text{as }s\to +\infty,
\end{align*}
modulo translation $s\mapsto s+\sigma$. We are going to prove
\begin{itemize}
    \item for $d=0$, (Theorem \ref{thm:exact-two-d0})
    \[
\#\mathcal{M}_0(a,a;\gamma_{0,0},\gamma_{1,0})=2.
\]
   \item for $0<d\ll1$, (Theorem \ref{thm:exact-two-small-d})
   \[
\#\mathcal{M}_d(a,b_d;\gamma_{0,d},\gamma_{1,d})=2.
\]
\end{itemize}
The small-$d$ Theorem \ref{thm:exact-two-small-d} is obtained by continuation from the two distinguished $d=0$ trajectories together with a no-escape/compactness argument.


\subsection{Case: $d=0$}\label{sec:case d=0}
{
In this subsection we prove the $d=0$ count
\[
\#\mathcal M_0(a,a;\gamma_{0,0},\gamma_{1,0})=2.
\]
The argument proceeds in two steps. First, we study the equation on a flow-invariant slice
$\Sigma \subset \mathbb H^2_{\mathbb C}$, which is isomorphic to $S^2$, and obtain exactly two solutions there. Second, we regard the same equation as a problem in the full space $\mathbb H^2_{\mathbb C}$, which we also call it the $d=0$ full system. By combining tail rigidity with backward uniqueness, we show that every full-space solution is forced to lie in the slice. Therefore, the full $d=0$ system has no solutions other than the two slice solutions. 

Finally we also analyze the $d=0$ regularity problem and construct an invertible argumented linear operator which serves for $d$-perturbation in the next subsection.}
 
\subsubsection{Heat flow on $S^2$}\label{sec:heat flow on S2}
Let
\[
\rho(z_0,z_1,z_2):=(\overline z_0,-\overline z_1,-\overline z_2)
\]
be the anti-holomorphic involution of $\HH^2_{\C}$. Let $S^2$ be the unit sphere in $\R^3$. Then we can identify 
\[
\Sigma:=\operatorname{Fix}(\rho)=Q(S^2),\qquad Q:=\operatorname{diag}(1,-i,-i).
\]

\begin{prop}
At $d=0$ and $\theta(0)=\pi$, the Morse vector field
\[
P_z(e^{i\pi}\eta \overline{z}_{tt})
=P_z(-\eta\overline{z}_{tt})
\]
is $\rho$-equivariant. The map
\[
Q:\Omega_aS^2\longrightarrow \mathcal P^{\C}_0\cap \Omega_a\Sigma,
\qquad x\longmapsto Qx,
\]
identifies the restricted Morse functional $F_{\pi}(Qx)$ with the based-loop energy
\[
E_{S^2}(x)=\frac12\int_0^\tau |x_t|^2\,\dd t.
\]
In particular, the restriction of the $d=0,\theta=\pi$ flow to $\Sigma$ is the endpoint-fixed heat flow on $S^2$,
\[
x_s=\nabla_t x_t,\qquad x(s,0)=x(s,\tau)=a.
\]
\end{prop}

\begin{proof}
We check invariance of the slice. Denote $R=\diag(1,-1,-1)$.
%\[R=\begin{pmatrix}
%    1 & 0 & 0\\
%    0 &-1 & 0\\
%    0 & 0& -1
%\end{pmatrix}\]
Using the projection formula
\[
P_z(W)=W-\frac{z^T\eta W}{|z|^2}\eta\overline z,
\]
and $R^T\eta R=\eta$, $|R\bar{z}|^2=|z|^2$,  the tangent map $D\rho_z$ of $\rho_z$ satisfies
\[
D\rho_z(P_zW)=R\overline{P_zW}=P_{\rho z}(D\rho_zW).
\]
Taking $W=-\eta\overline z_{tt}$, we get
\[
D\rho_zP_{z}(-\eta\bar{z}_{tt})
=
P_{\rho z}\bigl(-\eta\,\overline{(\rho z)}_{tt}\bigr).
\]
Hence the vector field is $\rho$-equivariant. Therefore its fixed-point set $\Sigma$ is invariant under the flow.

For $z=Qx$, we have
\[
z_t^T\eta z_t
=
-\dot x_0^2+(-i\dot x_1)^2+(-i\dot x_2)^2
=
-|x_t|^2.
\]
Therefore for $\theta=\pi$,
\[
F_{\pi}(Qx)=\re(-S(Qx))
=-
\frac12\int_0^\tau z_t^T\eta z_t\,\dd t
=
\frac12\int_0^\tau |x_t|^2\,\dd t=E_{S^2}(x).
\]
Because $Q$ is unitary for the induced Hermitian metric, the downward gradient flow of $F_\pi$ on the slice $\Sigma$ corresponds to the downward $L^2$-gradient flow of $E_{S^2}$ on $\Omega_aS^2$. Thus the restricted equation is
\begin{equation*}
 x_s=\nabla_t x_t=x_{tt}+|x_t|^2x,
 \qquad x(s,0)=x(s,\tau)=a.
\end{equation*}
And \[
\gamma_{0,0}\leftrightarrow x_0(t)\equiv a,
\qquad
\gamma_{1,0}\leftrightarrow x_{1,*}(t)=\Bigl(\cos\frac{2\pi t}{\tau},-\sin\frac{2\pi t}{\tau},0\Bigr). \qedhere
\]
\end{proof}


We next formulate the slice problem in the natural parabolic class on $S^2$.


\begin{defn}
Let $\mathcal{X}^\Sigma$ be the set of maps
\[
x:\R\times[0,\tau]\longrightarrow S^2
\]
such that, in any local chart along the image,
\[
x\in C^0_{\mathrm{loc}}(\R;H^1([0,\tau]))
\cap C^1_{\mathrm{loc}}(\R;L^2([0,\tau]))
\cap C^0_{\mathrm{loc}}(\R;H^2([0,\tau])),
\]
and $x$ satisfies 
\begin{equation}\label{eq:sphere-heat}
\begin{cases}
 x_s=\nabla_t x_t=x_{tt}+|x_t|^2x \quad \text{in }~L^2([0,\tau]) \text{ for every }~s\in\R,\\[0.1cm]
x(s,0)=x(s,\tau)=a,
\\[0.1cm]
\int_0^{+\infty}\norm{\partial_s x(s,\cdot)}_{L^2_t}^2\,\dd s<\infty.  
\end{cases}    
\end{equation}
The slice moduli space $\mathcal{M}^\Sigma_0$ consists of those $x\in \mathcal{X}^\Sigma$ such that
\[
\|x(s,\cdot)-x_{1,*}\|_{H^1([0,\tau])}\longrightarrow 0
\qquad (s\to-\infty),
\]
and
\[
\|x(s,\cdot)-x_0\|_{H^1([0,\tau])}\longrightarrow 0
\qquad (s\to+\infty),
\]
modulo the $\R$-translation in $s$.
%\[x(s,t)\longmapsto x(s+\sigma,t),\qquad \sigma\in\R.\]
\end{defn}


Let $T(t)$ be the unit tangent field of $x_{1,*}$ along the great circle and let $N(t)$ be a unit normal field in $T S^2$ along $x_{1,*}$.
As before, every Dirichlet variation field can be written as
\[
\xi(t)=\alpha(t)T(t)+\beta(t)N(t),
\qquad
\alpha(0)=\alpha(\tau)=\beta(0)=\beta(\tau)=0.
\]

\begin{lem}\label{lem:hessian-once-round}
The Hessian of the based-path energy at $x_{1,*}$ is
\[
\mathrm{Hess}_{x_{1,*}}E_{S^2}(\xi,\xi)
=\frac12\int_0^\tau\Bigl(\Abs{\alpha_t}^2+\Abs{\beta_t}^2-\Bigl(\frac{2\pi}{\tau}\Bigr)^2\Abs{\beta}^2\Bigr)\,\dd t.
\]
Furthermore, write 
\[
\alpha(t)=\sum_{n\ge 1} a_n\sin\frac{n\pi t}{\tau},
\qquad
\beta(t)=\sum_{n\ge 1} b_n\sin\frac{n\pi t}{\tau},
\]
then
\[
\mathrm{Hess}_{x_{1,*}}E_{S^2}
=\frac{\pi^2}{4\tau}\sum_{n\ge 1}\bigl(n^2 a_n^2+(n^2-4)b_n^2\bigr).
\]
And we have
\begin{enumerate}[label=\textup{(\roman*)}]
\setlength{\itemsep}{-1pt}
    \item all tangential modes have positive Hessian eigenvalues;
    \item in the normal sector the $n=1$ mode $\sin\frac{\pi t}{\tau}N(t)$ has negative Hessian eigenvalue;
    \item the normal $n=2$ mode $\sin\frac{2\pi t}{\tau}\,N(t)$ has zero Hessian eigenvalue;
    \item all normal modes with $n\ge 3$ have positive Hessian eigenvalues.
\end{enumerate}
In particular $x_{1,*}$ is Morse--Bott with one unstable direction and one Bott direction.
\end{lem}




\begin{proof} The proof is similar to the computation in Propositions \ref{prop:nondegenerate}, \ref{prop:linearization}, or just follows from the standard second variation formula on $S^2$. 
\end{proof}



Let $\mathcal{C}_B$ denote the Bott family of once-round great circles based at $a$
which contains $x_{1,*}$. By Lemma \ref{lem:hessian-once-round}, the zero mode $\sigma_B$ of the Hessian at
$x_{1,*}$ is precisely the normal $n=2$ Jacobi field. Geometrically, %this is the infinitesimal motion inside $\mathcal C$. Thus
\[
T_{x_{1,*}}\mathcal{C}_B=\operatorname{Span}\{\sigma_B\}.
\]
If $E^u$ denotes the negative eigenspace and $E^s$ the positive eigenspace, then
\[
T_{x_{1,*}}\Omega_a S^2
=
E^u\oplus \operatorname{Span}\{\sigma_B\}\oplus E^s,
\qquad
\dim E^u=1.
\]
We fix a local codimension-one slice transverse to the Bott family, that is
\[
\mathcal S^-\subset \Omega_aS^2,\qquad
T_{x_{1,*}}\mathcal S^-=E^u\oplus E^s.
\]

Now in the transverse slice $\mathcal S^-$, the Hessian at $x_{1,*}$ has exactly one
negative eigenvalue and no kernel. Then we obtain
a one-dimensional local strong unstable manifold
\[
W^{uu}_{\mathrm{loc}}(x_{1,*};\mathcal S^-)
\]
tangent to $E^u$ at $x_{1,*}$. Since it is one-dimensional, removing the endpoint
$x_{1,*}$ leaves exactly two local connected components, which we denote by $x_+$ and $x_-$.



\begin{prop}\label{prop:S2-global-existence}
Each of the two local trajectory germs $x_+$ and $x_-$ extends uniquely to a smooth
finite-energy solution
\[
x_\pm:\R\times[0,\tau]\longrightarrow S^2
\]
of
\[
x_s=\nabla_t x_t,\qquad x(s,0)=x(s,\tau)=a,
\]
such that
\begin{align*}
\|x_\pm(s,\cdot)-x_{1,*}\|_{H^1([0,\tau])}&\longrightarrow 0
\qquad (s\to-\infty),\\
\|x_\pm(s,\cdot)-x_0\|_{H^1([0,\tau])}&\longrightarrow 0
\qquad (s\to+\infty).
\end{align*}
\end{prop}

\begin{proof}
We prove the statement for one of the two local branches; the other one is identical.

First consider the negative end $s\rightarrow-\infty$. Applying the parabolic unstable-manifold theorem to the transverse slice $\mathcal S^-$, the local strong unstable branch is represented by
a solution defined for $s\ll0$, and it converges exponentially to $x_{1,*}$: for some $\delta>0$,
\[
\|x(s,\cdot)-x_{1,*}\|_{H^1_t}\le Ce^{\delta s}
\qquad (s\to-\infty).
\]


Then consider the positive end $s\rightarrow+\infty$. Choose now a finite time $s_0$ on this local branch and use $x(s_0,\cdot)$ as initial
data for the 
flow
\[
x_s=\nabla_t x_t,\qquad x(s,0)=x(s,\tau)=a.
\]
Since the source $[0,\tau]$ and the target $S^2$ are compact, by the standard continuation criterion
for semilinear parabolic equations, the endpoint-fixed heat flow on $S^2$ has a unique smooth solution on $[s_0,+\infty)$. The extended solution has finite energy:
\[
\int_{s_0}^{+\infty}\|x_s(s,\cdot)\|_{L^2_t}^2\,\dd s=
E_{S^2}(x(s_0,\cdot))-\lim_{s\to+\infty}E_{S^2}(x(s,\cdot))\le E_{S^2}(x(s_0,\cdot)).
\]
Together with the exponential convergence on the negative end, this gives
\[
\int_{-\infty}^{+\infty}\|x_s(s,\cdot)\|_{L^2_t}^2\,\dd s<\infty.
\]

It remains to identify the forward limit and to prove strong $H^1$-convergence. Since
$E(x(s,\cdot))$ is decreasing and bounded below, the limit
\[
E_{S^2,\infty}:=\lim_{s\to+\infty}E_{S^2}(x(s,\cdot))
\]
exists. Also, finite energy implies that there is a sequence $s_j\to+\infty$ such that
\[
\|x_s(s_j,\cdot)\|_{L^2_t}\longrightarrow0.
\]
The uniform energy bound, the compactness of $S^2$, and parabolic estimates on the
shifted strips $[s_j-1,s_j+1]\times[0,\tau]$ imply, after passing to a subsequence, that
\[
x(s_j+\sigma,t)\longrightarrow x_\infty(t)
\]
smoothly on compact subsets of $(-1,1)\times[0,\tau]$, with the limit independent of
$\sigma$. Passing to the limit in the equation gives
\[
\nabla_t (x_\infty)_t=0,\qquad x_\infty(0)=x_\infty(\tau)=a.
\]
Therefore $x_\infty$ is a based geodesic loop on $S^2$. It follows that $x_{\infty}$ has length $2m\pi$ for some $m\in\Z_{\ge1}$ if $x_{\infty}$ is nonconstant. Thus,
\[
\frac12\int_0^\tau |(x_\infty)_t|^2\,\dd t
=\frac{(2m\pi)^2}{2\tau}
\ge
\frac{2\pi^2}{\tau}=E_{S^2}(x_{1,*}).
\]
On the other hand,  for every finite
$s_0$ on the unstable branch, we have 
\[
E_{S^2}(x(s_0,\cdot))<E_{S^2}(x_{1,*})\le E_{S^2,\infty}.
\]
Contradiction. Thus the limiting geodesic must be constant:
\[
x_\infty=x_0,\qquad x_0(t)\equiv a.
\]
Thus every subsequential forward limit is the constant loop $x_0$, and consequently
\[
E_\infty=E(x_0)=0.
\]

Finally, since 
\[
\|x_t(s,\cdot)\|_{L^2_t}^2=2E(x(s,\cdot))\longrightarrow0.
\]
and 
\[
x(s,t)-x_0(t)=x(s,t)-a=\int_0^t x_t(s,r)\,dr.
\]
We have
\[
\|x(s,\cdot)-x_0\|_{L^2_t}
\le
C_\tau \|x_t(s,\cdot)\|_{L^2_t}
\longrightarrow0.
\]
Furthermore, we obtain
\[
\|x(s,\cdot)-x_0\|_{H^1_t}\longrightarrow0
\qquad (s\to+\infty). \qedhere
\]
\end{proof}


Thus, the slice moduli space $\mathcal{M}_0^{\Sigma}$ consists of exactly two elements, denoted by $[x_+]$, $[x_-]$. Equivalently, the $d=0$ flow on $\Sigma$ has exactly two finite-energy trajectories from $\gamma_{1,0}$ to $\gamma_{0,0}$, modulo $s$-translation, denoted by $[\nu_{0,+}]$, $[\nu_{0,-}]$.


\subsubsection{The full $d=0$ system}\label{subsec:full-d=0}
For the full system at $d=0$, we introduce real variables
\[
z_0=q_0+ip_0,
\qquad
z_1=p_1-iq_1,
\qquad
z_2=p_2-iq_2.
\]
Then $q=(q_0,q_1,q_2)\in \R^3$ and $p=(p_0,p_1,p_2)\in \R^3$, and the quadric equation becomes
\begin{equation}\label{eq:qp-constraints}
q\cdot q-p\cdot p=1,
\qquad
q\cdot p=0.
\end{equation}

At $d=0$, the fixed-endpoint-conditions are
\begin{equation}\label{eq:qp-boundary}
q(s,0)=q(s,\tau)=a,
\qquad
p(s,0)=p(s,\tau)=0.
\end{equation}
The slice $\Sigma=Q(S^2)$ is exactly given by $p\equiv 0$. At $\theta=\pi$, 
   \[F_{\pi}(q,p)=\frac{1}{2}\int_0^{\pi}(|q_t|^2-|p_t|^2)\,\dd t.\]




\begin{lem}\label{lem:qp-full-flow}
At $\theta=\pi$, the full downward $L^2$-gradient equation in $(q,p)$ becomes
\begin{equation}
\label{eq:qp-full-flow}
\begin{cases}
 q_s=q_{tt}+\lambda q+\mu p,\\
 p_s=-p_{tt}-\lambda p+\mu q,
\end{cases}
\end{equation}
where $\lambda,\mu$ are determined by
\begin{equation}
\label{eq:lambda-mu-pi}
\lambda M=q\cdot(q_s-q_{tt})- p\cdot(p_s+p_{tt}),
\qquad\mu M=p\cdot(q_s-q_{tt})+q\cdot(p_s+p_{tt}),
\end{equation}
with
\[
 M:=|q|^2+|p|^2\geq 1.
\]
Along classical trajectories,
\[
\frac{\dd}{\dd s}F_\pi(q(s,\cdot),p(s,\cdot))=
-\int_0^\tau \bigl(|q_s|^2+|p_s|^2\bigr)\,\dd t .
\]
\end{lem}

\begin{proof}
For a tangent variation $(\dot q,\dot p)$, the constraints give
\[
q\cdot \dot q-p\cdot \dot p=0,\qquad
q\cdot \dot p+p\cdot \dot q=0.
\]
Moreover, using the fixed endpoint condition,
\[
\dd F_\pi(q,p)(\dot q,\dot p)=\int_0^\tau (q_t\cdot \dot q_t-p_t\cdot \dot p_t)\,\dd t
=\int_0^\tau (-q_{tt}\cdot \dot q+p_{tt}\cdot \dot p)\,\dd t .
\]
Thus the ambient downward gradient is $(q_{tt},-p_{tt})$. Considering the normal directions
\[
(q,-p), \qquad (p,q),
\]
the full downward gradient of $(q,p)$ has the form
\[
(q_s,p_s)=(q_{tt},-p_{tt})+\lambda(q,-p)+\mu(p,q).
\]
Here $\lambda,\mu$ can be determined as follows. Along this flow, using $q\cdot p=0$, we have
\begin{align*}
    &q\cdot(q_s-q_{tt})-p\cdot(p_s+p_{tt})=q\cdot(\lambda q+\mu p)-p\cdot(-\lambda p+\mu q)=\lambda(|q|^2+|p|^2),\\
    &p\cdot(q_s-q_{tt})+q\cdot(p_s+p_{tt})=p\cdot(\lambda q+\mu p)+q\cdot(-\lambda p+\mu q)=\mu(|q|^2+|p|^2).
\end{align*}

Set
\[
 M=|q|^2+|p|^2\geq |q|^2-|p|^2=1.
\]
Taking the $s$-derivative of \eqref{eq:qp-constraints}, we have
\[
 q\cdot q_s-p\cdot p_s=0,
 \qquad
 q\cdot p_s+p\cdot q_s=0,
\]
Then we can rewrite \eqref{eq:lambda-mu-pi} as
\[
 \lambda M
 =
 -(q\cdot q_{tt}+p\cdot p_{tt}),
 \qquad
 \mu M
 =
 q\cdot p_{tt}-p\cdot q_{tt}.
\]

Finally, along a classical solution,
\begin{align*}
\frac{\dd}{\dd s}F_\pi(q(s,\cdot),p(s,\cdot))
=~&\int_0^\tau (-q_{tt}\cdot q_s+p_{tt}\cdot p_s)\,\dd t.\\
=~&-\int_0^\tau (|q_s|^2+|p_s|^2)\,\dd t+
\int_0^\tau\lambda(q\cdot q_s-p\cdot p_s)+\mu(p\cdot q_s+q\cdot p_s)\,\dd t\\
=~&-\int_0^\tau \bigl(|q_s|^2+|p_s|^2\bigr)\,\dd t. \qedhere
\end{align*}
\end{proof}



\begin{prop}\label{prop:morse-end-C0}
Let $(q,p)$ be a classical finite-energy solution of
\eqref{eq:qp-boundary}--\eqref{eq:qp-full-flow} on
$[S_*,+\infty)\times [0,\tau]$.  Assume that
\begin{equation}\label{eq:H1-positive-end}
\Norm{(q(s,\cdot),p(s,\cdot))-(a,0)}_{H^1([0,\tau])}\to 0
\qquad\text{as }s\to +\infty .
\end{equation}
Then for every integer $k\ge 0$, there exist constants
$C_k>0$, $\delta>0$, and $S_0\ge S_*$ such that
\[
\Norm{(q(s,\cdot),p(s,\cdot))-(a,0)}_{C^k([0,\tau])}
\le C_k e^{-\delta(s-S_0)}
\qquad (s\ge S_0).
\]
In particular,
\[
\Norm{\lambda(s,\cdot)}_{C^k([0,\tau])}
+
\Norm{\mu(s,\cdot)}_{C^k([0,\tau])}
\le C_k' e^{-\delta(s-S_0)}
\qquad (s\ge S_0).
\]
\end{prop}


\begin{proof}
By the strong $H^1$-asymptotics and the embedding
$H^1([0,\tau])\hookrightarrow C^0([0,\tau])$, after increasing $S_0$ we can
work in the fixed local chart near $(a,0)$ used above. Write
\[
 (q,p)=\Phi(v),\qquad v=(\xi,\zeta),\qquad \Phi(0)=(a,0).
\]
The boundary conditions give
\[
 \xi(s,0)=\xi(s,\tau)=0,
 \qquad
 \zeta(s,0)=\zeta(s,\tau)=0 .
\]

We first eliminate $\lambda$ and $\mu$. Lemma \ref{lem:qp-full-flow} says
\[
 \lambda
 =
 -\frac{q\cdot q_{tt}+p\cdot p_{tt}}{|q|^2+|p|^2},
 \qquad
 \mu
 =
 \frac{q\cdot p_{tt}-p\cdot q_{tt}}{|q|^2+|p|^2},\qquad
 |q|^2+|p|^2\geq 1 .
\]
Substituting these expressions into \eqref{eq:qp-boundary}--\eqref{eq:qp-full-flow} gives a closed equation
for $v$ of the form
\[
 v_s+A_\infty v=R(v),
\]
where
\[
 A_\infty(\xi,\zeta)=(-\xi_{tt},\zeta_{tt})
\]
with Dirichlet boundary conditions on $[0,\tau]$, and $R$ is the
smooth nonlinear remainder. In particular,
\[
 DR(0)=0 .
\]
Since $0\notin \operatorname{spec}(A_\infty)$, by  Lemma \ref{lem:Lyapunov--Perron}, there exists constants $C>0$, $\delta>0$.
\[\|v(s,\cdot)\|_{H^1([0,\tau])}<Ce^{-\delta(s-s_0)}\]
Applying the standard parabolic bootstrapping estimates on unit strips, the $H^1$-exponential decay upgrades to $H^m$-exponential decay for every $m\ge 1$. Therefore,
after increasing $S_0$ and the constants if necessary, we obtain
\[
\|v(s,\cdot)\|_{H^m([0,\tau])}
 \le C_m e^{-\delta(s-S_0)},
 \qquad s\ge S_0 .
\]
Since the chart is smooth and fixes the origin, this implies, by Sobolev embedding on $[0,\tau]$, that for every integer $k\geq 0$,
\[
 \|(q(s,\cdot),p(s,\cdot))-(a,0)\|_{C^k([0,\tau])}
 \leq
 C_k e^{-\delta(s-S_0)}
 \qquad
 (s\geq S_0).
\]

It remains only to translate this decay into decay of the multipliers. Using the
explicit formulas above and the $C^{k+2}$-decay of $(q,p)-(a,0)$, we have
\[
 \|\lambda(s,\cdot)\|_{C^k([0,\tau])}
 +
 \|\mu(s,\cdot)\|_{C^k([0,\tau])}
 \leq
 C'_k e^{-\delta(s-S_0)}
 \qquad
 (s\geq S_0).\qedhere
\]
\end{proof}



\begin{lem}\label{lem:forward-tail-rigidity}
Let $(q,p)$ be a classical finite-energy solution of \eqref{eq:qp-boundary}--\eqref{eq:qp-full-flow} on $[S_*,+\infty)\times [0,\tau]$, and assume that
\[
\Norm{(q(s,\cdot),p(s,\cdot))-(a,0)}_{H^1([0,\tau])}\to 0
\qquad\text{as }s\to +\infty.
\]
Then there exists $S_0\ge S_*$ such that
\[
p(s,t)\equiv 0
\qquad\text{for all }~s\ge S_0,
\quad 0\le t\le \tau.
\]
\end{lem}


\begin{proof}
By Proposition \ref{prop:morse-end-C0}, after increasing $S_0$ we have
\[
|\lambda(s,t)|\leq \frac{\pi^2}{2\tau^2}
\qquad
\text{for all }~s\geq S_0,\; t\in[0,\tau].
\]
Set
\[
 E_p(s):=\frac12\int_0^\tau |p(s,t)|^2\,\dd t .
\]
Using
\[
 p_s=-p_{tt}-\lambda p+\mu q,
\]
the boundary condition $p(s,0)=p(s,\tau)=0$, and the constraint $q\cdot p=0$,
we obtain
\[
\begin{aligned}
 E_p'(s)
 &=\int_0^\tau p\cdot p_s\,\dd t                                      =\int_0^\tau p\cdot(-p_{tt}-\lambda p+\mu q)\,\dd t                    \\
 &=\int_0^\tau |p_t|^2\,\dd t-\int_0^\tau \lambda |p|^2\,\dd t .
\end{aligned}
\]
Hence, by the Dirichlet Poincaré inequality,
\[
\begin{aligned}
E_p'(s)&\geq
\left(\frac{\pi^2}{\tau^2}-\frac{\pi^2}{2\tau^2}\right)
\int_0^\tau |p|^2\,\dd t=
\frac{\pi^2}{\tau^2} E_p(s).
\end{aligned}
\]
Thus, for $s\geq S_0$,
\[
 E_p'(s)\geq cE_p(s),
 \qquad
 c:=\frac{\pi^2}{\tau^2}>0.
\]
If $E_p(s_1)>0$ for some $s_1\geq S_0$, then Gronwall's inequality gives
\[
 E_p(s)\geq e^{c(s-s_1)}E_p(s_1)
 \qquad
 (s\geq s_1),
\]
which contradicts $E_p(s)\to0$ as $s\to+\infty$. Therefore
\[
 E_p(s)=0
 \qquad
 (s\geq S_0),
\]
and hence $p(s,t)\equiv0$ on the forward tail.
\end{proof}



Now tail rigidity gives $p\equiv 0$ on a half-cylinder. To conclude that the entire trajectory lies in the slice, we are going to use the backward-uniqueness result Proposition \ref{prop:backward-uniqueness-parabolic-system} again.

Lemma \ref{lem:forward-tail-rigidity} says that once a trajectory has entered the slice on a forward tail, we have $p\equiv 0$ there and $q$ solves the smooth sphere heat flow on the compact target $ S^2$. Expanding the full system \eqref{eq:qp-full-flow} around that slice trajectory and using the smooth dependence of $(\lambda,\mu)$ on $(q,p,q_t,p_t,q_{tt},p_{tt})$, we can rewrite the $p$-equation on the whole cylinder as
\[
p_s+p_{tt}=A_1(s,t)p_t+A_0(s,t)p.
\]
Therefore, we can apply Proposition \ref{prop:backward-uniqueness-parabolic-system} to $w=p$ and $[s_0,S_0]$ for any $s_0<S_0$, then let $s_0\rightarrow-\infty$. Then we obtain the following result.


\begin{prop}\label{prop:global-reduction-slice}
Let $z$ be a finite-energy classical $d=0$ trajectory from $\gamma_{1,0}$ to $\gamma_{0,0}$.
Then $z(s,t)\in \Sigma$ for all $(s,t)\in \R\times [0,\tau]$.
Hence the full ambient $d=0$ moduli space coincides with the slice moduli space.
\end{prop}


Combining Proposition \ref{prop:S2-global-existence} with Proposition \ref{prop:global-reduction-slice}, we obtain:

\begin{thm}\label{thm:exact-two-d0}
The full $d=0$ moduli space
\[
\mathcal{M}_0(a,a;\gamma_{0,0},\gamma_{1,0})
\]
consists of exactly two trajectories modulo translation, namely $[\nu_{0,+}]$ and $[\nu_{0,-}]$.
\end{thm}



\subsubsection{The $d=0$ augmented isomorphism}
To go from $d=0$ to small--$d$, we need not only the existence of the two $d=0$ trajectories, but also the corresponding linearized theory.



Let
\[
E_{0,\alpha}:=\nu_{0,\alpha}^{*}T\mathcal\HH_{\C}^2, \qquad \alpha\in\{+,-\}.
\]
Fix a small weight $0<\delta\ll1$ which is smaller than the corresponding nonzero spectral gap.  Define
\[
\begin{aligned}
X^{\rm fix}_{0,\alpha}
:=
\Bigl\{
  \xi\in
  W^{1,2}_{\delta}\bigl(\R_s;L^2([0,\tau],E_{0,\alpha})\bigr)
  \cap&L^2_{\delta}\bigl(\R_s;H^2_0([0,\tau]_t,E_{0,\alpha})\bigr)
  \ :\ \xi(s,0)=\xi(s,\tau)=0,\\[1mm]
  &\xi(s,\cdot)\to0
    \text{ as }|s|\to\infty,~\Pi_{\rm Bott}^{-}\xi=0,~\ell_{0,\alpha}(\xi)=0
\Bigr\}.
\end{aligned}
\]
Here $\Pi_{\rm Bott}^{-}\xi=0$ fixes the upper Morse--Bott coefficient to be zero, and
\[
\ell_{0,\alpha}(\xi):=\frac{
\langle \xi,\partial_s\nu_{0,\alpha}\rangle_{L^2_{\R}}}{\|\partial_s\nu_{0,\alpha}\|^2_{L^2_{\R}}}
\]
is the translation gauge. Define
\[
Y_{0,\alpha}:=L^2_{\delta}\bigl(\R_s;L^2([0,\tau]_t,E_{0,\alpha})\bigr)=
\left\{\eta(s,t)\in E_{0,\alpha}
: \int_{\R}e^{2\delta |s|}
\int_0^\tau |\eta(s,t)|^2\,dt\,ds<\infty
\right\}.
\]
The fixed-end linearized operator 
\[
D^{\rm fix}_{0,\alpha}:X^{\rm fix}_{0,\alpha}\longrightarrow
Y_{0,\alpha},
\]
is given by the linearization of the $d=0$ flow equation along
$\nu_{0,\alpha}$. Along the slice $\Sigma=Q(S^2)$,
\[
T\HH_{\C}^2\big|_{\Sigma}=T\Sigma\oplus iT\Sigma
\cong TS^2\oplus iTS^2.
\]
Then $D^{\rm fix}_{0,\alpha}$ splits into the tangential and normal blocks along the $S^2$-slice. 

For the tangential block, the kernel is given by $\partial_s\nu_{0,\alpha}$, which is removed by our translation gauge fixing condition $\ell_{0,\alpha}(\xi)=0$. That is,
\[\ker D^{\rm{fix},\top}_{0,\alpha}=0.\]

For the normal block, by the same positive-end rigidity Lemma \ref{lem:forward-tail-rigidity} and backward-uniqueness argument for $p$ in Subsection \ref{subsec:full-d=0}, we have
\[\ker D^{\rm{fix},\perp}_{0,\alpha}=0.\]

On the other hand, since we fixed the upper Bott coefficient to be zero in
$X^{\rm fix}_{0,\alpha}$, the missing target direction is the upper
soft obstruction associated with the normal $n=2$ Bott mode
\[
  \sigma_{B,\alpha}(t)
  =
  \sinh\left(\frac{2\pi i t}{\tau}\right)N
  =
  i\sin\left(\frac{2\pi t}{\tau}\right)N,
  \qquad
  N=(0,0,1).
\]


Let us consider an augmented linear map which fills the missing direction. Choose a smooth cutoff
$\chi_-(s)$ with $\chi_-(s)=1$ for $s\ll0$ and
$\chi_-(s)=0$ away from the upper end.  Set
\[
\widehat\sigma_{B,\alpha}(s,t):=
\chi_-(s)\sigma_{B,\alpha}(t).
\]

We claim that $D_{0,\alpha}\widehat\sigma_{B,\alpha}\in Y_{0,\alpha}$.
To prove it, we write the operator near the upper end in the form
\[
D_{0,\alpha}=\partial_s+A_\alpha(s),
\]
where $A_\alpha(s)$ converges exponentially to the asymptotic upper-end
operator $A_\alpha^-$. Since $\sigma_{B,\alpha}$ is the upper Bott mode, $A_\alpha^-\sigma_{B,\alpha}=0$. Then we have
\begin{align*}
D_{0,\alpha}\widehat\sigma_{B,\alpha}
&=D_{0,\alpha}\bigl(\chi_-(s)\sigma_{B,\alpha}(t)\bigr)\\
&=\chi_-'(s)\sigma_{B,\alpha}(t)+
\chi_-(s)A_\alpha(s)\sigma_{B,\alpha}(t) \\
&=\chi_-'(s)\sigma_{B,\alpha}(t)+
\chi_-(s)\bigl(A_\alpha(s)-A_\alpha^-\bigr)\sigma_{B,\alpha}(t).
\end{align*}
The first term is compactly supported in $s$. For the second term, the
exponential convergence of the tail gives, for some $\kappa>0$,
\[
\left\|\bigl(A_\alpha(s)-A_\alpha^-\bigr)\sigma_{B,\alpha}\right\|_{L^2_t}\le C e^{-\kappa |s|}
\qquad\text{as } s\to-\infty.
\]
Choosing the weight $\delta>0$ with $\delta<\kappa$, we obtain
\[
\int_{\mathbb R} e^{2\delta |s|}
\left\|D_{0,\alpha}\widehat\sigma_{B,\alpha}\right\|_{L^2_t}^2\,\dd s
<\infty.
\]
Thus we can define the following augmented linearized operator
\[
\widehat D_{0,\alpha}:X_{0,\alpha}^{\rm fix}\oplus\R\longrightarrow Y_{0,\alpha},\qquad
\widehat D_{0,\alpha}(\xi,\rho)
=D_{0,\alpha}^{\rm fix}\xi+
\rho\,D_{0,\alpha}\widehat\sigma_{B,\alpha}.
\]
By the above analysis, we have

\begin{lem}
\label{lem:upper-augmented-isomorphism}
The operator $\widehat D_{0,\alpha}$ is an isomorphism.
\end{lem}


\subsection{Case: $0<d\ll1$} 

We now extend our discussion from $d=0$ to small $d>0$ using perturbation method. The proof again has two parts. 

First, we have a small $d>0$ spectral analysis and construct a $d$--perturbative invertible augmented linear operator. The implicit function theorem, combined with a careful estimate of the soft eigenvalue direction, then produces two smooth solution branches for all sufficiently small $d>0$.

Second, we prove that these are no extra solutions by the uniform estimates, the no-escape and compactness argument. Then we get
\[
\#\mathcal M_d(a,b_d;\gamma_{0,d},\gamma_{1,d})=2
\]
for all sufficiently small $d>0$.



\subsubsection{The small $d$ spectral analysis}

Recall that at the critical path $\gamma_{j,d}$, we have obtained the linearization of the flow \eqref{eq:flow} with $\R$-linear asymptotic operator
\[
A_{j,d,\theta}=e^{i\theta}\,\mathcal{C}\,M_{\gamma_{j,d}}^{-1}\, L^{\C}_{j,d}.
\]

Applying the proof of Proposition \ref{prop:expconv} to the $j=0,1$ case, 
\begin{prop}\label{prop:exact-spectrum-small-d}
\begin{itemize}
    \item[(1)] At the lower end $j=0$, there is a uniform spectral gap:
\[
\mathrm{dist}(0,\sigma_{\R}(A_{0,d,\theta}))
\ge c_0>0,\qquad \text{for all } 0\le d\ll1.
\]
    \item[(2)] At the upper end $j=1$, the unique small normal mode is the normal
$n=2$ mode, for which
\[
\phi^{\mathrm{nor}}_{2,1}(d)
=\frac{4\pi^2+(d+2\pi i)^2}{\tau^2}
=\frac{d^2+4\pi i d}{\tau^2}.
\]
Hence
\[
\lambda_{\mathrm{soft}}(d)
:=|\phi^{\mathrm{nor}}_{2,1}(d)|
=\frac{d\sqrt{d^2+16\pi^2}}{\tau^2}
\sim\frac{4\pi}{\tau^2}d,\qquad d\to0^+.
\]
After removing this normal $n=2$ soft mode, the remaining spectrum at $j=1$ is uniformly bounded away from zero for all sufficiently small $d\ge0$.
\end{itemize}
\end{prop}


\begin{cor}\label{cor:fixed-d-exp-convergence}
There exists $d_0>0$ such that, for each fixed
$0<d\le d_0$ and each finite-energy trajectory (if it exists)
\[
u\in \mathcal M_d(a,b_d;\gamma_{0,d},\gamma_{1,d}),
\]
the following hold.

\begin{enumerate}
\item[(i)] As $s\to+\infty$, write
\[
u(s)=\exp_{\gamma_{0,d}}^{\C}\xi_+(s).
\]
Then, for every $m\ge0$ and every $0<\mu<c_0$, there exist
constants $C_{m,\mu}>0$ and $S_+>0$ such that
\[
\|\xi_+(s)\|_{H^m_t}\le C_{m,\mu}e^{-\mu s},\qquad s\ge S_+.
\]
\item[(ii)] As $s\to-\infty$, write
\[
u(s)=\exp_{\gamma_{1,d}}^{\C}\xi_-(s).
\]
Then, for every $m\ge0$ and every
$
0<\mu<\lambda_{\mathrm{soft}}(d)
$, there exist constants $C_{m,\mu}>0$ and $S_->0$ such that
\[
\|\xi_-(s)\|_{H^m_t}\le C_{m,\mu}e^{\mu s},
\qquad s\le -S_-.
\]
\end{enumerate}
\end{cor}



\subsubsection{Endpoint correction}

From $d=0$ to $d>0$, the endpoint at $t=\tau$ moves from $a$ to $b_d$.  To compare the $d$-flow equation with the $d=0$ trajectories in a fixed Banach space with homogeneous Dirichlet variations, we first fix this moving-endpoint
condition as follows.  Let
\[
B_r(z_0,z_1,z_2)
:=(\cosh r\,z_0+\sinh r\,z_1,
\sinh r\,z_0+\cosh r\,z_1,z_2)
\]
be the real Lorentz transformation. Then 
$B_d a=b_d$. For $\kappa(t)\in\PP_0^{\C}$, define
\begin{equation}\label{eq:Tkappa}
\mathcal T_d(\kappa)(t):=B_{dt/\tau}\kappa(t).
\end{equation}
Then we have
\[
\mathcal T_d:\mathcal P^{\C}_0\longrightarrow\mathcal P^{\C}_d,
\qquad\mathcal T_d(\kappa)(0)=a,
\qquad\mathcal T_d(\kappa)(\tau)=b_d,
\qquad\mathcal T_0=\mathrm{Id}.
\]
On bounded $C^2$-sets of paths,
\[
 \mathcal T_d(\kappa)=\kappa+O(d),
\qquad\partial_d\mathcal T_d(\kappa)\big|_{d=0}
=\frac{t}{\tau}K\kappa,
\]
where
\[
K(z_0,z_1,z_2):=(z_1,z_0,0).
\]
For each $\alpha\in\{+,-\}$, we define the approximate small-$d$ trajectory
\[
\nu^{\mathrm{app}}_{d,\alpha}
:=\mathcal T_d\nu_{0,\alpha}=\nu_{0,\alpha}+O(d).
\]
Then
\[\nu^{\mathrm{app}}_{d,\alpha}(s,0)=a,\qquad \nu^{\mathrm{app}}_{d,\alpha}(s,\tau)=b_d.\]

\begin{lem}\label{lem:endpoint-correction-orthogonal}
Let
\begin{equation}\label{eq:rho-alpha}
\rho_\alpha
:=\partial_d\nu^{\mathrm{app}}_{d,\alpha}\big|_{d=0}
=\frac{t}{\tau}K\nu_{0,\alpha}.    
\end{equation}
At the upper end $s\to-\infty$, we have
\[
\rho_\alpha(s,t)\longrightarrow
\frac{t}{\tau}K\gamma_{1,0}(t)=\frac{t}{\tau}T_{1,0}(t)
=\frac{t}{\tau}\left(\sinh\frac{2\pi i t}{\tau},
\cosh\frac{2\pi i t}{\tau},0\right).
\]
Thus the infinitesimal endpoint correction is orthogonal to 
the normal Bott direction $\sigma_{B,\alpha}(t)$.
\end{lem}


We now further formulate the local equation near $\nu^{\mathrm{app}}_{d,\alpha}$. 

\smallskip
Let $X^{\rm fix}_{d,\alpha}$ be the space of weighted sections along $v^{\rm app}_{d,\alpha}$ satisfying the same fixed-endpoint, decay,
upper-end Morse--Bott fixing, and phase-fixing conditions as $X^{\rm fix}_{0,\alpha}$, after identifying the bundles by parallel transport.
And let $Y_{d,\alpha}$ be the corresponding weighted $L^2$-space of sections along $v^{\rm app}_{d,\alpha}$.

Using the parallel transport along $\nu_{0,\alpha}$
and $\nu^{\rm app}_{d,\alpha}$, we have the identifications
\[
\mathcal P_{d,\alpha}:X^{\rm fix}_{0,\alpha}\longrightarrow
X^{\rm fix}_{d,\alpha},\qquad \Theta_{d,\alpha}:Y_{d,\alpha}\longrightarrow Y_{0,\alpha}.
\]
By construction, 
$\mathcal P_{d,\alpha}$ and $\Theta_{d,\alpha}$ depend $C^1$--smoothly on
$d$.  In particular,
\[
\mathcal P_{d,\alpha}=\operatorname{Id}+O(d),
\qquad\Theta_{d,\alpha}=\operatorname{Id}+O(d)
\]
as bounded maps on the corresponding weighted spaces.



%For $d>0$, let $\sigma_{B,\alpha,d}(t)$ be the 
%small $d$ upper soft eigenmode corresponding to $\lambda_{\rm soft}(d)$. Define
%\[\widehat\sigma_{B,\alpha,d}(s,t):=\chi_-(s)e^{\lambda_{\rm soft}(d)s}\sigma_{B,\alpha,d}(t),
%\qquad d>0.\]
%The limit as $d\rightarrow0^+$ recovers $\widehat\sigma_{B,\alpha}(s,t)$ that we used in the augmented isomorphism $\widehat{D}_{0,\alpha}$.


For
\[
(\zeta,\rho)\in X^{\rm fix}_{0,\alpha}\oplus\R,
\]
define the perturbation along $\nu^{\rm app}_{d,\alpha}$ by
\[
\mathcal I_{d,\alpha}(\zeta,\rho):=
\mathcal P_{d,\alpha}\zeta+
\rho\,\mathcal{P}_{d,\alpha}\widehat\sigma_{B,\alpha}=\zeta+\rho\widehat{\sigma}_{B,\alpha}+O(d).
\]
This map is $C^1$ in $d$.


Furthermore, we define the parameter-dependent augmented nonlinear map 
\[
\widehat{\mathcal F}_{\alpha}:
[0,\varepsilon)_d\times
\bigl(X^{\rm fix}_{0,\alpha}\oplus\R\bigr)
\longrightarrow Y_{0,\alpha},
\]
to be
\[
\widehat{\mathcal F}_{\alpha}(d,\zeta,\rho)
:=\Theta_{d,\alpha}\left[\mathcal F_d
\left(\exp_{\nu^{\rm app}_{d,\alpha}}^{\C}
\bigl(\mathcal I_{d,\alpha}(\zeta,\rho)\bigr)\right)
\right],\qquad \mathcal{F}_d(z):=\partial_sz-P_z(e^{i\theta(d)}\eta\bar{z}_{tt}).
\]
Here the exponential map is taken pointwise along
$\nu^{\rm app}_{d,\alpha}$. 

By construction, the augmented nonlinear map $\widehat{\mathcal F}_{\alpha}(d,\zeta,\rho)$ is $C^1$ in $(d,\zeta,\rho)$ near $(0,0,0)$ (in the one-sided sense at $d=0$).  Moreover, we have
\begin{itemize}
    \item \[
\widehat{\mathcal F}_{\alpha}(d,0,0)
=\widehat{\mathcal F}_{\alpha}(0,0,0)+O(d)=O(d),
\]
and the sharper estimate (proved in the following Lemma \ref{lem:quadratic-upper-bott})
\[
\Pi_{\rm Bott}^-\Bigl(\widehat{\mathcal F}_{\alpha}(d,0,0)\Bigr)=O(d^2)
\quad\text{in }~\R\sigma_{B,\alpha},
\qquad d\to0^+.
\]
    \item \[
D_{(\zeta,\rho)}
\widehat{\mathcal F}_{\alpha}(d,0,0)
=D_{(\zeta,\rho)}
\widehat{\mathcal F}_{\alpha}(0,0,0)
+O(d)
\]
as bounded operators. In particular,
\[
D_{(\zeta,\rho)}
\widehat{\mathcal F}_{\alpha}(0,0,0)(\xi,r)
=D^{\rm fix}_{0,\alpha}\xi
+r\,D_{0,\alpha}\widehat\sigma_{B,\alpha}.
\]
This is precisely our augmented linearized operator $\widehat{D}_{0,\alpha}$.
\end{itemize}


\begin{lem}\label{lem:quadratic-upper-bott}
For each $\alpha\in\{+,-\}$, we have
\[
\Pi_{\rm Bott}^-\Bigl(\widehat{\mathcal F}_{\alpha}(d,0,0)\Bigr)=O(d^2)
\quad\text{in }~ \R\,\sigma_{B,\alpha},
\qquad d\to0^+.
\]   
\end{lem}

\begin{proof}
At $(\zeta,\rho)=(0,0)$, we have $\mathcal I_{d,\alpha}(0,0)=0$, hence
\[
\widehat{\mathcal F}_{\alpha}(d,0,0)
=\Theta_{d,\alpha}
\left[\mathcal F_d(\nu^{\rm app}_{d,\alpha})
\right].
\]    
We can compute
\begin{align*}
\frac{\dd}{\dd d}\bigg|_{d=0}\Theta_{d,\alpha}
\left[\mathcal F_d(\nu^{\rm app}_{d,\alpha})
\right]
&=\left.\partial_d\Theta_{d,\alpha}\right|_{d=0}
\left[\mathcal F_0(\nu_{0,\alpha})\right]
+\Theta_{0,\alpha}\left[
D_z\mathcal F_0(\nu_{0,\alpha})\rho_\alpha
+\left.\partial_d\mathcal F_d(\nu_{0,\alpha})
\right|_{d=0}\right]\\
&=D_{0,\alpha}\rho_\alpha
+\left.\partial_d\mathcal F_d(\nu_{0,\alpha})
\right|_{d=0},
\end{align*}
where we use $\mathcal F_0(\nu_{0,\alpha})=0$ and $\Theta_{0,\alpha}=\Id$.

By Lemma \ref{lem:endpoint-correction-orthogonal}, we have
\[\Pi_{\rm Bott}^-(D_{0,\alpha}\rho_\alpha)=0.\]
On the other hand, the $d$-dependence in $\mathcal F_d$ at $\nu_{0,\alpha}$ is the factor
$e^{i\theta(d)}$, we get
\[
\left.\partial_d\mathcal F_d(\nu_{0,\alpha})\right|_{d=0}
=-P_z\left(i\theta'(0)e^{i\theta(0)}
\eta\overline{\nu_{0,\alpha}}_{tt}\right)=
\frac{i}{\pi}\partial_s\nu_{0,\alpha}.
\]
Hence, we also have
\[\Pi_{\rm Bott}^-\left(\frac{i}{\pi}\partial_s\nu_{0,\alpha}\right)=0.\]
In summary, we obtain
\[
\Pi_{\rm Bott}^-\Bigl(\widehat{\mathcal F}_{\alpha}(d,0,0)\Bigr)=\Pi_{\rm Bott}^-\Bigl(\Theta_{d,\alpha}\mathcal F_d(\nu^{\rm app}_{d,\alpha})\Bigr)=
O(d^2).\qedhere\]
\end{proof}




\begin{prop}\label{prop:at-least-two-small-d}
There exists $d_0>0$ and two smooth branches
\[
d\longmapsto [\nu_{d,+}],
\qquad d\longmapsto [\nu_{d,-}],
\qquad \quad 0<d<d_0,
\]
such that for $\alpha\in \{+,-\}$
\[
[\nu_{d,\alpha}]\in \mathcal{M}_d(a,b_d;\gamma_{0,d},\gamma_{1,d}),
\qquad
[\nu_{d,\alpha}]\to [\nu_{0,\alpha}]
\quad (d\to 0^+).
\]
In particular,
\[
\#\mathcal{M}_d(a,b_d;\gamma_{0,d},\gamma_{1,d})\ge 2\qquad (0<d<d_0).
\]
\end{prop}

\begin{proof}
By Lemma \ref{lem:upper-augmented-isomorphism}, $\widehat{D}_{0,\alpha}$ is an
isomorphism. The Banach-space implicit function theorem, applied to the fixed cut-off augmented map $\widehat{\mathcal{F}}_{\alpha}$, gives for all sufficiently small $d$, a unique small pair
$(\zeta_{d,\alpha},\rho_{d,\alpha})$ such that
\[
\widehat{\mathcal F}_\alpha
(d,\zeta_{d,\alpha},\rho_{d,\alpha})=0.
\]

Moreover, since
\[
\widehat{\mathcal F}_\alpha(d,0,0)=O(d)
\]
in $Y_{0,\alpha}$, and the inverse of the linearized operator remains
uniformly bounded near the isomorphism $\widehat D_{0,\alpha}$, we obtain
\begin{equation}\label{eq:zeta-rho}
\|\zeta_{d,\alpha}\|+|\rho_{d,\alpha}|\le C\|\widehat{\mathcal F}_\alpha(d,0)\|_{Y_{0,\alpha}}
\le Cd.
\end{equation}

Let
\[
z_{d,\alpha}:=\exp_{\nu^{\rm app}_{d,\alpha}}^{\C}
\Bigl(\mathcal I_{d,\alpha}\left(\zeta_{d,\alpha},\rho_{d,\alpha}
\right)\Bigr)
\]
be the small trajectory constructed above. Let
$\sigma_{{\rm soft},d}$ be the normalized upper $\lambda_{\rm soft}(d)$ eigenmode.
As $s\to-\infty$, we can expand
\[
Y_{d,\alpha}(s,t):=\left(\exp_{\gamma_{1,d}(t)}^{\C}\right)^{-1}(z_{d,\alpha}(s,t))=\beta_{d,\alpha}
e^{\lambda_{\rm soft}(d)s}
\sigma_{{\rm soft},d}(t)
+Y_{d,\alpha}^{\perp}(s,t),\qquad s\to-\infty,
\]
where
\[
\lambda_{\rm soft}(d)=c d+O(d^2),~~c>0;
\qquad
\Pi_{\rm soft}^-(d)Y_{d,\alpha}^{\perp}=0.
\]


Since
\[
\widehat{\mathcal F}_{\alpha}
\left(d,\zeta_{d,\alpha},\rho_{d,\alpha}\right)=0,
\]
we expand at $(\zeta,\rho)=(0,0)$:
\[
0=\widehat{\mathcal F}_{\alpha}(d,0,0)+D_{(\zeta,\rho)}
\widehat{\mathcal F}_{\alpha}(d,0,0)
\left(\zeta_{d,\alpha},\rho_{d,\alpha}
\right)+\mathcal N_{d,\alpha}
\left(\zeta_{d,\alpha},\rho_{d,\alpha}\right),
\]
where
\[
\left\|\mathcal N_{d,\alpha}(\zeta,\rho)\right\|
\le C\left(\|\zeta\|^2+|\rho|\,\|\zeta\|+|\rho|^2
\right).
\]

Projecting the equation to the upper soft channel gives
\[
0=\Pi_{\rm soft}^-(d)\widehat{\mathcal F}_{\alpha}(d,0,0)+\lambda_{\rm soft}(d)\beta_{d,\alpha}+
R_{d,\alpha}.
\]
Here $R_{d,\alpha}$ contains the transverse linear term and the nonlinear Taylor remainder.  


Let
\[
\mathscr L_{d,\alpha}
:=D_\zeta\widehat{\mathcal F}_\alpha(d,0,0)
:X^{\rm fix}_{0,\alpha}\longrightarrow Y_{0,\alpha}.
\]
Then
\[
\mathscr L_{d,\alpha}\zeta
=\left.\frac{d}{d\varepsilon}\right|_{\varepsilon=0}
\widehat{\mathcal F}_\alpha(d,\varepsilon\zeta,0),\qquad\text{and}\qquad \mathscr L_{0,\alpha}=D_{0,\alpha}^{\rm fix}.
\]

Moreover, by the $C^1$-dependence of
$\widehat{\mathcal F}_\alpha$ on $d$,
\[
\mathscr L_{d,\alpha}
=\mathscr L_{0,\alpha}+O(d)\quad\text{in }~\mathcal L
\left(X^{\rm fix}_{0,\alpha},Y_{0,\alpha}\right).
\]
After the fixed upper-end trivialization, the soft projection also
satisfies
\[
\Pi_{\rm soft}^-(d)=\Pi_{\rm Bott}^-+O(d).
\]
Hence
\[
\begin{aligned}
\Pi_{\rm soft}^-(d)\mathscr L_{d,\alpha}\zeta
&=\bigl(\Pi_{\rm Bott}^-+O(d)\bigr)
\bigl(\mathscr L_{0,\alpha}+O(d)\bigr)\zeta \\
&=\Pi_{\rm Bott}^-D_{0,\alpha}^{\rm fix}\zeta
+O(d)\|\zeta\|=O(d)\|\zeta\|.
\end{aligned}
\]
Thus the projected equation becomes
\[
\lambda_{\rm soft}(d)\beta_{d,\alpha}
=-\Pi_{\rm Bott}^-\widehat{\mathcal F}_{\alpha}(d,0,0)+O\left(d^2+d\|\zeta_{d,\alpha}\|
+\|\zeta_{d,\alpha}\|^2+|\rho_{d,\alpha}|\,\|\zeta_{d,\alpha}\|+|\rho_{d,\alpha}|^2\right).
\]
By Lemma \ref{lem:quadratic-upper-bott}, \eqref{eq:zeta-rho} and $\lambda_{\rm soft}(d)\sim d$, we obtain
\[
|\beta_{d,\alpha}|\le Cd.
\]
%In particular, the small upper eigenvalue does not produce an $O(1)$ Bott drift.

Then we get a finite-energy solution
\[
\nu_{d,\alpha}\in
\mathcal M_d(a,b_d;\gamma_{0,d},\gamma_{1,d}).
\]
Since the two $d=0$ trajectories $\nu_{0,+}$ and $\nu_{0,-}$ are
distinct, their small $d$ continuations remain distinct after
decreasing $d_0>0$ if necessary.  Therefore
\[
\#\mathcal M_d(a,b_d;\gamma_{0,d},\gamma_{1,d})\ge2
\]
for all sufficiently small $d>0$.    
\end{proof}




\subsubsection{The normal $O(d)$ control}

To prove exactness of the count, we need a small-$d$ no-escape result. The idea is to use the first-exit argument. To do that, let us establish some useful estimates.

\smallskip
First along the $d=0$ slice,  we will prove that the $d$-twisted vector field has a uniformly order-$d$ normal component. Recall from Section \ref{sec:heat flow on S2} 
\[
\Sigma=Q(S^2)\subset \HH^2_{\C}.
\]
For every smooth path $x:[0,\tau]\to S^2$, we regard $Qx$ as a path in
$\Sigma$. Along $Qx$, the Hermitian metric on $\HH^2_{\C}$ gives the
orthogonal splitting
\[
T_{Qx(t)}\HH^2_{\C}=T_{Qx(t)}\Sigma
\oplus N_{Qx(t)}\Sigma,\qquad 0\le t\le \tau .
\]
Let
\[
\pi_{Dx}^{\perp}:L^2\bigl([0,\tau],(Qx)^*T\HH^2_{\C}\bigr)
\longrightarrow L^2\bigl([0,\tau],(Qx)^*N\Sigma\bigr)
\]
denote the Hermitian orthogonal projection onto the normal component of $L^2$-sections along $Qx$.


For simplicity, denote the $d$-twisted Morse vector field by
\[
V_d(z):=P_z\bigl(e^{i\theta(d)}\eta \overline{z}_{tt}\bigr),\quad\text{with}\quad \theta(d)=\pi-\arctan\frac d\pi.
\]
Along the slice, its
normal component 
\[
V_d^\perp(Qx):=\pi_{Qx}^{\perp}\bigl(V_d(Qx)\bigr)
\]
measures the part of the $d$-twisted vector field which pushes the flow out of the $d=0$ slice $\Sigma$.

\begin{lem}\label{lem:order-d-almost-tangent}
With the notations above, there exists a constant $C>0$, independent of $x$ on bounded $C^2$-sets
and independent of all sufficiently small $d\ge 0$, such that
\[
\bigl\|V_d^\perp(Qx)\bigr\|_{L_t^2}\le Cd .
\]
%In particular, $V_d^\perp(Qx)$ vanishes at $d=0$ and depends smoothly on $d$. %Hence the normal forcing along the entire $d=0$ model is uniformly of order $d$.
\end{lem}

\begin{proof}
At $d=0$, we have $\theta(0)=\pi$, and  the preceding analysis in Section \ref{sec:heat flow on S2} shows
\[
V_0(Qx)\in L^2\bigl([0,\tau],(Qx)^*T\Sigma\bigr).
\]
Applying the normal projection gives
\[
V_0^\perp(Qx)=
\pi_{Qx}^{\perp}\bigl(V_0(Qx)\bigr)=0.
\]

Now fix a bounded $C^2$-set of slice paths. On such a set, the map
\[
(d,z)\longmapsto V_d(z)
=P_z\bigl(e^{i\theta(d)}\eta\overline{z}_{tt}\bigr)
\]
depends smoothly on $d$ and on $z$. Moreover, the Hermitian orthogonal splitting
\[
T_{Qx(t)}\HH^2_{\C}=T_{Qx(t)}\Sigma
\oplus N_{Qx(t)}\Sigma
\]
varies smoothly with the base point $Qx(t)$. Hence the normal projection
$\pi_{Qx}^{\perp}$ also varies smoothly with $x$. Therefore the map
\[
(d,x)\longmapsto V_d^\perp(Qx)
:=\pi_{Qx}^{\perp}\bigl(V_d(Qx)\bigr)
\]
is smooth on bounded $C^2$-sets of slice paths.

Since this map vanishes 
at $d=0$, the mean-value theorem in the
parameter $d$ gives
\[
\bigl\|V_d^\perp(Qx)\bigr\|_{L_t^2}=
\bigl\|V_d^\perp(Qx)-V_0^\perp(Qx)\bigr\|_{L_t^2} 
\le d\,\sup_{0\le \delta\le d}\left\|\frac{\partial}{\partial \delta}V_\delta^\perp(Qx)\right\|_{L_t^2}.
\]
On a fixed bounded $C^2$-set of paths, the standard compactness argument
shows that the last supremum is bounded by a constant $C$. Thus
\[
\bigl\|V_d^\perp(Qx)\bigr\|_{L_t^2}\le Cd. \qedhere
\]
\end{proof}





\subsubsection{The small-$d$ no-escape theorem}


\smallskip
Set
\[
\mathcal{K}_0:=\{\gamma_{1,0}\}\cup
\{\nu_{0,+}(s,\cdot):s\in \R\}\cup
\{\nu_{0,-}(s,\cdot):s\in \R\}\cup
\{\gamma_{0,0}\}.
\]
It is the compactified $d=0$ model. We use the following three kinds of local charts:
\begin{itemize}
\item[(a)] near $\gamma_{1,0}$, a Morse--Bott end chart, with the Bott direction separated from the transverse directions;
\item[(b)] along $\nu_{0,\pm}$, the tubular neighborhood charts;
\item[(c)] near $\gamma_{0,0}$, an ordinary Morse end chart.
\end{itemize}
The asymptotic expansions at $\gamma_{1,0}$, $\gamma_{0,0}$, established in the previous
subsections, will be used to glue these charts with uniform constants.

Now fix an interval $I\subset\R_s$ on which one of the above local
charts is used. We denote by
\[
\kappa_I:I\times[0,\tau]\longrightarrow \HH^2_{\C}
\]
the corresponding $d=0$ reference path in $\mathcal{K}_0$. Using the endpoint correction \eqref{eq:Tkappa}, set
\[
\kappa_{I,d}:=\mathcal T_d(\kappa_I).
\]
Then
\[
\kappa_{I,d}(s,0)=a,\qquad
\kappa_{I,d}(s,\tau)=b_d,\qquad
\kappa_{I,d}=\kappa_I+O(d).
\]
A nearby $d$-trajectory is written in the form
\[
z(s,\cdot)=\exp_{\kappa_{I,d}(s,\cdot)}^{\C}Y_I(s,\cdot),\qquad s\in I.
\]
Then the fixed-endpoint condition for $z$ becomes the
homogeneous Dirichlet condition
\[
Y_I(s,0)=0,\qquad Y_I(s,\tau)=0.
\]

Let $E_{I,d}:=\kappa_{I,d}^{*}T\HH^2_{\C}$. Since $\kappa_{I,d}$ is $O(d)$-close to $\kappa_I$, we identify
$E_{I,d}$ with
\[
E_I:=E_{I,0}=\kappa_I^{*}T\HH^2_{\C}
\]
by parallel transport along the short curves
$r\mapsto \kappa_{I,r}(s,t)$, $0\le r\le d$. Under this identification,
all $d$-dependent changes of the reference path contribute only $O(d)$
terms. %and will be included in the forcing term $f_d$.

Define $\mathcal X(I)$ to be the
space of sections $Y_I$ of $E_I$ such that
\[
Y_I\in H^1(I;L_t^2(E_I))\cap
L^2(I;H_t^2(E_I)\cap H^1_{0,t}(E_I))
\cap L^\infty(I;H^1_{0,t}(E_I)),
\]
satisfying 
\begin{enumerate}[label=\textup{(\alph*)}]
\setlength{\itemsep}{-1pt}
    \item near the Morse--Bott end $\gamma_{1,0}$, 
    \[Y_I \perp \sigma_B;\]
    \item on a middle chart around $\nu_{0,\alpha}$,
    \[Y_I \perp \partial_s\nu_{0,\alpha};\]
    \item near the Morse end $\gamma_{0,0}$, $Y_I$ satisfies the decaying Morse-end condition.
\end{enumerate}
Here the orthogonality is taken in the relevant $L^2$-sense. The norm is
\[
\|Y_I\|_{\mathcal X(I)}:=\|Y_I\|_{H^1(I;L_t^2)}+
\|Y_I\|_{L^2(I;H_t^2)}+\|Y_I\|_{L^\infty(I;H_t^1)}.
\]
%The target space is
%\[\mathcal Y(I):=L^2(I;L_t^2(E_I)),\qquad\|F\|_{\mathcal Y(I)}:=\|F\|_{L^2(I;L_t^2)}.\]
%On half-infinite end intervals, the same notation will be used for the
%corresponding exponentially weighted versions of these spaces. %, with a fixed weight chosen smaller than the relevant spectral gap after the Bott or translation direction has been removed.


For each chart, let $L_I(s)$ denote the $d=0$ linearized operator along the reference map $\kappa_I$. Then we can write the full $d=0$ linearized operator in the $I$-chart as
\[
\nabla_s^{\kappa_I}-L_I(s).
\]

The target space is
\[
\mathcal Y(I):=L^2(I;L_t^2(E_I)),\qquad
\|F\|_{\mathcal Y(I)}:=\|F\|_{L^2(I;L_t^2)}.
\]
On half-infinite end intervals, the same notation will be used for the
corresponding exponentially weighted versions of these spaces.


\begin{prop}\label{prop:YIestimate}
There exist $\rho_*>0$, constants $C_0,C_1,C_2>0$, and a system of local
charts as above such that the following holds. Let $0<\rho\le\rho_*$ and let $d\ge0$ be sufficiently small. If a $d$-trajectory $z$
\[
z(s,\cdot)=\exp_{\kappa_{I,d}(s,\cdot)}Y_I(s,\cdot),\qquad s\in I.
\]
stays in the $\rho$-neighborhood of $\mathcal{K}_0$ on an interval $I$, then 
$Y_I\in\mathcal X(I)$ satisfies
\[
\nabla_s^{\kappa_I}Y_I-L_I(s)Y_I
=f_{I,d}(s)+R_{I,d}(Y_I).
\]
The two remainders (forcing and nonlinear reminder respectively) satisfy
\[
\|f_{I,d}\|_{\mathcal Y(I)}\le C_0d,
\]
and
\[
\|R_{I,d}(Y_I)\|_{\mathcal Y(I)}
\le C_1\left(\rho\|Y_I\|_{\mathcal X(I)}
+\|Y_I\|_{\mathcal X(I)}^2\right).
\]
Moreover, %the sliced linearized operator satisfies the uniform inverse estimate
\[
\|Y_I\|_{\mathcal X(I)}\le
C_2\|\nabla_s^{\kappa_I}Y_I-L_I(s)Y_I\|_{\mathcal Y(I)}.
\]
On every bounded interval $I=[-R,R]$, bounded subsets of $\mathcal X(I)$
are relatively compact in $C^0(I;H_t^1)$.
\end{prop}


\begin{proof}
Write the trajectory in the chosen chart as
\[
z(s,\cdot)=\exp_{\kappa_{I,d}(s,\cdot)}^{\C}Y_I(s,\cdot),\qquad Y_I(s,0)=Y_I(s,\tau)=0 .
\]
Substituting this expression into the $d$-flow equation $\partial_s z=V_d(z)$ and expanding around the reference trajectory $\kappa_I$ gives
\[
\nabla^{\kappa_I}_{s}Y_I-L_I(s)Y_I
=f_{I,d}(s)+R_{I,d}(Y_I).
\]
Here 
\begin{itemize}
\setlength{\itemsep}{-1pt}
    \item $L_I(s)$ is the $d=0$ sliced linearized operator along $\kappa_I$.
    \item The term $f_{I,d}$ is the $Y_I$-independent forcing term:
    \[
 f_{I,d}
 =\bigl(V_d(\kappa_{I,d})-\partial_s\kappa_{I,d}\bigr)^{\perp}.
\]
It therefore measures how far $\kappa_{I,d}$ is from being an exact
$d$-flow trajectory. By the estimates in Lemma \ref{lem:order-d-almost-tangent} for $(V_d-V_0)^{\perp}$ and for
$\kappa_{I,d}-\kappa_I$, we have uniformly
\[
\|f_{I,d}\|_{\mathcal Y(I)}\le C_0 d.
\]
    \item The term $R_{I,d}(Y_I)$ is the remaining nonlinear part. Thus $R_{I,d}(0)=0$, it is quadratic in $Y_I$ at $d=0$,
and for small $d$ it satisfies
\[
\|R_{I,d}(Y_I)\|_{\mathcal Y(I)}\le C_1\bigl(\rho\,\|Y_I\|_{\mathcal X(I)}+\|Y_I\|_{\mathcal X(I)}^2\bigr).
\]
\end{itemize}


It only remains to recall the uniform estimate for the sliced linearized
operator.  On the two ends this is the standard Morse--Bott estimate after
removing the Bott direction.  On the compact middle part it follows from
the regularity and nondegeneracy of the two $d=0$ trajectories
$\nu_{0,+}$ and $\nu_{0,-}$. Since the compactified model
$\mathcal K_0$ is covered by finitely many such charts, we obtain, after
possibly decreasing $\rho_*$,
\[
\|Y_I\|_{\mathcal X(I)}\le C_2\bigl\|
\nabla^{\kappa_I}_{s}Y_I-L_I(s)Y_I
\bigr\|_{\mathcal Y(I)} .
\]

Finally, if $I=[-R,R]$ is bounded, the definition of the
$\mathcal X(I)$-norm controls
\[
H^1(I;L_t^2)\cap L^\infty(I;H_t^1)\cap L^2(I;H_t^2).
\]
Since $H_t^2\Subset H_t^1\hookrightarrow L_t^2$ on the compact
$t$-interval, the Rellich--Aubin--Lions compactness theorem implies that
bounded subsets of $\mathcal X(I)$ are relatively compact in $C^0(I;H_t^1)$.
\end{proof}


\begin{prop}\label{prop:no-escape-small-d}
There exist $d_*>0$ and a neighborhood $\mathcal N(\mathcal{K}_0)\subset \PP^{\C}_0$
of the compactified $d=0$ model such that the following holds. For every
$0<d\le d_*$, every finite-energy trajectory
\[
z\in \mathcal{M}_d(a,b_d;\gamma_{0,d},\gamma_{1,d})
\]
satisfies
\[
\mathcal T_d^{-1}z(s,\cdot)\in \mathcal N(\mathcal{K}_0)
\qquad\text{for all }s\in\R.
\]
\end{prop}

\begin{proof}
We prove it by a first-exit contradiction.

Let $\rho^*>0$, $C_0,C_1,C_2>0$ be the constants in Proposition \ref{prop:YIestimate}.
Choose $0<\rho<\rho^*$ so small that $C_2C_1\rho<1/2$. Suppose that the theorem were false. Then there exist $d_n\to0^+$ and
trajectories
\[z_n\in \mathcal{M}_{d_n}(a,b_{d_n};\gamma_{0,d_n},\gamma_{1,d_n})\]
such that the pulled-back trajectories
\[\widetilde z_n:=\mathcal T_{d_n}^{-1}z_n\in \PP^{\C}_0 \quad\text{leave }~\mathcal N_\rho(\mathcal{K}_0).\]
As $s\to-\infty$, 
\[z_n(s,\cdot)\to\gamma_{1,d_n},\qquad
\mathcal{T}_{d_n}^{-1}\gamma_{1,d_n}\to\gamma_{1,0}
,\]
then the path $\widetilde z_n(s,\cdot)$ lies in $\mathcal N_\rho(\mathcal{K}_0)$ for
all $s\ll0$. Hence there is a first exit time $s_n^*$
such that
\[
\widetilde z_n((-\infty,s_n^*])\subset\mathcal N_\rho(\mathcal{K}_0),\qquad
\widetilde z_n(s_n^*)\in\partial\mathcal N_\rho(\mathcal{K}_0).
\]
Applying Proposition \ref{prop:YIestimate} on the interval
$ I_n:=(-\infty,s_n^*]$, we can write
\[
z_n(s,\cdot)=\exp_{\kappa_{I_n,d_n}(s,\cdot)}^{\C}
Y_n(s,\cdot),\qquad s\in I_n,
\]
where $Y_n\in\mathcal X(I_n)$ satisfies
\[
\nabla_s^{\kappa_{I_n}}Y_n-L_{I_n}(s)Y_n=
f_{I_n,d_n}(s)+R_{I_n,d_n}(Y_n).
\]
Furthermore, for each $n$ and each $S\le s_n^*$, set
\[
A_n(S):=\|Y_n\|_{\mathcal X((-\infty,S])}.
\]
Again, by Proposition \ref{prop:YIestimate} on the interval $(-\infty,S]$ and our assumption $C_2C_1\rho<1/2$,
\[
\frac12 A_n(S)\le Cd_n+C A_n(S)^2 .
\]
This scalar inequality implies that, for $d_n$ sufficiently small, we have
\[
A_n(S)\le C'd_n\quad\text{ or }\quad A_n(S)\ge c_0.
\]
In particular, $c_0>0$ is independent of $n$.

For each fixed $n$, the convergence
\[
z_n(s,\cdot)\to\gamma_{1,d_n}\qquad (s\to-\infty)
\]
implies that %, after applying the endpoint correction and the upper-end chart,
the normal coordinate satisfies
\[
A_n(S)\to0 \qquad (S\to-\infty).
\]
Hence there exists $S_n\le s_n^*$ such that
\[
A_n(S_n)\le C'd_n.
\]
Since $S\mapsto A_n(S)$ is continuous from the right on finite subintervals
and monotone nondecreasing, the value $A_n(S)$ cannot pass from the small
branch to the large branch without crossing the forbidden gap between them.
Therefore
\[
A_n(S)\le C'd_n \qquad\text{for all }S\le s_n^*.
\]


By the definition of the $\mathcal X$-norm, we have
\[
\sup_{s\le s_n^*}\|Y_n(s,\cdot)\|_{H_t^1}
\le C''d_n.
\]
The complexified exponential map and the endpoint correction $\mathcal T_{d_n}$
are smooth and uniformly $C^1$-close to the identity for $d_n$ small.
Therefore
\[
\operatorname{dist}_{H_t^1}
\bigl(\mathcal T_{d_n}^{-1}z_n(s,\cdot),\mathcal{K}_0\bigr)
\le C''d_n\qquad\text{for all }s\le s_n^*.
\]
For $n$ sufficiently large, $C''d_n<\rho/2$. Hence
\[
\widetilde z_n((-\infty,s_n^*])
\subset\mathcal N_{\rho/2}(K_0).
\]
This contradicts the first exit time assumption
\[
\widetilde z_n(s_n^*)\in\partial\mathcal N_\rho(\mathcal{K}_0).\qedhere
\]
\end{proof}

\subsubsection{Compactness as $d_n\to 0^+$ and the exact-two count}

We now combine no-escape with local continuation. A limiting sequence could
fail to give the desired $d=0$ trajectory only by breaking, or by developing
a zero-energy drift along the upper Bott family at $d=0$. The first
phenomenon is excluded, since %excluded by the action filtration: $F_{\theta(d)}$ strictly decreases on nonconstant flow levels, while $G_{\theta(d)}$ is preserved, and
for $0<d\ll1$ the no-escape neighborhood contains no intermediate critical
point $\zeta$ with
\[
F_{\theta(d)}(\gamma_{0,d})< F_{\theta(d)}(\zeta)
< F_{\theta(d)}(\gamma_{1,d}),\qquad
G_{\theta(d)}(\zeta)=G_{\theta(d)}(\gamma_{0,d})
=G_{\theta(d)}(\gamma_{1,d}).
\]
The second phenomenon is excluded by our estimate $|\beta_{d_n,\alpha}|\leq Cd_n$ in Proposition \ref{prop:at-least-two-small-d}.


\begin{prop}\label{prop:compactness-dn-to-0}
Let $d_n\to 0^+$ and choose
\[
[u_n]\in \mathcal{M}_{d_n}(a,b_{d_n};\gamma_{0,d_n},\gamma_{1,d_n}).
\]
Then, after $s$-translation, a subsequence converges in $C^\infty_{\mathrm{loc}}$ to an element of
\[
\mathcal{M}_0(a,a;\gamma_{0,0},\gamma_{1,0})=
\{[\nu_{0,+}],[\nu_{0,-}]\}.
\]
\end{prop}

\begin{proof}
Choose representatives $u_n$ and translate them in the $s$-variable so
that, for a fixed regular value
\[
c_*\in\bigl(F_{\theta(0)}(\gamma_{0,0}),F_{\theta(0)}(\gamma_{1,0})\bigr),
\]
we have
\[
F_{\theta(d_n)}(u_n(0,\cdot))=c_* .
\]
For $n$ large, $c_*$ still lies strictly between the two endpoint critical
values.

By Proposition \ref{prop:no-escape-small-d}, after applying the endpoint correction
$\mathcal T_{d_n}^{-1}$, the trajectories $u_n$ remain in a fixed
neighborhood of $\mathcal{K}_0$. And by Proposition \ref{prop:YIestimate},
after passing to a subsequence, we obtain convergence on each bounded strip.
A diagonal argument gives a limit
\[
u_\infty:\R\times[0,\tau]\longrightarrow \HH^2_{\C}.
\]

Since $d_n\to0$ and $\mathcal T_{d_n}\to\mathrm{Id}$, the limit $u_\infty$ solves the $d=0$ flow
equation. Standard local parabolic bootstrapping improves the $C_{\mathrm{loc}}^{\infty}$ convergence. 

The normalization also passes to the limit and gives
\[
F_{\theta(0)}(u_\infty(0,\cdot))=c_*.
\]
Thus, $u_\infty$ is nonconstant. Moreover, the energy identity gives
\[
\int_{\R}\|\partial_su_n(s,\cdot)\|_{L_t^2}^2\,\dd s
=F_{\theta(d_n)}(\gamma_{1,d_n})-F_{\theta(d_n)}(\gamma_{0,d_n}).
\]
By lower semicontinuity, $u_\infty$ is a finite-energy $d=0$ trajectory.
The preceding discussion rules out broken intermediate levels,
and the transversality of the endpoint branch to the upper Bott family rules
out zero-energy Bott drift. Therefore the upper and lower limits of
$u_\infty$ are forced to be
\[
\gamma_{1,0}\quad\text{and}\quad\gamma_{0,0}.
\]
Thus
\[
[u_\infty]\in M_0(a,a;\gamma_{0,0},\gamma_{1,0}).
\]
By Theorem \ref{thm:exact-two-d0},
\[
M_0(a,a;\gamma_{0,0},\gamma_{1,0})
=\{[\nu_{0,+}],[\nu_{0,-}]\}.
\]
Hence, after translating in $s$ and passing to a subsequence,
$[u_n]$ converges in $C^\infty_{\mathrm{loc}}$ to $[\nu_{0,+}]$ or $[\nu_{0,-}]$.
\end{proof}


\begin{thm}\label{thm:exact-two-small-d}
There exists $\varepsilon>0$ such that
\[
\#\mathcal{M}_d(a,b_d;\gamma_{0,d},\gamma_{1,d})=2
\qquad\text{for all }0<d<\varepsilon.
\]
\end{thm}

\begin{proof}

Proposition \ref{prop:at-least-two-small-d} gives at least two trajectories for all small $d>0$. Suppose there are additional trajectories $[w_n]$ with $d_n\to 0^+$. By Proposition \ref{prop:compactness-dn-to-0}, after translation and subsequence extraction we get a limit
\[
[w_\infty]\in \mathcal{M}_0(a,a;\gamma_{0,0},\gamma_{1,0})
=\{[\nu_{0,+}],[\nu_{0,-}]\}.
\]
Then the local implicit-function theorem from Proposition \ref{prop:at-least-two-small-d} forces $[w_n]$ to lie on the corresponding local branch for all sufficiently large $n$, contradicting the assumption that the $[w_n]$ is distinct from those two branches.
Therefore no additional small-$d$ trajectories exist.
\end{proof}



\bigskip



\begin{thebibliography}{0}


\bibitem{AG} D. N. Akhiezer and S. G. Gindikin.
\textit{On Stein extensions of real symmetric spaces}.
Math. Ann. 286 (1990), No. 1-3, 1-12.


\bibitem{Az} H. Azad.
\textit{Levi-curvature of manifolds with a Stein rational fibration}.
Manuscr. Math. 50 (1985), 269-311.


\bibitem{Balser}
W. Balser.
\textit{From Divergent Power Series to Analytic Functions:
Theory and Application of Multisummable Power Series}.
Lecture Notes in Mathematics, Vol. 1582, Springer-Verlag, Berlin--Heidelberg, 1994.


%\bibitem{BaerPfaeffle}
%C. B\"ar and F. Pf\"affle.
%\textit{Wiener measures on Riemannian manifolds and the Feynman--Kac formula}.
%arXiv:1108.5082, 2011.


%\bibitem{BGV} N. Berline, E. Getzler and M. Vergne.
%\textit{Heat Kernels and Dirac Operators}.
%Springer-Verlag, Berlin, 1992.

\bibitem{BiswasDumitrescuHRM} I. Biswas and S. Dumitrescu.
\textit{Holomorphic Riemannian metric and the fundamental group},
Bulletin de la Société Mathématique de France \textbf{147} (2019), no.3, 455--468.


\bibitem{BleisteinHandelsman} N. Bleistein and R. A. Handelsman,
\textit{Asymptotic Expansions of Integrals}.
Dover Publications, 1986.



%\bibitem{BHH} D. Burns, S. Halverscheid and R. Hind.
%\textit{The geometry of Grauert tubes and complexifcation of symmetric spaces}
%Duke Math. J. 118(3) (2003), 465-491.


%\bibitem{Cam} R. Camporesi.
%\textit{Harmonic analysis and propagators on homogeneous spaces}.
%Phys. Rep. 196 (1990), no. 1--2, 1--134.


%\bibitem{Cam2} R. Camporesi.
%\textit{On the analytic continuation of the Minakshisundaram--Pleijel zeta function for compact symmetric spaces of rank one}.
%J. Math. Anal. Appl. 214 (1997), no. 2, 524--549.


\bibitem{Costin}
O. Costin.
\textit{Asymptotics and Borel Summability}.
Monographs and Surveys in Pure and Applied Mathematics, Vol. 141,
Chapman \& Hall/CRC, Boca Raton, FL, 2008.


\bibitem{CPT} O. Costin, H. Park and Y. Takei.
\textit{Borel summability of the heat equation with variable coefficients}.
J. Differential Equations 252 (2012), no. 4, 3076--3092.


%\bibitem{DKOR} N. Dondi, I. Kalogerakis, D. Orlando and S. Reffert.
%\textit{Resurgence of the large-charge expansion}.
%JHEP 05 (2021), 035.


%\bibitem{DG} J. J. Duistermaat and V. W. Guillemin.
%\textit{The spectrum of positive elliptic operators and periodic bicharacteristics}.
%Invent. Math. 29 (1975), 39--79.


\bibitem{Du} G. V. Dunne.
\textit{Borel Summation and Analytic Continuation of the Heat Kernel on Hyperbolic Space}.
In \textit{Peter Suranyi Festschrift: A Life in Quantum Field Theory}, pp. 167--189, World Scientific, 2022.


\bibitem{Ecalle1981I}
J.~Écalle,
\textit{Les fonctions résurgentes. Tome I: Les algèbres de fonctions résurgentes},
Publications Mathématiques d'Orsay, 81-05,
Université de Paris-Sud, Département de Mathématique, Orsay, 1981.


\bibitem{Ecalle1981II}
J.~Écalle,
\textit{Les fonctions résurgentes. Tome II: Les fonctions résurgentes appliquées à l'itération},
Publications Mathématiques d'Orsay, 81-06,
Université de Paris-Sud, Département de Mathématique, Orsay, 1981.



\bibitem{Ecalle1985III}
J.~Écalle,
\textit{Les fonctions résurgentes. Tome III: L'équation du pont et la classification analytique des objets locaux},
Publications Mathématiques d'Orsay, 85-05,
Université de Paris-Sud, Département de Mathématique, Orsay, 1985.


\bibitem{ES} J. Eells, Jr. and J. H. Sampson.
\textit{Harmonic mappings of Riemannian manifolds}.
Amer. J. Math. 86 (1964), No. 1, 109--160.


%\bibitem{Evans} L. C. Evans.
%\textit{Partial Differential Equations}.
%Graduate Studies in Mathematics, Vol. 19, American Mathematical Society, Providence, RI, 2nd edition, 2010.


\bibitem{Gr} H. Grauert.
\textit{On Levi's problem and the imbedding of real-analytic manifolds}.
Ann. Math. 68 (1958), 460-472.


\bibitem{GS} V. Guillemin and M. Stenzel.
\textit{Grauert tubes and the homogeneous Monge-Amp\'ere equation}.
J. Differential Geom. 34 (1991), no.2, 561-570.


%\bibitem{GW} S. Gukov and E. Witten.
%\textit{Branes and Quantization}.
%arXiv:0809.0305v2.


\bibitem{Ham} R. S. Hamilton.
\textit{Harmonic maps of manifolds with boundary}.
Lecture Notes in Mathematics, Vol. 471. Springer-Verlag, Berlin--New York, 1975.


\bibitem{Har} T. Harg\'e.
\textit{Borel summation of the small time expansion of the heat kernel. The scalar potential case.}.
arXiv:1301.7742, 2013.


\bibitem{Har2} T. Harg\'e.
\textit{Borel summation of the small time expansion of the heat kernel with a vector potential}.
arXiv:1302.0604, 2013.

%\bibitem{HW} R. Harvey and R. O. Wells.
%\textit{Zero sets of non-negative strictly plurisubharmonic functions}.
%Math. Ann. 201 (1973), 165-170.


\bibitem{Henry1981}
D. Henry.
\textit{Geometric Theory of Semilinear Parabolic Equations}.
Lecture Notes in Mathematics, Vol. 840. Springer-Verlag, Berlin--New York, 1981.


\bibitem{HilgertNeeb2012}
J. Hilgert and K.-H. Neeb.
\textit{Structure and Geometry of Lie Groups}.
Springer Monographs in Mathematics. Springer, New York, 2012.


\bibitem{Hochschild1965}
G. Hochschild.
\textit{The Structure of Lie Groups}.
Holden-Day, San Francisco, 1965.


\bibitem{Hochschild1966}
G. Hochschild.
\textit{Complexification of real analytic groups}.
Trans. Amer. Math. Soc. 125 (1966), no. 3, 406--413.


%\bibitem{Hsu} E. P. Hsu.
%\textit{Stochastic Analysis on Manifolds}.
%Graduate Studies in Mathematics 38, American Mathematical Society, Providence, RI, 2002.


\bibitem{HS} A. Huckleberry and D. Snow.
\textit{A classification of strictly pseudoconcave homogeneous manifold}.
Ann. Scuola Norm. Sup. Pisa Cl. Sci. (4) 8 (1981), no. 2, 231--255.

%\bibitem{Kac}
%M. Kac.
%\textit{On distributions of certain Wiener functionals}.
%Transactions of the American Mathematical Society, Vol. 65, pp. 1--13, 1949.


%\bibitem{KM} S.-J. Kan and D. Ma.
%\textit{On the rigidity of Grauert tubes over locally symmetric spaces}.
%J. Reine Angew. Math. 524 (2000), 205-225.


\bibitem{Kont} M. Kontsevich.
\textit{Private discussion}, YMSC \& BIMSA, 2023.

\bibitem{Kont2020} M. Kontsevich.
\textit{Talk Slides: Exponential integrals, Lefschetz thimbles and linear resurgence, 30 June, 2020}. 

\bibitem{KS-2022}
M.~Kontsevich, and Y.~Soibelman. \textit{Analyticity and resurgence in wall-crossing formulas.} Lett. Math. Phys. 112 (2022), no. 2, Paper No. 32.


\bibitem{KS-2024}
M.~Kontsevich, and Y.~Soibelman. \textit{Holomorphic Floer theory I: exponential integrals in finite and infinite dimensions.} arXiv preprint arXiv:2402.07343 [math.SG], 2024.


\bibitem{Ku} R. S. Kulkarni.
\textit{On complexifications of differentiable manifolds}.
Invent. Math. 44 (1978), 49--64.


\bibitem{Le} L. Lempert.
\textit{Complex structures on the tangent bundle of Riemannian manifolds}.
In \textit{Complex Analysis and Geometry}, The University Series in Mathematics, Springer/Plenum, 1993, pp. 235--251.


\bibitem{LS} L. Lempert and R. Sz\"oke.
\textit{Global solutions of the homogeneous complex Monge-Amp\'ere equation and complex structures on the tangent bundle of Riemannian manifolds}.
Math. Ann. 290 (1991), no.4, 689--712.


\bibitem{LiLiTang2026}
S.~Li, Y.~Li, and X.~Tang,
\textit{Picard--Lefschetz Theory and Alien Calculus: A Case Study},
arXiv:2605.08867v1 [math-ph], 2026.


\bibitem{Lin1986}
X.-B. Lin.
\textit{Exponential dichotomies and homoclinic orbits in functional differential equations}.
J. Differential Equations 63 (1986), no. 2, 227--254.


\bibitem{LM1960} J.-L. Lions and B. Malgrange.
\textit{Sur l'unicit{\'e} r{\'e}trograde dans les probl{\`e}mes mixtes paraboliques}.
Math. Scand. 8 (1960), 277--286.


\bibitem{LMS} D. A. Lutz, M. Miyake and R. Sch\"afke.
\textit{On the Borel summability of divergent solutions of the heat equation}.
Nagoya Math. J. 154 (1999), 1--29.


\bibitem{MitschiSauzin2016}
C.~Mitschi and D.~Sauzin,
\textit{Divergent Series, Summability and Resurgence I:
Monodromy and Resurgence},
Lecture Notes in Mathematics, Vol.~2153, Springer, Cham, 2016.


\bibitem{MN} A. Morimoto and T. Nagano.
\textit{On pseudo-conformal transformations of hypersurfaces}.
J. Math. Soc. Japan 15 (1963), no.3, 289-300.


\bibitem{Neeb2000}
K.-H. Neeb.
\textit{Holomorphy and Convexity in Lie Theory}.
de Gruyter Expositions in Mathematics, Vol. 28. 
Walter de Gruyter, Berlin, 2000.

\bibitem{PW} G. Patrizio and P-M. Wong.
\textit{Stein manifolds with compact symmetric center}.
Math. Ann. 289 (1991), no. 3, 355--382.


\bibitem{PSS1997}
D. Peterhof, B. Sandstede and A. Scheel.
\textit{Exponential dichotomies for solitary-wave solutions of semilinear elliptic equations on infinite cylinders}.
J. Differential Equations 140 (1997), no. 2, 266--308.


\bibitem{PhamSaddlepoint}
F. Pham,
\textit{Vanishing homologies and the $n$ variable saddlepoint method}. 
In Singularities, Part 2 (Arcata, Calif., 1981), Proc. Sympos. Pure Math., Vol. 40, Amer. Math. Soc., 1983, pp. 319--333.


\bibitem{SauzinSplitting}
D. Sauzin.
\textit{Resurgent functions and splitting problems}.
RIMS K\^oky\^uroku \textbf{1493} (2006), 48--117.
%\bibitem{Sauzin}
%D. Sauzin.
%\textit{Introduction to 1-summability and resurgence}.
%arXiv:1405.0356, 2014.


\bibitem{Schnaubelt1999}
R. Schnaubelt.
\textit{Sufficient conditions for exponential stability and dichotomy of evolution equations}.
Forum Math. 11 (1999), no. 5, 543--566.


%\bibitem{Sim} B. Simon.
%\textit{Functional Integration and Quantum Physics}.
%Second edition, AMS Chelsea Publishing, Providence, RI, 2005.


\bibitem{Sz} R. Sz\"oke.
\textit{Complex structures on tangent bundles of Riemannian manifolds}.
Math. Ann. 291 (1991), no. 3, 409--428.


\bibitem{BW} H. Whitney and F. Bruhat.
\textit{Quelques propri\'et\'es fondamentales des ensembles analytiques-r\'eels}.
Comment. Math. Helv. 33 (1959) 132-160.


\bibitem{Witten2010}
E. Witten.
\textit{A new look at the path integral of quantum mechanics}.
In \textit{Surveys in Differential Geometry. Volume XV. Perspectives in Mathematics and Physics}, Int. Press, Somerville, MA, 2011, 345--419.



\bibitem{Wu-thesis}
Y.~Wu, \textit{Resurgence theory on the heat
kernel of Riemann surfaces}. Undegraduate thesis (2024), Qiuzhen college, Tsinghua University. 

\end{thebibliography}



\end{document}
